\begin{python}
from scripts.calc_10 import *
\end{python}

\section{Определение спектра периодического входного сигнала}

Для периодического сигнала спектральные характеристики представляют собой набор дискретных гармоник, кратных основной частоте сигнала $\omega_1$. Для заданного меандра с периодом $T = 2t_{\text{и}} = \pv[.2]{ex10['T']}$ основная частота равна:
\begin{equation}
    \omega_1 = \frac{2\pi}{T} = \frac{2\pi}{\pv[.2]{ex10['T']}} = \pv[.3]{ex10['w1']}
\end{equation}

Коэффициенты ряда Фурье для симметричного двуполярного меандра определяются аналитически:
\begin{itemize}
    \item Постоянная составляющая $A_0 = 0$.
    \item Амплитуды гармоник: $A_k = \frac{4 I_m}{k\pi}$ для нечетных $k$, и $A_k = 0$ для четных $k$.
    \item Начальные фазы: $\Phi_k = 0$ для всех гармоник.
\end{itemize}

Дискретные амплитудный и фазовый спектры, рассчитанные для $N=\pv[.0]{ex10['N']}$ гармоник, приведены на \cref{fig:plot_disc_spec_A,fig:plot_disc_spec_Phi}.

\begin{figure}[H]
    \centering
    %% Creator: Matplotlib, PGF backend
%%
%% To include the figure in your LaTeX document, write
%%   \input{<filename>.pgf}
%%
%% Make sure the required packages are loaded in your preamble
%%   \usepackage{pgf}
%%
%% Also ensure that all the required font packages are loaded; for instance,
%% the lmodern package is sometimes necessary when using math font.
%%   \usepackage{lmodern}
%%
%% Figures using additional raster images can only be included by \input if
%% they are in the same directory as the main LaTeX file. For loading figures
%% from other directories you can use the `import` package
%%   \usepackage{import}
%%
%% and then include the figures with
%%   \import{<path to file>}{<filename>.pgf}
%%
%% Matplotlib used the following preamble
%%   \def\mathdefault#1{#1}
%%   \everymath=\expandafter{\the\everymath\displaystyle}
%%   \IfFileExists{scrextend.sty}{
%%     \usepackage[fontsize=12.000000pt]{scrextend}
%%   }{
%%     \renewcommand{\normalsize}{\fontsize{12.000000}{14.400000}\selectfont}
%%     \normalsize
%%   }
%%   \usepackage{fontspec}\setmainfont{Times New Roman}\usepackage[english,russian]{babel}\usepackage{amsmath}
%%   \ifdefined\pdftexversion\else  % non-pdftex case.
%%     \usepackage{fontspec}
%%     \setmainfont{times.ttf}[Path=\detokenize{/usr/local/share/fonts/truetype/times-new-roman/}]
%%     \setsansfont{DejaVuSans.ttf}[Path=\detokenize{/usr/share/matplotlib/mpl-data/fonts/ttf/}]
%%     \setmonofont{DejaVuSansMono.ttf}[Path=\detokenize{/usr/share/matplotlib/mpl-data/fonts/ttf/}]
%%   \fi
%%   \makeatletter\@ifpackageloaded{underscore}{}{\usepackage[strings]{underscore}}\makeatother
%%
\begingroup%
\makeatletter%
\begin{pgfpicture}%
\pgfpathrectangle{\pgfpointorigin}{\pgfqpoint{6.490000in}{3.740000in}}%
\pgfusepath{use as bounding box, clip}%
\begin{pgfscope}%
\pgfsetbuttcap%
\pgfsetmiterjoin%
\definecolor{currentfill}{rgb}{1.000000,1.000000,1.000000}%
\pgfsetfillcolor{currentfill}%
\pgfsetlinewidth{0.000000pt}%
\definecolor{currentstroke}{rgb}{1.000000,1.000000,1.000000}%
\pgfsetstrokecolor{currentstroke}%
\pgfsetdash{}{0pt}%
\pgfpathmoveto{\pgfqpoint{0.000000in}{0.000000in}}%
\pgfpathlineto{\pgfqpoint{6.490000in}{0.000000in}}%
\pgfpathlineto{\pgfqpoint{6.490000in}{3.740000in}}%
\pgfpathlineto{\pgfqpoint{0.000000in}{3.740000in}}%
\pgfpathlineto{\pgfqpoint{0.000000in}{0.000000in}}%
\pgfpathclose%
\pgfusepath{fill}%
\end{pgfscope}%
\begin{pgfscope}%
\pgfsetbuttcap%
\pgfsetmiterjoin%
\definecolor{currentfill}{rgb}{1.000000,1.000000,1.000000}%
\pgfsetfillcolor{currentfill}%
\pgfsetlinewidth{0.000000pt}%
\definecolor{currentstroke}{rgb}{0.000000,0.000000,0.000000}%
\pgfsetstrokecolor{currentstroke}%
\pgfsetstrokeopacity{0.000000}%
\pgfsetdash{}{0pt}%
\pgfpathmoveto{\pgfqpoint{0.514058in}{0.456901in}}%
\pgfpathlineto{\pgfqpoint{6.440000in}{0.456901in}}%
\pgfpathlineto{\pgfqpoint{6.440000in}{3.690000in}}%
\pgfpathlineto{\pgfqpoint{0.514058in}{3.690000in}}%
\pgfpathlineto{\pgfqpoint{0.514058in}{0.456901in}}%
\pgfpathclose%
\pgfusepath{fill}%
\end{pgfscope}%
\begin{pgfscope}%
\pgfpathrectangle{\pgfqpoint{0.514058in}{0.456901in}}{\pgfqpoint{5.925942in}{3.233099in}}%
\pgfusepath{clip}%
\pgfsetbuttcap%
\pgfsetroundjoin%
\pgfsetlinewidth{0.501875pt}%
\definecolor{currentstroke}{rgb}{0.501961,0.501961,0.501961}%
\pgfsetstrokecolor{currentstroke}%
\pgfsetstrokeopacity{0.400000}%
\pgfsetdash{{1.850000pt}{0.800000pt}}{0.000000pt}%
\pgfpathmoveto{\pgfqpoint{0.783419in}{0.456901in}}%
\pgfpathlineto{\pgfqpoint{0.783419in}{3.690000in}}%
\pgfusepath{stroke}%
\end{pgfscope}%
\begin{pgfscope}%
\pgfsetbuttcap%
\pgfsetroundjoin%
\definecolor{currentfill}{rgb}{1.000000,1.000000,1.000000}%
\pgfsetfillcolor{currentfill}%
\pgfsetlinewidth{0.803000pt}%
\definecolor{currentstroke}{rgb}{0.000000,0.000000,0.000000}%
\pgfsetstrokecolor{currentstroke}%
\pgfsetdash{}{0pt}%
\pgfsys@defobject{currentmarker}{\pgfqpoint{0.000000in}{0.000000in}}{\pgfqpoint{0.000000in}{0.048611in}}{%
\pgfpathmoveto{\pgfqpoint{0.000000in}{0.000000in}}%
\pgfpathlineto{\pgfqpoint{0.000000in}{0.048611in}}%
\pgfusepath{stroke,fill}%
}%
\begin{pgfscope}%
\pgfsys@transformshift{0.783419in}{0.456901in}%
\pgfsys@useobject{currentmarker}{}%
\end{pgfscope}%
\end{pgfscope}%
\begin{pgfscope}%
\definecolor{textcolor}{rgb}{0.000000,0.000000,0.000000}%
\pgfsetstrokecolor{textcolor}%
\pgfsetfillcolor{textcolor}%
\pgftext[x=0.783419in,y=0.408290in,,top]{\color{textcolor}{\rmfamily\fontsize{12.000000}{14.400000}\selectfont\catcode`\^=\active\def^{\ifmmode\sp\else\^{}\fi}\catcode`\%=\active\def%{\%}$\mathdefault{0}$}}%
\end{pgfscope}%
\begin{pgfscope}%
\pgfpathrectangle{\pgfqpoint{0.514058in}{0.456901in}}{\pgfqpoint{5.925942in}{3.233099in}}%
\pgfusepath{clip}%
\pgfsetbuttcap%
\pgfsetroundjoin%
\pgfsetlinewidth{0.501875pt}%
\definecolor{currentstroke}{rgb}{0.501961,0.501961,0.501961}%
\pgfsetstrokecolor{currentstroke}%
\pgfsetstrokeopacity{0.400000}%
\pgfsetdash{{1.850000pt}{0.800000pt}}{0.000000pt}%
\pgfpathmoveto{\pgfqpoint{1.860863in}{0.456901in}}%
\pgfpathlineto{\pgfqpoint{1.860863in}{3.690000in}}%
\pgfusepath{stroke}%
\end{pgfscope}%
\begin{pgfscope}%
\pgfsetbuttcap%
\pgfsetroundjoin%
\definecolor{currentfill}{rgb}{1.000000,1.000000,1.000000}%
\pgfsetfillcolor{currentfill}%
\pgfsetlinewidth{0.803000pt}%
\definecolor{currentstroke}{rgb}{0.000000,0.000000,0.000000}%
\pgfsetstrokecolor{currentstroke}%
\pgfsetdash{}{0pt}%
\pgfsys@defobject{currentmarker}{\pgfqpoint{0.000000in}{0.000000in}}{\pgfqpoint{0.000000in}{0.048611in}}{%
\pgfpathmoveto{\pgfqpoint{0.000000in}{0.000000in}}%
\pgfpathlineto{\pgfqpoint{0.000000in}{0.048611in}}%
\pgfusepath{stroke,fill}%
}%
\begin{pgfscope}%
\pgfsys@transformshift{1.860863in}{0.456901in}%
\pgfsys@useobject{currentmarker}{}%
\end{pgfscope}%
\end{pgfscope}%
\begin{pgfscope}%
\definecolor{textcolor}{rgb}{0.000000,0.000000,0.000000}%
\pgfsetstrokecolor{textcolor}%
\pgfsetfillcolor{textcolor}%
\pgftext[x=1.860863in,y=0.408290in,,top]{\color{textcolor}{\rmfamily\fontsize{12.000000}{14.400000}\selectfont\catcode`\^=\active\def^{\ifmmode\sp\else\^{}\fi}\catcode`\%=\active\def%{\%}$\mathdefault{1}$}}%
\end{pgfscope}%
\begin{pgfscope}%
\pgfpathrectangle{\pgfqpoint{0.514058in}{0.456901in}}{\pgfqpoint{5.925942in}{3.233099in}}%
\pgfusepath{clip}%
\pgfsetbuttcap%
\pgfsetroundjoin%
\pgfsetlinewidth{0.501875pt}%
\definecolor{currentstroke}{rgb}{0.501961,0.501961,0.501961}%
\pgfsetstrokecolor{currentstroke}%
\pgfsetstrokeopacity{0.400000}%
\pgfsetdash{{1.850000pt}{0.800000pt}}{0.000000pt}%
\pgfpathmoveto{\pgfqpoint{4.015751in}{0.456901in}}%
\pgfpathlineto{\pgfqpoint{4.015751in}{3.690000in}}%
\pgfusepath{stroke}%
\end{pgfscope}%
\begin{pgfscope}%
\pgfsetbuttcap%
\pgfsetroundjoin%
\definecolor{currentfill}{rgb}{1.000000,1.000000,1.000000}%
\pgfsetfillcolor{currentfill}%
\pgfsetlinewidth{0.803000pt}%
\definecolor{currentstroke}{rgb}{0.000000,0.000000,0.000000}%
\pgfsetstrokecolor{currentstroke}%
\pgfsetdash{}{0pt}%
\pgfsys@defobject{currentmarker}{\pgfqpoint{0.000000in}{0.000000in}}{\pgfqpoint{0.000000in}{0.048611in}}{%
\pgfpathmoveto{\pgfqpoint{0.000000in}{0.000000in}}%
\pgfpathlineto{\pgfqpoint{0.000000in}{0.048611in}}%
\pgfusepath{stroke,fill}%
}%
\begin{pgfscope}%
\pgfsys@transformshift{4.015751in}{0.456901in}%
\pgfsys@useobject{currentmarker}{}%
\end{pgfscope}%
\end{pgfscope}%
\begin{pgfscope}%
\definecolor{textcolor}{rgb}{0.000000,0.000000,0.000000}%
\pgfsetstrokecolor{textcolor}%
\pgfsetfillcolor{textcolor}%
\pgftext[x=4.015751in,y=0.408290in,,top]{\color{textcolor}{\rmfamily\fontsize{12.000000}{14.400000}\selectfont\catcode`\^=\active\def^{\ifmmode\sp\else\^{}\fi}\catcode`\%=\active\def%{\%}$\mathdefault{3}$}}%
\end{pgfscope}%
\begin{pgfscope}%
\pgfpathrectangle{\pgfqpoint{0.514058in}{0.456901in}}{\pgfqpoint{5.925942in}{3.233099in}}%
\pgfusepath{clip}%
\pgfsetbuttcap%
\pgfsetroundjoin%
\pgfsetlinewidth{0.501875pt}%
\definecolor{currentstroke}{rgb}{0.501961,0.501961,0.501961}%
\pgfsetstrokecolor{currentstroke}%
\pgfsetstrokeopacity{0.400000}%
\pgfsetdash{{1.850000pt}{0.800000pt}}{0.000000pt}%
\pgfpathmoveto{\pgfqpoint{6.170639in}{0.456901in}}%
\pgfpathlineto{\pgfqpoint{6.170639in}{3.690000in}}%
\pgfusepath{stroke}%
\end{pgfscope}%
\begin{pgfscope}%
\pgfsetbuttcap%
\pgfsetroundjoin%
\definecolor{currentfill}{rgb}{1.000000,1.000000,1.000000}%
\pgfsetfillcolor{currentfill}%
\pgfsetlinewidth{0.803000pt}%
\definecolor{currentstroke}{rgb}{0.000000,0.000000,0.000000}%
\pgfsetstrokecolor{currentstroke}%
\pgfsetdash{}{0pt}%
\pgfsys@defobject{currentmarker}{\pgfqpoint{0.000000in}{0.000000in}}{\pgfqpoint{0.000000in}{0.048611in}}{%
\pgfpathmoveto{\pgfqpoint{0.000000in}{0.000000in}}%
\pgfpathlineto{\pgfqpoint{0.000000in}{0.048611in}}%
\pgfusepath{stroke,fill}%
}%
\begin{pgfscope}%
\pgfsys@transformshift{6.170639in}{0.456901in}%
\pgfsys@useobject{currentmarker}{}%
\end{pgfscope}%
\end{pgfscope}%
\begin{pgfscope}%
\definecolor{textcolor}{rgb}{0.000000,0.000000,0.000000}%
\pgfsetstrokecolor{textcolor}%
\pgfsetfillcolor{textcolor}%
\pgftext[x=6.170639in,y=0.408290in,,top]{\color{textcolor}{\rmfamily\fontsize{12.000000}{14.400000}\selectfont\catcode`\^=\active\def^{\ifmmode\sp\else\^{}\fi}\catcode`\%=\active\def%{\%}$\mathdefault{5}$}}%
\end{pgfscope}%
\begin{pgfscope}%
\pgfpathrectangle{\pgfqpoint{0.514058in}{0.456901in}}{\pgfqpoint{5.925942in}{3.233099in}}%
\pgfusepath{clip}%
\pgfsetbuttcap%
\pgfsetroundjoin%
\pgfsetlinewidth{0.501875pt}%
\definecolor{currentstroke}{rgb}{0.501961,0.501961,0.501961}%
\pgfsetstrokecolor{currentstroke}%
\pgfsetstrokeopacity{0.400000}%
\pgfsetdash{{1.850000pt}{0.800000pt}}{0.000000pt}%
\pgfpathmoveto{\pgfqpoint{0.567930in}{0.456901in}}%
\pgfpathlineto{\pgfqpoint{0.567930in}{3.690000in}}%
\pgfusepath{stroke}%
\end{pgfscope}%
\begin{pgfscope}%
\pgfsetbuttcap%
\pgfsetroundjoin%
\definecolor{currentfill}{rgb}{1.000000,1.000000,1.000000}%
\pgfsetfillcolor{currentfill}%
\pgfsetlinewidth{0.602250pt}%
\definecolor{currentstroke}{rgb}{0.000000,0.000000,0.000000}%
\pgfsetstrokecolor{currentstroke}%
\pgfsetdash{}{0pt}%
\pgfsys@defobject{currentmarker}{\pgfqpoint{0.000000in}{0.000000in}}{\pgfqpoint{0.000000in}{0.027778in}}{%
\pgfpathmoveto{\pgfqpoint{0.000000in}{0.000000in}}%
\pgfpathlineto{\pgfqpoint{0.000000in}{0.027778in}}%
\pgfusepath{stroke,fill}%
}%
\begin{pgfscope}%
\pgfsys@transformshift{0.567930in}{0.456901in}%
\pgfsys@useobject{currentmarker}{}%
\end{pgfscope}%
\end{pgfscope}%
\begin{pgfscope}%
\pgfpathrectangle{\pgfqpoint{0.514058in}{0.456901in}}{\pgfqpoint{5.925942in}{3.233099in}}%
\pgfusepath{clip}%
\pgfsetbuttcap%
\pgfsetroundjoin%
\pgfsetlinewidth{0.501875pt}%
\definecolor{currentstroke}{rgb}{0.501961,0.501961,0.501961}%
\pgfsetstrokecolor{currentstroke}%
\pgfsetstrokeopacity{0.400000}%
\pgfsetdash{{1.850000pt}{0.800000pt}}{0.000000pt}%
\pgfpathmoveto{\pgfqpoint{0.998908in}{0.456901in}}%
\pgfpathlineto{\pgfqpoint{0.998908in}{3.690000in}}%
\pgfusepath{stroke}%
\end{pgfscope}%
\begin{pgfscope}%
\pgfsetbuttcap%
\pgfsetroundjoin%
\definecolor{currentfill}{rgb}{1.000000,1.000000,1.000000}%
\pgfsetfillcolor{currentfill}%
\pgfsetlinewidth{0.602250pt}%
\definecolor{currentstroke}{rgb}{0.000000,0.000000,0.000000}%
\pgfsetstrokecolor{currentstroke}%
\pgfsetdash{}{0pt}%
\pgfsys@defobject{currentmarker}{\pgfqpoint{0.000000in}{0.000000in}}{\pgfqpoint{0.000000in}{0.027778in}}{%
\pgfpathmoveto{\pgfqpoint{0.000000in}{0.000000in}}%
\pgfpathlineto{\pgfqpoint{0.000000in}{0.027778in}}%
\pgfusepath{stroke,fill}%
}%
\begin{pgfscope}%
\pgfsys@transformshift{0.998908in}{0.456901in}%
\pgfsys@useobject{currentmarker}{}%
\end{pgfscope}%
\end{pgfscope}%
\begin{pgfscope}%
\pgfpathrectangle{\pgfqpoint{0.514058in}{0.456901in}}{\pgfqpoint{5.925942in}{3.233099in}}%
\pgfusepath{clip}%
\pgfsetbuttcap%
\pgfsetroundjoin%
\pgfsetlinewidth{0.501875pt}%
\definecolor{currentstroke}{rgb}{0.501961,0.501961,0.501961}%
\pgfsetstrokecolor{currentstroke}%
\pgfsetstrokeopacity{0.400000}%
\pgfsetdash{{1.850000pt}{0.800000pt}}{0.000000pt}%
\pgfpathmoveto{\pgfqpoint{1.214396in}{0.456901in}}%
\pgfpathlineto{\pgfqpoint{1.214396in}{3.690000in}}%
\pgfusepath{stroke}%
\end{pgfscope}%
\begin{pgfscope}%
\pgfsetbuttcap%
\pgfsetroundjoin%
\definecolor{currentfill}{rgb}{1.000000,1.000000,1.000000}%
\pgfsetfillcolor{currentfill}%
\pgfsetlinewidth{0.602250pt}%
\definecolor{currentstroke}{rgb}{0.000000,0.000000,0.000000}%
\pgfsetstrokecolor{currentstroke}%
\pgfsetdash{}{0pt}%
\pgfsys@defobject{currentmarker}{\pgfqpoint{0.000000in}{0.000000in}}{\pgfqpoint{0.000000in}{0.027778in}}{%
\pgfpathmoveto{\pgfqpoint{0.000000in}{0.000000in}}%
\pgfpathlineto{\pgfqpoint{0.000000in}{0.027778in}}%
\pgfusepath{stroke,fill}%
}%
\begin{pgfscope}%
\pgfsys@transformshift{1.214396in}{0.456901in}%
\pgfsys@useobject{currentmarker}{}%
\end{pgfscope}%
\end{pgfscope}%
\begin{pgfscope}%
\pgfpathrectangle{\pgfqpoint{0.514058in}{0.456901in}}{\pgfqpoint{5.925942in}{3.233099in}}%
\pgfusepath{clip}%
\pgfsetbuttcap%
\pgfsetroundjoin%
\pgfsetlinewidth{0.501875pt}%
\definecolor{currentstroke}{rgb}{0.501961,0.501961,0.501961}%
\pgfsetstrokecolor{currentstroke}%
\pgfsetstrokeopacity{0.400000}%
\pgfsetdash{{1.850000pt}{0.800000pt}}{0.000000pt}%
\pgfpathmoveto{\pgfqpoint{1.429885in}{0.456901in}}%
\pgfpathlineto{\pgfqpoint{1.429885in}{3.690000in}}%
\pgfusepath{stroke}%
\end{pgfscope}%
\begin{pgfscope}%
\pgfsetbuttcap%
\pgfsetroundjoin%
\definecolor{currentfill}{rgb}{1.000000,1.000000,1.000000}%
\pgfsetfillcolor{currentfill}%
\pgfsetlinewidth{0.602250pt}%
\definecolor{currentstroke}{rgb}{0.000000,0.000000,0.000000}%
\pgfsetstrokecolor{currentstroke}%
\pgfsetdash{}{0pt}%
\pgfsys@defobject{currentmarker}{\pgfqpoint{0.000000in}{0.000000in}}{\pgfqpoint{0.000000in}{0.027778in}}{%
\pgfpathmoveto{\pgfqpoint{0.000000in}{0.000000in}}%
\pgfpathlineto{\pgfqpoint{0.000000in}{0.027778in}}%
\pgfusepath{stroke,fill}%
}%
\begin{pgfscope}%
\pgfsys@transformshift{1.429885in}{0.456901in}%
\pgfsys@useobject{currentmarker}{}%
\end{pgfscope}%
\end{pgfscope}%
\begin{pgfscope}%
\pgfpathrectangle{\pgfqpoint{0.514058in}{0.456901in}}{\pgfqpoint{5.925942in}{3.233099in}}%
\pgfusepath{clip}%
\pgfsetbuttcap%
\pgfsetroundjoin%
\pgfsetlinewidth{0.501875pt}%
\definecolor{currentstroke}{rgb}{0.501961,0.501961,0.501961}%
\pgfsetstrokecolor{currentstroke}%
\pgfsetstrokeopacity{0.400000}%
\pgfsetdash{{1.850000pt}{0.800000pt}}{0.000000pt}%
\pgfpathmoveto{\pgfqpoint{1.645374in}{0.456901in}}%
\pgfpathlineto{\pgfqpoint{1.645374in}{3.690000in}}%
\pgfusepath{stroke}%
\end{pgfscope}%
\begin{pgfscope}%
\pgfsetbuttcap%
\pgfsetroundjoin%
\definecolor{currentfill}{rgb}{1.000000,1.000000,1.000000}%
\pgfsetfillcolor{currentfill}%
\pgfsetlinewidth{0.602250pt}%
\definecolor{currentstroke}{rgb}{0.000000,0.000000,0.000000}%
\pgfsetstrokecolor{currentstroke}%
\pgfsetdash{}{0pt}%
\pgfsys@defobject{currentmarker}{\pgfqpoint{0.000000in}{0.000000in}}{\pgfqpoint{0.000000in}{0.027778in}}{%
\pgfpathmoveto{\pgfqpoint{0.000000in}{0.000000in}}%
\pgfpathlineto{\pgfqpoint{0.000000in}{0.027778in}}%
\pgfusepath{stroke,fill}%
}%
\begin{pgfscope}%
\pgfsys@transformshift{1.645374in}{0.456901in}%
\pgfsys@useobject{currentmarker}{}%
\end{pgfscope}%
\end{pgfscope}%
\begin{pgfscope}%
\pgfpathrectangle{\pgfqpoint{0.514058in}{0.456901in}}{\pgfqpoint{5.925942in}{3.233099in}}%
\pgfusepath{clip}%
\pgfsetbuttcap%
\pgfsetroundjoin%
\pgfsetlinewidth{0.501875pt}%
\definecolor{currentstroke}{rgb}{0.501961,0.501961,0.501961}%
\pgfsetstrokecolor{currentstroke}%
\pgfsetstrokeopacity{0.400000}%
\pgfsetdash{{1.850000pt}{0.800000pt}}{0.000000pt}%
\pgfpathmoveto{\pgfqpoint{2.076352in}{0.456901in}}%
\pgfpathlineto{\pgfqpoint{2.076352in}{3.690000in}}%
\pgfusepath{stroke}%
\end{pgfscope}%
\begin{pgfscope}%
\pgfsetbuttcap%
\pgfsetroundjoin%
\definecolor{currentfill}{rgb}{1.000000,1.000000,1.000000}%
\pgfsetfillcolor{currentfill}%
\pgfsetlinewidth{0.602250pt}%
\definecolor{currentstroke}{rgb}{0.000000,0.000000,0.000000}%
\pgfsetstrokecolor{currentstroke}%
\pgfsetdash{}{0pt}%
\pgfsys@defobject{currentmarker}{\pgfqpoint{0.000000in}{0.000000in}}{\pgfqpoint{0.000000in}{0.027778in}}{%
\pgfpathmoveto{\pgfqpoint{0.000000in}{0.000000in}}%
\pgfpathlineto{\pgfqpoint{0.000000in}{0.027778in}}%
\pgfusepath{stroke,fill}%
}%
\begin{pgfscope}%
\pgfsys@transformshift{2.076352in}{0.456901in}%
\pgfsys@useobject{currentmarker}{}%
\end{pgfscope}%
\end{pgfscope}%
\begin{pgfscope}%
\pgfpathrectangle{\pgfqpoint{0.514058in}{0.456901in}}{\pgfqpoint{5.925942in}{3.233099in}}%
\pgfusepath{clip}%
\pgfsetbuttcap%
\pgfsetroundjoin%
\pgfsetlinewidth{0.501875pt}%
\definecolor{currentstroke}{rgb}{0.501961,0.501961,0.501961}%
\pgfsetstrokecolor{currentstroke}%
\pgfsetstrokeopacity{0.400000}%
\pgfsetdash{{1.850000pt}{0.800000pt}}{0.000000pt}%
\pgfpathmoveto{\pgfqpoint{2.291840in}{0.456901in}}%
\pgfpathlineto{\pgfqpoint{2.291840in}{3.690000in}}%
\pgfusepath{stroke}%
\end{pgfscope}%
\begin{pgfscope}%
\pgfsetbuttcap%
\pgfsetroundjoin%
\definecolor{currentfill}{rgb}{1.000000,1.000000,1.000000}%
\pgfsetfillcolor{currentfill}%
\pgfsetlinewidth{0.602250pt}%
\definecolor{currentstroke}{rgb}{0.000000,0.000000,0.000000}%
\pgfsetstrokecolor{currentstroke}%
\pgfsetdash{}{0pt}%
\pgfsys@defobject{currentmarker}{\pgfqpoint{0.000000in}{0.000000in}}{\pgfqpoint{0.000000in}{0.027778in}}{%
\pgfpathmoveto{\pgfqpoint{0.000000in}{0.000000in}}%
\pgfpathlineto{\pgfqpoint{0.000000in}{0.027778in}}%
\pgfusepath{stroke,fill}%
}%
\begin{pgfscope}%
\pgfsys@transformshift{2.291840in}{0.456901in}%
\pgfsys@useobject{currentmarker}{}%
\end{pgfscope}%
\end{pgfscope}%
\begin{pgfscope}%
\pgfpathrectangle{\pgfqpoint{0.514058in}{0.456901in}}{\pgfqpoint{5.925942in}{3.233099in}}%
\pgfusepath{clip}%
\pgfsetbuttcap%
\pgfsetroundjoin%
\pgfsetlinewidth{0.501875pt}%
\definecolor{currentstroke}{rgb}{0.501961,0.501961,0.501961}%
\pgfsetstrokecolor{currentstroke}%
\pgfsetstrokeopacity{0.400000}%
\pgfsetdash{{1.850000pt}{0.800000pt}}{0.000000pt}%
\pgfpathmoveto{\pgfqpoint{2.507329in}{0.456901in}}%
\pgfpathlineto{\pgfqpoint{2.507329in}{3.690000in}}%
\pgfusepath{stroke}%
\end{pgfscope}%
\begin{pgfscope}%
\pgfsetbuttcap%
\pgfsetroundjoin%
\definecolor{currentfill}{rgb}{1.000000,1.000000,1.000000}%
\pgfsetfillcolor{currentfill}%
\pgfsetlinewidth{0.602250pt}%
\definecolor{currentstroke}{rgb}{0.000000,0.000000,0.000000}%
\pgfsetstrokecolor{currentstroke}%
\pgfsetdash{}{0pt}%
\pgfsys@defobject{currentmarker}{\pgfqpoint{0.000000in}{0.000000in}}{\pgfqpoint{0.000000in}{0.027778in}}{%
\pgfpathmoveto{\pgfqpoint{0.000000in}{0.000000in}}%
\pgfpathlineto{\pgfqpoint{0.000000in}{0.027778in}}%
\pgfusepath{stroke,fill}%
}%
\begin{pgfscope}%
\pgfsys@transformshift{2.507329in}{0.456901in}%
\pgfsys@useobject{currentmarker}{}%
\end{pgfscope}%
\end{pgfscope}%
\begin{pgfscope}%
\pgfpathrectangle{\pgfqpoint{0.514058in}{0.456901in}}{\pgfqpoint{5.925942in}{3.233099in}}%
\pgfusepath{clip}%
\pgfsetbuttcap%
\pgfsetroundjoin%
\pgfsetlinewidth{0.501875pt}%
\definecolor{currentstroke}{rgb}{0.501961,0.501961,0.501961}%
\pgfsetstrokecolor{currentstroke}%
\pgfsetstrokeopacity{0.400000}%
\pgfsetdash{{1.850000pt}{0.800000pt}}{0.000000pt}%
\pgfpathmoveto{\pgfqpoint{2.722818in}{0.456901in}}%
\pgfpathlineto{\pgfqpoint{2.722818in}{3.690000in}}%
\pgfusepath{stroke}%
\end{pgfscope}%
\begin{pgfscope}%
\pgfsetbuttcap%
\pgfsetroundjoin%
\definecolor{currentfill}{rgb}{1.000000,1.000000,1.000000}%
\pgfsetfillcolor{currentfill}%
\pgfsetlinewidth{0.602250pt}%
\definecolor{currentstroke}{rgb}{0.000000,0.000000,0.000000}%
\pgfsetstrokecolor{currentstroke}%
\pgfsetdash{}{0pt}%
\pgfsys@defobject{currentmarker}{\pgfqpoint{0.000000in}{0.000000in}}{\pgfqpoint{0.000000in}{0.027778in}}{%
\pgfpathmoveto{\pgfqpoint{0.000000in}{0.000000in}}%
\pgfpathlineto{\pgfqpoint{0.000000in}{0.027778in}}%
\pgfusepath{stroke,fill}%
}%
\begin{pgfscope}%
\pgfsys@transformshift{2.722818in}{0.456901in}%
\pgfsys@useobject{currentmarker}{}%
\end{pgfscope}%
\end{pgfscope}%
\begin{pgfscope}%
\pgfpathrectangle{\pgfqpoint{0.514058in}{0.456901in}}{\pgfqpoint{5.925942in}{3.233099in}}%
\pgfusepath{clip}%
\pgfsetbuttcap%
\pgfsetroundjoin%
\pgfsetlinewidth{0.501875pt}%
\definecolor{currentstroke}{rgb}{0.501961,0.501961,0.501961}%
\pgfsetstrokecolor{currentstroke}%
\pgfsetstrokeopacity{0.400000}%
\pgfsetdash{{1.850000pt}{0.800000pt}}{0.000000pt}%
\pgfpathmoveto{\pgfqpoint{2.938307in}{0.456901in}}%
\pgfpathlineto{\pgfqpoint{2.938307in}{3.690000in}}%
\pgfusepath{stroke}%
\end{pgfscope}%
\begin{pgfscope}%
\pgfsetbuttcap%
\pgfsetroundjoin%
\definecolor{currentfill}{rgb}{1.000000,1.000000,1.000000}%
\pgfsetfillcolor{currentfill}%
\pgfsetlinewidth{0.602250pt}%
\definecolor{currentstroke}{rgb}{0.000000,0.000000,0.000000}%
\pgfsetstrokecolor{currentstroke}%
\pgfsetdash{}{0pt}%
\pgfsys@defobject{currentmarker}{\pgfqpoint{0.000000in}{0.000000in}}{\pgfqpoint{0.000000in}{0.027778in}}{%
\pgfpathmoveto{\pgfqpoint{0.000000in}{0.000000in}}%
\pgfpathlineto{\pgfqpoint{0.000000in}{0.027778in}}%
\pgfusepath{stroke,fill}%
}%
\begin{pgfscope}%
\pgfsys@transformshift{2.938307in}{0.456901in}%
\pgfsys@useobject{currentmarker}{}%
\end{pgfscope}%
\end{pgfscope}%
\begin{pgfscope}%
\pgfpathrectangle{\pgfqpoint{0.514058in}{0.456901in}}{\pgfqpoint{5.925942in}{3.233099in}}%
\pgfusepath{clip}%
\pgfsetbuttcap%
\pgfsetroundjoin%
\pgfsetlinewidth{0.501875pt}%
\definecolor{currentstroke}{rgb}{0.501961,0.501961,0.501961}%
\pgfsetstrokecolor{currentstroke}%
\pgfsetstrokeopacity{0.400000}%
\pgfsetdash{{1.850000pt}{0.800000pt}}{0.000000pt}%
\pgfpathmoveto{\pgfqpoint{3.153796in}{0.456901in}}%
\pgfpathlineto{\pgfqpoint{3.153796in}{3.690000in}}%
\pgfusepath{stroke}%
\end{pgfscope}%
\begin{pgfscope}%
\pgfsetbuttcap%
\pgfsetroundjoin%
\definecolor{currentfill}{rgb}{1.000000,1.000000,1.000000}%
\pgfsetfillcolor{currentfill}%
\pgfsetlinewidth{0.602250pt}%
\definecolor{currentstroke}{rgb}{0.000000,0.000000,0.000000}%
\pgfsetstrokecolor{currentstroke}%
\pgfsetdash{}{0pt}%
\pgfsys@defobject{currentmarker}{\pgfqpoint{0.000000in}{0.000000in}}{\pgfqpoint{0.000000in}{0.027778in}}{%
\pgfpathmoveto{\pgfqpoint{0.000000in}{0.000000in}}%
\pgfpathlineto{\pgfqpoint{0.000000in}{0.027778in}}%
\pgfusepath{stroke,fill}%
}%
\begin{pgfscope}%
\pgfsys@transformshift{3.153796in}{0.456901in}%
\pgfsys@useobject{currentmarker}{}%
\end{pgfscope}%
\end{pgfscope}%
\begin{pgfscope}%
\pgfpathrectangle{\pgfqpoint{0.514058in}{0.456901in}}{\pgfqpoint{5.925942in}{3.233099in}}%
\pgfusepath{clip}%
\pgfsetbuttcap%
\pgfsetroundjoin%
\pgfsetlinewidth{0.501875pt}%
\definecolor{currentstroke}{rgb}{0.501961,0.501961,0.501961}%
\pgfsetstrokecolor{currentstroke}%
\pgfsetstrokeopacity{0.400000}%
\pgfsetdash{{1.850000pt}{0.800000pt}}{0.000000pt}%
\pgfpathmoveto{\pgfqpoint{3.369284in}{0.456901in}}%
\pgfpathlineto{\pgfqpoint{3.369284in}{3.690000in}}%
\pgfusepath{stroke}%
\end{pgfscope}%
\begin{pgfscope}%
\pgfsetbuttcap%
\pgfsetroundjoin%
\definecolor{currentfill}{rgb}{1.000000,1.000000,1.000000}%
\pgfsetfillcolor{currentfill}%
\pgfsetlinewidth{0.602250pt}%
\definecolor{currentstroke}{rgb}{0.000000,0.000000,0.000000}%
\pgfsetstrokecolor{currentstroke}%
\pgfsetdash{}{0pt}%
\pgfsys@defobject{currentmarker}{\pgfqpoint{0.000000in}{0.000000in}}{\pgfqpoint{0.000000in}{0.027778in}}{%
\pgfpathmoveto{\pgfqpoint{0.000000in}{0.000000in}}%
\pgfpathlineto{\pgfqpoint{0.000000in}{0.027778in}}%
\pgfusepath{stroke,fill}%
}%
\begin{pgfscope}%
\pgfsys@transformshift{3.369284in}{0.456901in}%
\pgfsys@useobject{currentmarker}{}%
\end{pgfscope}%
\end{pgfscope}%
\begin{pgfscope}%
\pgfpathrectangle{\pgfqpoint{0.514058in}{0.456901in}}{\pgfqpoint{5.925942in}{3.233099in}}%
\pgfusepath{clip}%
\pgfsetbuttcap%
\pgfsetroundjoin%
\pgfsetlinewidth{0.501875pt}%
\definecolor{currentstroke}{rgb}{0.501961,0.501961,0.501961}%
\pgfsetstrokecolor{currentstroke}%
\pgfsetstrokeopacity{0.400000}%
\pgfsetdash{{1.850000pt}{0.800000pt}}{0.000000pt}%
\pgfpathmoveto{\pgfqpoint{3.584773in}{0.456901in}}%
\pgfpathlineto{\pgfqpoint{3.584773in}{3.690000in}}%
\pgfusepath{stroke}%
\end{pgfscope}%
\begin{pgfscope}%
\pgfsetbuttcap%
\pgfsetroundjoin%
\definecolor{currentfill}{rgb}{1.000000,1.000000,1.000000}%
\pgfsetfillcolor{currentfill}%
\pgfsetlinewidth{0.602250pt}%
\definecolor{currentstroke}{rgb}{0.000000,0.000000,0.000000}%
\pgfsetstrokecolor{currentstroke}%
\pgfsetdash{}{0pt}%
\pgfsys@defobject{currentmarker}{\pgfqpoint{0.000000in}{0.000000in}}{\pgfqpoint{0.000000in}{0.027778in}}{%
\pgfpathmoveto{\pgfqpoint{0.000000in}{0.000000in}}%
\pgfpathlineto{\pgfqpoint{0.000000in}{0.027778in}}%
\pgfusepath{stroke,fill}%
}%
\begin{pgfscope}%
\pgfsys@transformshift{3.584773in}{0.456901in}%
\pgfsys@useobject{currentmarker}{}%
\end{pgfscope}%
\end{pgfscope}%
\begin{pgfscope}%
\pgfpathrectangle{\pgfqpoint{0.514058in}{0.456901in}}{\pgfqpoint{5.925942in}{3.233099in}}%
\pgfusepath{clip}%
\pgfsetbuttcap%
\pgfsetroundjoin%
\pgfsetlinewidth{0.501875pt}%
\definecolor{currentstroke}{rgb}{0.501961,0.501961,0.501961}%
\pgfsetstrokecolor{currentstroke}%
\pgfsetstrokeopacity{0.400000}%
\pgfsetdash{{1.850000pt}{0.800000pt}}{0.000000pt}%
\pgfpathmoveto{\pgfqpoint{3.800262in}{0.456901in}}%
\pgfpathlineto{\pgfqpoint{3.800262in}{3.690000in}}%
\pgfusepath{stroke}%
\end{pgfscope}%
\begin{pgfscope}%
\pgfsetbuttcap%
\pgfsetroundjoin%
\definecolor{currentfill}{rgb}{1.000000,1.000000,1.000000}%
\pgfsetfillcolor{currentfill}%
\pgfsetlinewidth{0.602250pt}%
\definecolor{currentstroke}{rgb}{0.000000,0.000000,0.000000}%
\pgfsetstrokecolor{currentstroke}%
\pgfsetdash{}{0pt}%
\pgfsys@defobject{currentmarker}{\pgfqpoint{0.000000in}{0.000000in}}{\pgfqpoint{0.000000in}{0.027778in}}{%
\pgfpathmoveto{\pgfqpoint{0.000000in}{0.000000in}}%
\pgfpathlineto{\pgfqpoint{0.000000in}{0.027778in}}%
\pgfusepath{stroke,fill}%
}%
\begin{pgfscope}%
\pgfsys@transformshift{3.800262in}{0.456901in}%
\pgfsys@useobject{currentmarker}{}%
\end{pgfscope}%
\end{pgfscope}%
\begin{pgfscope}%
\pgfpathrectangle{\pgfqpoint{0.514058in}{0.456901in}}{\pgfqpoint{5.925942in}{3.233099in}}%
\pgfusepath{clip}%
\pgfsetbuttcap%
\pgfsetroundjoin%
\pgfsetlinewidth{0.501875pt}%
\definecolor{currentstroke}{rgb}{0.501961,0.501961,0.501961}%
\pgfsetstrokecolor{currentstroke}%
\pgfsetstrokeopacity{0.400000}%
\pgfsetdash{{1.850000pt}{0.800000pt}}{0.000000pt}%
\pgfpathmoveto{\pgfqpoint{4.231240in}{0.456901in}}%
\pgfpathlineto{\pgfqpoint{4.231240in}{3.690000in}}%
\pgfusepath{stroke}%
\end{pgfscope}%
\begin{pgfscope}%
\pgfsetbuttcap%
\pgfsetroundjoin%
\definecolor{currentfill}{rgb}{1.000000,1.000000,1.000000}%
\pgfsetfillcolor{currentfill}%
\pgfsetlinewidth{0.602250pt}%
\definecolor{currentstroke}{rgb}{0.000000,0.000000,0.000000}%
\pgfsetstrokecolor{currentstroke}%
\pgfsetdash{}{0pt}%
\pgfsys@defobject{currentmarker}{\pgfqpoint{0.000000in}{0.000000in}}{\pgfqpoint{0.000000in}{0.027778in}}{%
\pgfpathmoveto{\pgfqpoint{0.000000in}{0.000000in}}%
\pgfpathlineto{\pgfqpoint{0.000000in}{0.027778in}}%
\pgfusepath{stroke,fill}%
}%
\begin{pgfscope}%
\pgfsys@transformshift{4.231240in}{0.456901in}%
\pgfsys@useobject{currentmarker}{}%
\end{pgfscope}%
\end{pgfscope}%
\begin{pgfscope}%
\pgfpathrectangle{\pgfqpoint{0.514058in}{0.456901in}}{\pgfqpoint{5.925942in}{3.233099in}}%
\pgfusepath{clip}%
\pgfsetbuttcap%
\pgfsetroundjoin%
\pgfsetlinewidth{0.501875pt}%
\definecolor{currentstroke}{rgb}{0.501961,0.501961,0.501961}%
\pgfsetstrokecolor{currentstroke}%
\pgfsetstrokeopacity{0.400000}%
\pgfsetdash{{1.850000pt}{0.800000pt}}{0.000000pt}%
\pgfpathmoveto{\pgfqpoint{4.446729in}{0.456901in}}%
\pgfpathlineto{\pgfqpoint{4.446729in}{3.690000in}}%
\pgfusepath{stroke}%
\end{pgfscope}%
\begin{pgfscope}%
\pgfsetbuttcap%
\pgfsetroundjoin%
\definecolor{currentfill}{rgb}{1.000000,1.000000,1.000000}%
\pgfsetfillcolor{currentfill}%
\pgfsetlinewidth{0.602250pt}%
\definecolor{currentstroke}{rgb}{0.000000,0.000000,0.000000}%
\pgfsetstrokecolor{currentstroke}%
\pgfsetdash{}{0pt}%
\pgfsys@defobject{currentmarker}{\pgfqpoint{0.000000in}{0.000000in}}{\pgfqpoint{0.000000in}{0.027778in}}{%
\pgfpathmoveto{\pgfqpoint{0.000000in}{0.000000in}}%
\pgfpathlineto{\pgfqpoint{0.000000in}{0.027778in}}%
\pgfusepath{stroke,fill}%
}%
\begin{pgfscope}%
\pgfsys@transformshift{4.446729in}{0.456901in}%
\pgfsys@useobject{currentmarker}{}%
\end{pgfscope}%
\end{pgfscope}%
\begin{pgfscope}%
\pgfpathrectangle{\pgfqpoint{0.514058in}{0.456901in}}{\pgfqpoint{5.925942in}{3.233099in}}%
\pgfusepath{clip}%
\pgfsetbuttcap%
\pgfsetroundjoin%
\pgfsetlinewidth{0.501875pt}%
\definecolor{currentstroke}{rgb}{0.501961,0.501961,0.501961}%
\pgfsetstrokecolor{currentstroke}%
\pgfsetstrokeopacity{0.400000}%
\pgfsetdash{{1.850000pt}{0.800000pt}}{0.000000pt}%
\pgfpathmoveto{\pgfqpoint{4.662217in}{0.456901in}}%
\pgfpathlineto{\pgfqpoint{4.662217in}{3.690000in}}%
\pgfusepath{stroke}%
\end{pgfscope}%
\begin{pgfscope}%
\pgfsetbuttcap%
\pgfsetroundjoin%
\definecolor{currentfill}{rgb}{1.000000,1.000000,1.000000}%
\pgfsetfillcolor{currentfill}%
\pgfsetlinewidth{0.602250pt}%
\definecolor{currentstroke}{rgb}{0.000000,0.000000,0.000000}%
\pgfsetstrokecolor{currentstroke}%
\pgfsetdash{}{0pt}%
\pgfsys@defobject{currentmarker}{\pgfqpoint{0.000000in}{0.000000in}}{\pgfqpoint{0.000000in}{0.027778in}}{%
\pgfpathmoveto{\pgfqpoint{0.000000in}{0.000000in}}%
\pgfpathlineto{\pgfqpoint{0.000000in}{0.027778in}}%
\pgfusepath{stroke,fill}%
}%
\begin{pgfscope}%
\pgfsys@transformshift{4.662217in}{0.456901in}%
\pgfsys@useobject{currentmarker}{}%
\end{pgfscope}%
\end{pgfscope}%
\begin{pgfscope}%
\pgfpathrectangle{\pgfqpoint{0.514058in}{0.456901in}}{\pgfqpoint{5.925942in}{3.233099in}}%
\pgfusepath{clip}%
\pgfsetbuttcap%
\pgfsetroundjoin%
\pgfsetlinewidth{0.501875pt}%
\definecolor{currentstroke}{rgb}{0.501961,0.501961,0.501961}%
\pgfsetstrokecolor{currentstroke}%
\pgfsetstrokeopacity{0.400000}%
\pgfsetdash{{1.850000pt}{0.800000pt}}{0.000000pt}%
\pgfpathmoveto{\pgfqpoint{4.877706in}{0.456901in}}%
\pgfpathlineto{\pgfqpoint{4.877706in}{3.690000in}}%
\pgfusepath{stroke}%
\end{pgfscope}%
\begin{pgfscope}%
\pgfsetbuttcap%
\pgfsetroundjoin%
\definecolor{currentfill}{rgb}{1.000000,1.000000,1.000000}%
\pgfsetfillcolor{currentfill}%
\pgfsetlinewidth{0.602250pt}%
\definecolor{currentstroke}{rgb}{0.000000,0.000000,0.000000}%
\pgfsetstrokecolor{currentstroke}%
\pgfsetdash{}{0pt}%
\pgfsys@defobject{currentmarker}{\pgfqpoint{0.000000in}{0.000000in}}{\pgfqpoint{0.000000in}{0.027778in}}{%
\pgfpathmoveto{\pgfqpoint{0.000000in}{0.000000in}}%
\pgfpathlineto{\pgfqpoint{0.000000in}{0.027778in}}%
\pgfusepath{stroke,fill}%
}%
\begin{pgfscope}%
\pgfsys@transformshift{4.877706in}{0.456901in}%
\pgfsys@useobject{currentmarker}{}%
\end{pgfscope}%
\end{pgfscope}%
\begin{pgfscope}%
\pgfpathrectangle{\pgfqpoint{0.514058in}{0.456901in}}{\pgfqpoint{5.925942in}{3.233099in}}%
\pgfusepath{clip}%
\pgfsetbuttcap%
\pgfsetroundjoin%
\pgfsetlinewidth{0.501875pt}%
\definecolor{currentstroke}{rgb}{0.501961,0.501961,0.501961}%
\pgfsetstrokecolor{currentstroke}%
\pgfsetstrokeopacity{0.400000}%
\pgfsetdash{{1.850000pt}{0.800000pt}}{0.000000pt}%
\pgfpathmoveto{\pgfqpoint{5.093195in}{0.456901in}}%
\pgfpathlineto{\pgfqpoint{5.093195in}{3.690000in}}%
\pgfusepath{stroke}%
\end{pgfscope}%
\begin{pgfscope}%
\pgfsetbuttcap%
\pgfsetroundjoin%
\definecolor{currentfill}{rgb}{1.000000,1.000000,1.000000}%
\pgfsetfillcolor{currentfill}%
\pgfsetlinewidth{0.602250pt}%
\definecolor{currentstroke}{rgb}{0.000000,0.000000,0.000000}%
\pgfsetstrokecolor{currentstroke}%
\pgfsetdash{}{0pt}%
\pgfsys@defobject{currentmarker}{\pgfqpoint{0.000000in}{0.000000in}}{\pgfqpoint{0.000000in}{0.027778in}}{%
\pgfpathmoveto{\pgfqpoint{0.000000in}{0.000000in}}%
\pgfpathlineto{\pgfqpoint{0.000000in}{0.027778in}}%
\pgfusepath{stroke,fill}%
}%
\begin{pgfscope}%
\pgfsys@transformshift{5.093195in}{0.456901in}%
\pgfsys@useobject{currentmarker}{}%
\end{pgfscope}%
\end{pgfscope}%
\begin{pgfscope}%
\pgfpathrectangle{\pgfqpoint{0.514058in}{0.456901in}}{\pgfqpoint{5.925942in}{3.233099in}}%
\pgfusepath{clip}%
\pgfsetbuttcap%
\pgfsetroundjoin%
\pgfsetlinewidth{0.501875pt}%
\definecolor{currentstroke}{rgb}{0.501961,0.501961,0.501961}%
\pgfsetstrokecolor{currentstroke}%
\pgfsetstrokeopacity{0.400000}%
\pgfsetdash{{1.850000pt}{0.800000pt}}{0.000000pt}%
\pgfpathmoveto{\pgfqpoint{5.308684in}{0.456901in}}%
\pgfpathlineto{\pgfqpoint{5.308684in}{3.690000in}}%
\pgfusepath{stroke}%
\end{pgfscope}%
\begin{pgfscope}%
\pgfsetbuttcap%
\pgfsetroundjoin%
\definecolor{currentfill}{rgb}{1.000000,1.000000,1.000000}%
\pgfsetfillcolor{currentfill}%
\pgfsetlinewidth{0.602250pt}%
\definecolor{currentstroke}{rgb}{0.000000,0.000000,0.000000}%
\pgfsetstrokecolor{currentstroke}%
\pgfsetdash{}{0pt}%
\pgfsys@defobject{currentmarker}{\pgfqpoint{0.000000in}{0.000000in}}{\pgfqpoint{0.000000in}{0.027778in}}{%
\pgfpathmoveto{\pgfqpoint{0.000000in}{0.000000in}}%
\pgfpathlineto{\pgfqpoint{0.000000in}{0.027778in}}%
\pgfusepath{stroke,fill}%
}%
\begin{pgfscope}%
\pgfsys@transformshift{5.308684in}{0.456901in}%
\pgfsys@useobject{currentmarker}{}%
\end{pgfscope}%
\end{pgfscope}%
\begin{pgfscope}%
\pgfpathrectangle{\pgfqpoint{0.514058in}{0.456901in}}{\pgfqpoint{5.925942in}{3.233099in}}%
\pgfusepath{clip}%
\pgfsetbuttcap%
\pgfsetroundjoin%
\pgfsetlinewidth{0.501875pt}%
\definecolor{currentstroke}{rgb}{0.501961,0.501961,0.501961}%
\pgfsetstrokecolor{currentstroke}%
\pgfsetstrokeopacity{0.400000}%
\pgfsetdash{{1.850000pt}{0.800000pt}}{0.000000pt}%
\pgfpathmoveto{\pgfqpoint{5.524173in}{0.456901in}}%
\pgfpathlineto{\pgfqpoint{5.524173in}{3.690000in}}%
\pgfusepath{stroke}%
\end{pgfscope}%
\begin{pgfscope}%
\pgfsetbuttcap%
\pgfsetroundjoin%
\definecolor{currentfill}{rgb}{1.000000,1.000000,1.000000}%
\pgfsetfillcolor{currentfill}%
\pgfsetlinewidth{0.602250pt}%
\definecolor{currentstroke}{rgb}{0.000000,0.000000,0.000000}%
\pgfsetstrokecolor{currentstroke}%
\pgfsetdash{}{0pt}%
\pgfsys@defobject{currentmarker}{\pgfqpoint{0.000000in}{0.000000in}}{\pgfqpoint{0.000000in}{0.027778in}}{%
\pgfpathmoveto{\pgfqpoint{0.000000in}{0.000000in}}%
\pgfpathlineto{\pgfqpoint{0.000000in}{0.027778in}}%
\pgfusepath{stroke,fill}%
}%
\begin{pgfscope}%
\pgfsys@transformshift{5.524173in}{0.456901in}%
\pgfsys@useobject{currentmarker}{}%
\end{pgfscope}%
\end{pgfscope}%
\begin{pgfscope}%
\pgfpathrectangle{\pgfqpoint{0.514058in}{0.456901in}}{\pgfqpoint{5.925942in}{3.233099in}}%
\pgfusepath{clip}%
\pgfsetbuttcap%
\pgfsetroundjoin%
\pgfsetlinewidth{0.501875pt}%
\definecolor{currentstroke}{rgb}{0.501961,0.501961,0.501961}%
\pgfsetstrokecolor{currentstroke}%
\pgfsetstrokeopacity{0.400000}%
\pgfsetdash{{1.850000pt}{0.800000pt}}{0.000000pt}%
\pgfpathmoveto{\pgfqpoint{5.739661in}{0.456901in}}%
\pgfpathlineto{\pgfqpoint{5.739661in}{3.690000in}}%
\pgfusepath{stroke}%
\end{pgfscope}%
\begin{pgfscope}%
\pgfsetbuttcap%
\pgfsetroundjoin%
\definecolor{currentfill}{rgb}{1.000000,1.000000,1.000000}%
\pgfsetfillcolor{currentfill}%
\pgfsetlinewidth{0.602250pt}%
\definecolor{currentstroke}{rgb}{0.000000,0.000000,0.000000}%
\pgfsetstrokecolor{currentstroke}%
\pgfsetdash{}{0pt}%
\pgfsys@defobject{currentmarker}{\pgfqpoint{0.000000in}{0.000000in}}{\pgfqpoint{0.000000in}{0.027778in}}{%
\pgfpathmoveto{\pgfqpoint{0.000000in}{0.000000in}}%
\pgfpathlineto{\pgfqpoint{0.000000in}{0.027778in}}%
\pgfusepath{stroke,fill}%
}%
\begin{pgfscope}%
\pgfsys@transformshift{5.739661in}{0.456901in}%
\pgfsys@useobject{currentmarker}{}%
\end{pgfscope}%
\end{pgfscope}%
\begin{pgfscope}%
\pgfpathrectangle{\pgfqpoint{0.514058in}{0.456901in}}{\pgfqpoint{5.925942in}{3.233099in}}%
\pgfusepath{clip}%
\pgfsetbuttcap%
\pgfsetroundjoin%
\pgfsetlinewidth{0.501875pt}%
\definecolor{currentstroke}{rgb}{0.501961,0.501961,0.501961}%
\pgfsetstrokecolor{currentstroke}%
\pgfsetstrokeopacity{0.400000}%
\pgfsetdash{{1.850000pt}{0.800000pt}}{0.000000pt}%
\pgfpathmoveto{\pgfqpoint{5.955150in}{0.456901in}}%
\pgfpathlineto{\pgfqpoint{5.955150in}{3.690000in}}%
\pgfusepath{stroke}%
\end{pgfscope}%
\begin{pgfscope}%
\pgfsetbuttcap%
\pgfsetroundjoin%
\definecolor{currentfill}{rgb}{1.000000,1.000000,1.000000}%
\pgfsetfillcolor{currentfill}%
\pgfsetlinewidth{0.602250pt}%
\definecolor{currentstroke}{rgb}{0.000000,0.000000,0.000000}%
\pgfsetstrokecolor{currentstroke}%
\pgfsetdash{}{0pt}%
\pgfsys@defobject{currentmarker}{\pgfqpoint{0.000000in}{0.000000in}}{\pgfqpoint{0.000000in}{0.027778in}}{%
\pgfpathmoveto{\pgfqpoint{0.000000in}{0.000000in}}%
\pgfpathlineto{\pgfqpoint{0.000000in}{0.027778in}}%
\pgfusepath{stroke,fill}%
}%
\begin{pgfscope}%
\pgfsys@transformshift{5.955150in}{0.456901in}%
\pgfsys@useobject{currentmarker}{}%
\end{pgfscope}%
\end{pgfscope}%
\begin{pgfscope}%
\pgfpathrectangle{\pgfqpoint{0.514058in}{0.456901in}}{\pgfqpoint{5.925942in}{3.233099in}}%
\pgfusepath{clip}%
\pgfsetbuttcap%
\pgfsetroundjoin%
\pgfsetlinewidth{0.501875pt}%
\definecolor{currentstroke}{rgb}{0.501961,0.501961,0.501961}%
\pgfsetstrokecolor{currentstroke}%
\pgfsetstrokeopacity{0.400000}%
\pgfsetdash{{1.850000pt}{0.800000pt}}{0.000000pt}%
\pgfpathmoveto{\pgfqpoint{6.386128in}{0.456901in}}%
\pgfpathlineto{\pgfqpoint{6.386128in}{3.690000in}}%
\pgfusepath{stroke}%
\end{pgfscope}%
\begin{pgfscope}%
\pgfsetbuttcap%
\pgfsetroundjoin%
\definecolor{currentfill}{rgb}{1.000000,1.000000,1.000000}%
\pgfsetfillcolor{currentfill}%
\pgfsetlinewidth{0.602250pt}%
\definecolor{currentstroke}{rgb}{0.000000,0.000000,0.000000}%
\pgfsetstrokecolor{currentstroke}%
\pgfsetdash{}{0pt}%
\pgfsys@defobject{currentmarker}{\pgfqpoint{0.000000in}{0.000000in}}{\pgfqpoint{0.000000in}{0.027778in}}{%
\pgfpathmoveto{\pgfqpoint{0.000000in}{0.000000in}}%
\pgfpathlineto{\pgfqpoint{0.000000in}{0.027778in}}%
\pgfusepath{stroke,fill}%
}%
\begin{pgfscope}%
\pgfsys@transformshift{6.386128in}{0.456901in}%
\pgfsys@useobject{currentmarker}{}%
\end{pgfscope}%
\end{pgfscope}%
\begin{pgfscope}%
\definecolor{textcolor}{rgb}{0.000000,0.000000,0.000000}%
\pgfsetstrokecolor{textcolor}%
\pgfsetfillcolor{textcolor}%
\pgftext[x=3.477029in,y=0.201367in,,top]{\color{textcolor}{\rmfamily\fontsize{12.000000}{14.400000}\selectfont\catcode`\^=\active\def^{\ifmmode\sp\else\^{}\fi}\catcode`\%=\active\def%{\%}$k$}}%
\end{pgfscope}%
\begin{pgfscope}%
\pgfpathrectangle{\pgfqpoint{0.514058in}{0.456901in}}{\pgfqpoint{5.925942in}{3.233099in}}%
\pgfusepath{clip}%
\pgfsetbuttcap%
\pgfsetroundjoin%
\pgfsetlinewidth{0.501875pt}%
\definecolor{currentstroke}{rgb}{0.501961,0.501961,0.501961}%
\pgfsetstrokecolor{currentstroke}%
\pgfsetstrokeopacity{0.400000}%
\pgfsetdash{{1.850000pt}{0.800000pt}}{0.000000pt}%
\pgfpathmoveto{\pgfqpoint{0.514058in}{0.603860in}}%
\pgfpathlineto{\pgfqpoint{6.440000in}{0.603860in}}%
\pgfusepath{stroke}%
\end{pgfscope}%
\begin{pgfscope}%
\pgfsetbuttcap%
\pgfsetroundjoin%
\definecolor{currentfill}{rgb}{1.000000,1.000000,1.000000}%
\pgfsetfillcolor{currentfill}%
\pgfsetlinewidth{0.803000pt}%
\definecolor{currentstroke}{rgb}{0.000000,0.000000,0.000000}%
\pgfsetstrokecolor{currentstroke}%
\pgfsetdash{}{0pt}%
\pgfsys@defobject{currentmarker}{\pgfqpoint{0.000000in}{0.000000in}}{\pgfqpoint{0.048611in}{0.000000in}}{%
\pgfpathmoveto{\pgfqpoint{0.000000in}{0.000000in}}%
\pgfpathlineto{\pgfqpoint{0.048611in}{0.000000in}}%
\pgfusepath{stroke,fill}%
}%
\begin{pgfscope}%
\pgfsys@transformshift{0.514058in}{0.603860in}%
\pgfsys@useobject{currentmarker}{}%
\end{pgfscope}%
\end{pgfscope}%
\begin{pgfscope}%
\definecolor{textcolor}{rgb}{0.000000,0.000000,0.000000}%
\pgfsetstrokecolor{textcolor}%
\pgfsetfillcolor{textcolor}%
\pgftext[x=0.256923in, y=0.545999in, left, base]{\color{textcolor}{\rmfamily\fontsize{12.000000}{14.400000}\selectfont\catcode`\^=\active\def^{\ifmmode\sp\else\^{}\fi}\catcode`\%=\active\def%{\%}$\mathdefault{0.0}$}}%
\end{pgfscope}%
\begin{pgfscope}%
\pgfpathrectangle{\pgfqpoint{0.514058in}{0.456901in}}{\pgfqpoint{5.925942in}{3.233099in}}%
\pgfusepath{clip}%
\pgfsetbuttcap%
\pgfsetroundjoin%
\pgfsetlinewidth{0.501875pt}%
\definecolor{currentstroke}{rgb}{0.501961,0.501961,0.501961}%
\pgfsetstrokecolor{currentstroke}%
\pgfsetstrokeopacity{0.400000}%
\pgfsetdash{{1.850000pt}{0.800000pt}}{0.000000pt}%
\pgfpathmoveto{\pgfqpoint{0.514058in}{1.180967in}}%
\pgfpathlineto{\pgfqpoint{6.440000in}{1.180967in}}%
\pgfusepath{stroke}%
\end{pgfscope}%
\begin{pgfscope}%
\pgfsetbuttcap%
\pgfsetroundjoin%
\definecolor{currentfill}{rgb}{1.000000,1.000000,1.000000}%
\pgfsetfillcolor{currentfill}%
\pgfsetlinewidth{0.803000pt}%
\definecolor{currentstroke}{rgb}{0.000000,0.000000,0.000000}%
\pgfsetstrokecolor{currentstroke}%
\pgfsetdash{}{0pt}%
\pgfsys@defobject{currentmarker}{\pgfqpoint{0.000000in}{0.000000in}}{\pgfqpoint{0.048611in}{0.000000in}}{%
\pgfpathmoveto{\pgfqpoint{0.000000in}{0.000000in}}%
\pgfpathlineto{\pgfqpoint{0.048611in}{0.000000in}}%
\pgfusepath{stroke,fill}%
}%
\begin{pgfscope}%
\pgfsys@transformshift{0.514058in}{1.180967in}%
\pgfsys@useobject{currentmarker}{}%
\end{pgfscope}%
\end{pgfscope}%
\begin{pgfscope}%
\definecolor{textcolor}{rgb}{0.000000,0.000000,0.000000}%
\pgfsetstrokecolor{textcolor}%
\pgfsetfillcolor{textcolor}%
\pgftext[x=0.256923in, y=1.123105in, left, base]{\color{textcolor}{\rmfamily\fontsize{12.000000}{14.400000}\selectfont\catcode`\^=\active\def^{\ifmmode\sp\else\^{}\fi}\catcode`\%=\active\def%{\%}$\mathdefault{0.5}$}}%
\end{pgfscope}%
\begin{pgfscope}%
\pgfpathrectangle{\pgfqpoint{0.514058in}{0.456901in}}{\pgfqpoint{5.925942in}{3.233099in}}%
\pgfusepath{clip}%
\pgfsetbuttcap%
\pgfsetroundjoin%
\pgfsetlinewidth{0.501875pt}%
\definecolor{currentstroke}{rgb}{0.501961,0.501961,0.501961}%
\pgfsetstrokecolor{currentstroke}%
\pgfsetstrokeopacity{0.400000}%
\pgfsetdash{{1.850000pt}{0.800000pt}}{0.000000pt}%
\pgfpathmoveto{\pgfqpoint{0.514058in}{1.758074in}}%
\pgfpathlineto{\pgfqpoint{6.440000in}{1.758074in}}%
\pgfusepath{stroke}%
\end{pgfscope}%
\begin{pgfscope}%
\pgfsetbuttcap%
\pgfsetroundjoin%
\definecolor{currentfill}{rgb}{1.000000,1.000000,1.000000}%
\pgfsetfillcolor{currentfill}%
\pgfsetlinewidth{0.803000pt}%
\definecolor{currentstroke}{rgb}{0.000000,0.000000,0.000000}%
\pgfsetstrokecolor{currentstroke}%
\pgfsetdash{}{0pt}%
\pgfsys@defobject{currentmarker}{\pgfqpoint{0.000000in}{0.000000in}}{\pgfqpoint{0.048611in}{0.000000in}}{%
\pgfpathmoveto{\pgfqpoint{0.000000in}{0.000000in}}%
\pgfpathlineto{\pgfqpoint{0.048611in}{0.000000in}}%
\pgfusepath{stroke,fill}%
}%
\begin{pgfscope}%
\pgfsys@transformshift{0.514058in}{1.758074in}%
\pgfsys@useobject{currentmarker}{}%
\end{pgfscope}%
\end{pgfscope}%
\begin{pgfscope}%
\definecolor{textcolor}{rgb}{0.000000,0.000000,0.000000}%
\pgfsetstrokecolor{textcolor}%
\pgfsetfillcolor{textcolor}%
\pgftext[x=0.256923in, y=1.700212in, left, base]{\color{textcolor}{\rmfamily\fontsize{12.000000}{14.400000}\selectfont\catcode`\^=\active\def^{\ifmmode\sp\else\^{}\fi}\catcode`\%=\active\def%{\%}$\mathdefault{1.0}$}}%
\end{pgfscope}%
\begin{pgfscope}%
\pgfpathrectangle{\pgfqpoint{0.514058in}{0.456901in}}{\pgfqpoint{5.925942in}{3.233099in}}%
\pgfusepath{clip}%
\pgfsetbuttcap%
\pgfsetroundjoin%
\pgfsetlinewidth{0.501875pt}%
\definecolor{currentstroke}{rgb}{0.501961,0.501961,0.501961}%
\pgfsetstrokecolor{currentstroke}%
\pgfsetstrokeopacity{0.400000}%
\pgfsetdash{{1.850000pt}{0.800000pt}}{0.000000pt}%
\pgfpathmoveto{\pgfqpoint{0.514058in}{2.335180in}}%
\pgfpathlineto{\pgfqpoint{6.440000in}{2.335180in}}%
\pgfusepath{stroke}%
\end{pgfscope}%
\begin{pgfscope}%
\pgfsetbuttcap%
\pgfsetroundjoin%
\definecolor{currentfill}{rgb}{1.000000,1.000000,1.000000}%
\pgfsetfillcolor{currentfill}%
\pgfsetlinewidth{0.803000pt}%
\definecolor{currentstroke}{rgb}{0.000000,0.000000,0.000000}%
\pgfsetstrokecolor{currentstroke}%
\pgfsetdash{}{0pt}%
\pgfsys@defobject{currentmarker}{\pgfqpoint{0.000000in}{0.000000in}}{\pgfqpoint{0.048611in}{0.000000in}}{%
\pgfpathmoveto{\pgfqpoint{0.000000in}{0.000000in}}%
\pgfpathlineto{\pgfqpoint{0.048611in}{0.000000in}}%
\pgfusepath{stroke,fill}%
}%
\begin{pgfscope}%
\pgfsys@transformshift{0.514058in}{2.335180in}%
\pgfsys@useobject{currentmarker}{}%
\end{pgfscope}%
\end{pgfscope}%
\begin{pgfscope}%
\definecolor{textcolor}{rgb}{0.000000,0.000000,0.000000}%
\pgfsetstrokecolor{textcolor}%
\pgfsetfillcolor{textcolor}%
\pgftext[x=0.256923in, y=2.277319in, left, base]{\color{textcolor}{\rmfamily\fontsize{12.000000}{14.400000}\selectfont\catcode`\^=\active\def^{\ifmmode\sp\else\^{}\fi}\catcode`\%=\active\def%{\%}$\mathdefault{1.5}$}}%
\end{pgfscope}%
\begin{pgfscope}%
\pgfpathrectangle{\pgfqpoint{0.514058in}{0.456901in}}{\pgfqpoint{5.925942in}{3.233099in}}%
\pgfusepath{clip}%
\pgfsetbuttcap%
\pgfsetroundjoin%
\pgfsetlinewidth{0.501875pt}%
\definecolor{currentstroke}{rgb}{0.501961,0.501961,0.501961}%
\pgfsetstrokecolor{currentstroke}%
\pgfsetstrokeopacity{0.400000}%
\pgfsetdash{{1.850000pt}{0.800000pt}}{0.000000pt}%
\pgfpathmoveto{\pgfqpoint{0.514058in}{2.912287in}}%
\pgfpathlineto{\pgfqpoint{6.440000in}{2.912287in}}%
\pgfusepath{stroke}%
\end{pgfscope}%
\begin{pgfscope}%
\pgfsetbuttcap%
\pgfsetroundjoin%
\definecolor{currentfill}{rgb}{1.000000,1.000000,1.000000}%
\pgfsetfillcolor{currentfill}%
\pgfsetlinewidth{0.803000pt}%
\definecolor{currentstroke}{rgb}{0.000000,0.000000,0.000000}%
\pgfsetstrokecolor{currentstroke}%
\pgfsetdash{}{0pt}%
\pgfsys@defobject{currentmarker}{\pgfqpoint{0.000000in}{0.000000in}}{\pgfqpoint{0.048611in}{0.000000in}}{%
\pgfpathmoveto{\pgfqpoint{0.000000in}{0.000000in}}%
\pgfpathlineto{\pgfqpoint{0.048611in}{0.000000in}}%
\pgfusepath{stroke,fill}%
}%
\begin{pgfscope}%
\pgfsys@transformshift{0.514058in}{2.912287in}%
\pgfsys@useobject{currentmarker}{}%
\end{pgfscope}%
\end{pgfscope}%
\begin{pgfscope}%
\definecolor{textcolor}{rgb}{0.000000,0.000000,0.000000}%
\pgfsetstrokecolor{textcolor}%
\pgfsetfillcolor{textcolor}%
\pgftext[x=0.256923in, y=2.854426in, left, base]{\color{textcolor}{\rmfamily\fontsize{12.000000}{14.400000}\selectfont\catcode`\^=\active\def^{\ifmmode\sp\else\^{}\fi}\catcode`\%=\active\def%{\%}$\mathdefault{2.0}$}}%
\end{pgfscope}%
\begin{pgfscope}%
\pgfpathrectangle{\pgfqpoint{0.514058in}{0.456901in}}{\pgfqpoint{5.925942in}{3.233099in}}%
\pgfusepath{clip}%
\pgfsetbuttcap%
\pgfsetroundjoin%
\pgfsetlinewidth{0.501875pt}%
\definecolor{currentstroke}{rgb}{0.501961,0.501961,0.501961}%
\pgfsetstrokecolor{currentstroke}%
\pgfsetstrokeopacity{0.400000}%
\pgfsetdash{{1.850000pt}{0.800000pt}}{0.000000pt}%
\pgfpathmoveto{\pgfqpoint{0.514058in}{3.489394in}}%
\pgfpathlineto{\pgfqpoint{6.440000in}{3.489394in}}%
\pgfusepath{stroke}%
\end{pgfscope}%
\begin{pgfscope}%
\pgfsetbuttcap%
\pgfsetroundjoin%
\definecolor{currentfill}{rgb}{1.000000,1.000000,1.000000}%
\pgfsetfillcolor{currentfill}%
\pgfsetlinewidth{0.803000pt}%
\definecolor{currentstroke}{rgb}{0.000000,0.000000,0.000000}%
\pgfsetstrokecolor{currentstroke}%
\pgfsetdash{}{0pt}%
\pgfsys@defobject{currentmarker}{\pgfqpoint{0.000000in}{0.000000in}}{\pgfqpoint{0.048611in}{0.000000in}}{%
\pgfpathmoveto{\pgfqpoint{0.000000in}{0.000000in}}%
\pgfpathlineto{\pgfqpoint{0.048611in}{0.000000in}}%
\pgfusepath{stroke,fill}%
}%
\begin{pgfscope}%
\pgfsys@transformshift{0.514058in}{3.489394in}%
\pgfsys@useobject{currentmarker}{}%
\end{pgfscope}%
\end{pgfscope}%
\begin{pgfscope}%
\definecolor{textcolor}{rgb}{0.000000,0.000000,0.000000}%
\pgfsetstrokecolor{textcolor}%
\pgfsetfillcolor{textcolor}%
\pgftext[x=0.256923in, y=3.431533in, left, base]{\color{textcolor}{\rmfamily\fontsize{12.000000}{14.400000}\selectfont\catcode`\^=\active\def^{\ifmmode\sp\else\^{}\fi}\catcode`\%=\active\def%{\%}$\mathdefault{2.5}$}}%
\end{pgfscope}%
\begin{pgfscope}%
\definecolor{textcolor}{rgb}{0.000000,0.000000,0.000000}%
\pgfsetstrokecolor{textcolor}%
\pgfsetfillcolor{textcolor}%
\pgftext[x=0.201367in,y=2.073450in,,bottom,rotate=90.000000]{\color{textcolor}{\rmfamily\fontsize{12.000000}{14.400000}\selectfont\catcode`\^=\active\def^{\ifmmode\sp\else\^{}\fi}\catcode`\%=\active\def%{\%}$A_{1k}$}}%
\end{pgfscope}%
\begin{pgfscope}%
\pgfpathrectangle{\pgfqpoint{0.514058in}{0.456901in}}{\pgfqpoint{5.925942in}{3.233099in}}%
\pgfusepath{clip}%
\pgfsetbuttcap%
\pgfsetroundjoin%
\pgfsetlinewidth{2.007500pt}%
\definecolor{currentstroke}{rgb}{0.000000,0.000000,0.000000}%
\pgfsetstrokecolor{currentstroke}%
\pgfsetdash{{7.400000pt}{3.200000pt}}{0.000000pt}%
\pgfpathmoveto{\pgfqpoint{0.783419in}{0.603860in}}%
\pgfpathlineto{\pgfqpoint{0.783419in}{0.603860in}}%
\pgfusepath{stroke}%
\end{pgfscope}%
\begin{pgfscope}%
\pgfpathrectangle{\pgfqpoint{0.514058in}{0.456901in}}{\pgfqpoint{5.925942in}{3.233099in}}%
\pgfusepath{clip}%
\pgfsetbuttcap%
\pgfsetroundjoin%
\pgfsetlinewidth{2.007500pt}%
\definecolor{currentstroke}{rgb}{0.000000,0.000000,0.000000}%
\pgfsetstrokecolor{currentstroke}%
\pgfsetdash{{7.400000pt}{3.200000pt}}{0.000000pt}%
\pgfpathmoveto{\pgfqpoint{1.860863in}{0.603860in}}%
\pgfpathlineto{\pgfqpoint{1.860863in}{3.543041in}}%
\pgfusepath{stroke}%
\end{pgfscope}%
\begin{pgfscope}%
\pgfpathrectangle{\pgfqpoint{0.514058in}{0.456901in}}{\pgfqpoint{5.925942in}{3.233099in}}%
\pgfusepath{clip}%
\pgfsetbuttcap%
\pgfsetroundjoin%
\pgfsetlinewidth{2.007500pt}%
\definecolor{currentstroke}{rgb}{0.000000,0.000000,0.000000}%
\pgfsetstrokecolor{currentstroke}%
\pgfsetdash{{7.400000pt}{3.200000pt}}{0.000000pt}%
\pgfpathmoveto{\pgfqpoint{4.015751in}{0.603860in}}%
\pgfpathlineto{\pgfqpoint{4.015751in}{1.583587in}}%
\pgfusepath{stroke}%
\end{pgfscope}%
\begin{pgfscope}%
\pgfpathrectangle{\pgfqpoint{0.514058in}{0.456901in}}{\pgfqpoint{5.925942in}{3.233099in}}%
\pgfusepath{clip}%
\pgfsetbuttcap%
\pgfsetroundjoin%
\pgfsetlinewidth{2.007500pt}%
\definecolor{currentstroke}{rgb}{0.000000,0.000000,0.000000}%
\pgfsetstrokecolor{currentstroke}%
\pgfsetdash{{7.400000pt}{3.200000pt}}{0.000000pt}%
\pgfpathmoveto{\pgfqpoint{6.170639in}{0.603860in}}%
\pgfpathlineto{\pgfqpoint{6.170639in}{1.191696in}}%
\pgfusepath{stroke}%
\end{pgfscope}%
\begin{pgfscope}%
\pgfpathrectangle{\pgfqpoint{0.514058in}{0.456901in}}{\pgfqpoint{5.925942in}{3.233099in}}%
\pgfusepath{clip}%
\pgfsetbuttcap%
\pgfsetroundjoin%
\definecolor{currentfill}{rgb}{1.000000,1.000000,1.000000}%
\pgfsetfillcolor{currentfill}%
\pgfsetlinewidth{1.505625pt}%
\definecolor{currentstroke}{rgb}{0.000000,0.000000,0.000000}%
\pgfsetstrokecolor{currentstroke}%
\pgfsetdash{}{0pt}%
\pgfsys@defobject{currentmarker}{\pgfqpoint{-0.027778in}{-0.027778in}}{\pgfqpoint{0.027778in}{0.027778in}}{%
\pgfpathmoveto{\pgfqpoint{0.000000in}{-0.027778in}}%
\pgfpathcurveto{\pgfqpoint{0.007367in}{-0.027778in}}{\pgfqpoint{0.014433in}{-0.024851in}}{\pgfqpoint{0.019642in}{-0.019642in}}%
\pgfpathcurveto{\pgfqpoint{0.024851in}{-0.014433in}}{\pgfqpoint{0.027778in}{-0.007367in}}{\pgfqpoint{0.027778in}{0.000000in}}%
\pgfpathcurveto{\pgfqpoint{0.027778in}{0.007367in}}{\pgfqpoint{0.024851in}{0.014433in}}{\pgfqpoint{0.019642in}{0.019642in}}%
\pgfpathcurveto{\pgfqpoint{0.014433in}{0.024851in}}{\pgfqpoint{0.007367in}{0.027778in}}{\pgfqpoint{0.000000in}{0.027778in}}%
\pgfpathcurveto{\pgfqpoint{-0.007367in}{0.027778in}}{\pgfqpoint{-0.014433in}{0.024851in}}{\pgfqpoint{-0.019642in}{0.019642in}}%
\pgfpathcurveto{\pgfqpoint{-0.024851in}{0.014433in}}{\pgfqpoint{-0.027778in}{0.007367in}}{\pgfqpoint{-0.027778in}{0.000000in}}%
\pgfpathcurveto{\pgfqpoint{-0.027778in}{-0.007367in}}{\pgfqpoint{-0.024851in}{-0.014433in}}{\pgfqpoint{-0.019642in}{-0.019642in}}%
\pgfpathcurveto{\pgfqpoint{-0.014433in}{-0.024851in}}{\pgfqpoint{-0.007367in}{-0.027778in}}{\pgfqpoint{0.000000in}{-0.027778in}}%
\pgfpathlineto{\pgfqpoint{0.000000in}{-0.027778in}}%
\pgfpathclose%
\pgfusepath{stroke,fill}%
}%
\begin{pgfscope}%
\pgfsys@transformshift{0.783419in}{0.603860in}%
\pgfsys@useobject{currentmarker}{}%
\end{pgfscope}%
\begin{pgfscope}%
\pgfsys@transformshift{1.860863in}{3.543041in}%
\pgfsys@useobject{currentmarker}{}%
\end{pgfscope}%
\begin{pgfscope}%
\pgfsys@transformshift{4.015751in}{1.583587in}%
\pgfsys@useobject{currentmarker}{}%
\end{pgfscope}%
\begin{pgfscope}%
\pgfsys@transformshift{6.170639in}{1.191696in}%
\pgfsys@useobject{currentmarker}{}%
\end{pgfscope}%
\end{pgfscope}%
\begin{pgfscope}%
\pgfpathrectangle{\pgfqpoint{0.514058in}{0.456901in}}{\pgfqpoint{5.925942in}{3.233099in}}%
\pgfusepath{clip}%
\pgfsetrectcap%
\pgfsetroundjoin%
\pgfsetlinewidth{2.007500pt}%
\definecolor{currentstroke}{rgb}{0.000000,0.000000,0.000000}%
\pgfsetstrokecolor{currentstroke}%
\pgfsetdash{}{0pt}%
\pgfpathmoveto{\pgfqpoint{0.783419in}{0.603860in}}%
\pgfpathlineto{\pgfqpoint{6.170639in}{0.603860in}}%
\pgfusepath{stroke}%
\end{pgfscope}%
\begin{pgfscope}%
\pgfsetrectcap%
\pgfsetmiterjoin%
\pgfsetlinewidth{0.602250pt}%
\definecolor{currentstroke}{rgb}{0.000000,0.000000,0.000000}%
\pgfsetstrokecolor{currentstroke}%
\pgfsetdash{}{0pt}%
\pgfpathmoveto{\pgfqpoint{0.514058in}{0.456901in}}%
\pgfpathlineto{\pgfqpoint{0.514058in}{3.690000in}}%
\pgfusepath{stroke}%
\end{pgfscope}%
\begin{pgfscope}%
\pgfsetrectcap%
\pgfsetmiterjoin%
\pgfsetlinewidth{0.602250pt}%
\definecolor{currentstroke}{rgb}{0.000000,0.000000,0.000000}%
\pgfsetstrokecolor{currentstroke}%
\pgfsetdash{}{0pt}%
\pgfpathmoveto{\pgfqpoint{0.514058in}{0.456901in}}%
\pgfpathlineto{\pgfqpoint{6.440000in}{0.456901in}}%
\pgfusepath{stroke}%
\end{pgfscope}%
\end{pgfpicture}%
\makeatother%
\endgroup%

    \caption{Дискретный амплитудный спектр входного сигнала}
    \label{fig:plot_disc_spec_A}
\end{figure}

\begin{figure}[H]
    \centering
    %% Creator: Matplotlib, PGF backend
%%
%% To include the figure in your LaTeX document, write
%%   \input{<filename>.pgf}
%%
%% Make sure the required packages are loaded in your preamble
%%   \usepackage{pgf}
%%
%% Also ensure that all the required font packages are loaded; for instance,
%% the lmodern package is sometimes necessary when using math font.
%%   \usepackage{lmodern}
%%
%% Figures using additional raster images can only be included by \input if
%% they are in the same directory as the main LaTeX file. For loading figures
%% from other directories you can use the `import` package
%%   \usepackage{import}
%%
%% and then include the figures with
%%   \import{<path to file>}{<filename>.pgf}
%%
%% Matplotlib used the following preamble
%%   \def\mathdefault#1{#1}
%%   \everymath=\expandafter{\the\everymath\displaystyle}
%%   \IfFileExists{scrextend.sty}{
%%     \usepackage[fontsize=12.000000pt]{scrextend}
%%   }{
%%     \renewcommand{\normalsize}{\fontsize{12.000000}{14.400000}\selectfont}
%%     \normalsize
%%   }
%%   \usepackage{fontspec}\setmainfont{Times New Roman}\usepackage[english,russian]{babel}\usepackage{amsmath}
%%   \ifdefined\pdftexversion\else  % non-pdftex case.
%%     \usepackage{fontspec}
%%     \setmainfont{times.ttf}[Path=\detokenize{/usr/local/share/fonts/truetype/times-new-roman/}]
%%     \setsansfont{DejaVuSans.ttf}[Path=\detokenize{/usr/share/matplotlib/mpl-data/fonts/ttf/}]
%%     \setmonofont{DejaVuSansMono.ttf}[Path=\detokenize{/usr/share/matplotlib/mpl-data/fonts/ttf/}]
%%   \fi
%%   \makeatletter\@ifpackageloaded{underscore}{}{\usepackage[strings]{underscore}}\makeatother
%%
\begingroup%
\makeatletter%
\begin{pgfpicture}%
\pgfpathrectangle{\pgfpointorigin}{\pgfqpoint{6.490000in}{3.740000in}}%
\pgfusepath{use as bounding box, clip}%
\begin{pgfscope}%
\pgfsetbuttcap%
\pgfsetmiterjoin%
\definecolor{currentfill}{rgb}{1.000000,1.000000,1.000000}%
\pgfsetfillcolor{currentfill}%
\pgfsetlinewidth{0.000000pt}%
\definecolor{currentstroke}{rgb}{1.000000,1.000000,1.000000}%
\pgfsetstrokecolor{currentstroke}%
\pgfsetdash{}{0pt}%
\pgfpathmoveto{\pgfqpoint{0.000000in}{0.000000in}}%
\pgfpathlineto{\pgfqpoint{6.490000in}{0.000000in}}%
\pgfpathlineto{\pgfqpoint{6.490000in}{3.740000in}}%
\pgfpathlineto{\pgfqpoint{0.000000in}{3.740000in}}%
\pgfpathlineto{\pgfqpoint{0.000000in}{0.000000in}}%
\pgfpathclose%
\pgfusepath{fill}%
\end{pgfscope}%
\begin{pgfscope}%
\pgfsetbuttcap%
\pgfsetmiterjoin%
\definecolor{currentfill}{rgb}{1.000000,1.000000,1.000000}%
\pgfsetfillcolor{currentfill}%
\pgfsetlinewidth{0.000000pt}%
\definecolor{currentstroke}{rgb}{0.000000,0.000000,0.000000}%
\pgfsetstrokecolor{currentstroke}%
\pgfsetstrokeopacity{0.000000}%
\pgfsetdash{}{0pt}%
\pgfpathmoveto{\pgfqpoint{0.725873in}{0.456901in}}%
\pgfpathlineto{\pgfqpoint{6.440000in}{0.456901in}}%
\pgfpathlineto{\pgfqpoint{6.440000in}{3.632139in}}%
\pgfpathlineto{\pgfqpoint{0.725873in}{3.632139in}}%
\pgfpathlineto{\pgfqpoint{0.725873in}{0.456901in}}%
\pgfpathclose%
\pgfusepath{fill}%
\end{pgfscope}%
\begin{pgfscope}%
\pgfpathrectangle{\pgfqpoint{0.725873in}{0.456901in}}{\pgfqpoint{5.714127in}{3.175238in}}%
\pgfusepath{clip}%
\pgfsetbuttcap%
\pgfsetroundjoin%
\pgfsetlinewidth{0.501875pt}%
\definecolor{currentstroke}{rgb}{0.501961,0.501961,0.501961}%
\pgfsetstrokecolor{currentstroke}%
\pgfsetstrokeopacity{0.400000}%
\pgfsetdash{{1.850000pt}{0.800000pt}}{0.000000pt}%
\pgfpathmoveto{\pgfqpoint{0.985606in}{0.456901in}}%
\pgfpathlineto{\pgfqpoint{0.985606in}{3.632139in}}%
\pgfusepath{stroke}%
\end{pgfscope}%
\begin{pgfscope}%
\pgfsetbuttcap%
\pgfsetroundjoin%
\definecolor{currentfill}{rgb}{1.000000,1.000000,1.000000}%
\pgfsetfillcolor{currentfill}%
\pgfsetlinewidth{0.803000pt}%
\definecolor{currentstroke}{rgb}{0.000000,0.000000,0.000000}%
\pgfsetstrokecolor{currentstroke}%
\pgfsetdash{}{0pt}%
\pgfsys@defobject{currentmarker}{\pgfqpoint{0.000000in}{0.000000in}}{\pgfqpoint{0.000000in}{0.048611in}}{%
\pgfpathmoveto{\pgfqpoint{0.000000in}{0.000000in}}%
\pgfpathlineto{\pgfqpoint{0.000000in}{0.048611in}}%
\pgfusepath{stroke,fill}%
}%
\begin{pgfscope}%
\pgfsys@transformshift{0.985606in}{0.456901in}%
\pgfsys@useobject{currentmarker}{}%
\end{pgfscope}%
\end{pgfscope}%
\begin{pgfscope}%
\definecolor{textcolor}{rgb}{0.000000,0.000000,0.000000}%
\pgfsetstrokecolor{textcolor}%
\pgfsetfillcolor{textcolor}%
\pgftext[x=0.985606in,y=0.408290in,,top]{\color{textcolor}{\rmfamily\fontsize{12.000000}{14.400000}\selectfont\catcode`\^=\active\def^{\ifmmode\sp\else\^{}\fi}\catcode`\%=\active\def%{\%}$\mathdefault{0}$}}%
\end{pgfscope}%
\begin{pgfscope}%
\pgfpathrectangle{\pgfqpoint{0.725873in}{0.456901in}}{\pgfqpoint{5.714127in}{3.175238in}}%
\pgfusepath{clip}%
\pgfsetbuttcap%
\pgfsetroundjoin%
\pgfsetlinewidth{0.501875pt}%
\definecolor{currentstroke}{rgb}{0.501961,0.501961,0.501961}%
\pgfsetstrokecolor{currentstroke}%
\pgfsetstrokeopacity{0.400000}%
\pgfsetdash{{1.850000pt}{0.800000pt}}{0.000000pt}%
\pgfpathmoveto{\pgfqpoint{2.024538in}{0.456901in}}%
\pgfpathlineto{\pgfqpoint{2.024538in}{3.632139in}}%
\pgfusepath{stroke}%
\end{pgfscope}%
\begin{pgfscope}%
\pgfsetbuttcap%
\pgfsetroundjoin%
\definecolor{currentfill}{rgb}{1.000000,1.000000,1.000000}%
\pgfsetfillcolor{currentfill}%
\pgfsetlinewidth{0.803000pt}%
\definecolor{currentstroke}{rgb}{0.000000,0.000000,0.000000}%
\pgfsetstrokecolor{currentstroke}%
\pgfsetdash{}{0pt}%
\pgfsys@defobject{currentmarker}{\pgfqpoint{0.000000in}{0.000000in}}{\pgfqpoint{0.000000in}{0.048611in}}{%
\pgfpathmoveto{\pgfqpoint{0.000000in}{0.000000in}}%
\pgfpathlineto{\pgfqpoint{0.000000in}{0.048611in}}%
\pgfusepath{stroke,fill}%
}%
\begin{pgfscope}%
\pgfsys@transformshift{2.024538in}{0.456901in}%
\pgfsys@useobject{currentmarker}{}%
\end{pgfscope}%
\end{pgfscope}%
\begin{pgfscope}%
\definecolor{textcolor}{rgb}{0.000000,0.000000,0.000000}%
\pgfsetstrokecolor{textcolor}%
\pgfsetfillcolor{textcolor}%
\pgftext[x=2.024538in,y=0.408290in,,top]{\color{textcolor}{\rmfamily\fontsize{12.000000}{14.400000}\selectfont\catcode`\^=\active\def^{\ifmmode\sp\else\^{}\fi}\catcode`\%=\active\def%{\%}$\mathdefault{1}$}}%
\end{pgfscope}%
\begin{pgfscope}%
\pgfpathrectangle{\pgfqpoint{0.725873in}{0.456901in}}{\pgfqpoint{5.714127in}{3.175238in}}%
\pgfusepath{clip}%
\pgfsetbuttcap%
\pgfsetroundjoin%
\pgfsetlinewidth{0.501875pt}%
\definecolor{currentstroke}{rgb}{0.501961,0.501961,0.501961}%
\pgfsetstrokecolor{currentstroke}%
\pgfsetstrokeopacity{0.400000}%
\pgfsetdash{{1.850000pt}{0.800000pt}}{0.000000pt}%
\pgfpathmoveto{\pgfqpoint{4.102403in}{0.456901in}}%
\pgfpathlineto{\pgfqpoint{4.102403in}{3.632139in}}%
\pgfusepath{stroke}%
\end{pgfscope}%
\begin{pgfscope}%
\pgfsetbuttcap%
\pgfsetroundjoin%
\definecolor{currentfill}{rgb}{1.000000,1.000000,1.000000}%
\pgfsetfillcolor{currentfill}%
\pgfsetlinewidth{0.803000pt}%
\definecolor{currentstroke}{rgb}{0.000000,0.000000,0.000000}%
\pgfsetstrokecolor{currentstroke}%
\pgfsetdash{}{0pt}%
\pgfsys@defobject{currentmarker}{\pgfqpoint{0.000000in}{0.000000in}}{\pgfqpoint{0.000000in}{0.048611in}}{%
\pgfpathmoveto{\pgfqpoint{0.000000in}{0.000000in}}%
\pgfpathlineto{\pgfqpoint{0.000000in}{0.048611in}}%
\pgfusepath{stroke,fill}%
}%
\begin{pgfscope}%
\pgfsys@transformshift{4.102403in}{0.456901in}%
\pgfsys@useobject{currentmarker}{}%
\end{pgfscope}%
\end{pgfscope}%
\begin{pgfscope}%
\definecolor{textcolor}{rgb}{0.000000,0.000000,0.000000}%
\pgfsetstrokecolor{textcolor}%
\pgfsetfillcolor{textcolor}%
\pgftext[x=4.102403in,y=0.408290in,,top]{\color{textcolor}{\rmfamily\fontsize{12.000000}{14.400000}\selectfont\catcode`\^=\active\def^{\ifmmode\sp\else\^{}\fi}\catcode`\%=\active\def%{\%}$\mathdefault{3}$}}%
\end{pgfscope}%
\begin{pgfscope}%
\pgfpathrectangle{\pgfqpoint{0.725873in}{0.456901in}}{\pgfqpoint{5.714127in}{3.175238in}}%
\pgfusepath{clip}%
\pgfsetbuttcap%
\pgfsetroundjoin%
\pgfsetlinewidth{0.501875pt}%
\definecolor{currentstroke}{rgb}{0.501961,0.501961,0.501961}%
\pgfsetstrokecolor{currentstroke}%
\pgfsetstrokeopacity{0.400000}%
\pgfsetdash{{1.850000pt}{0.800000pt}}{0.000000pt}%
\pgfpathmoveto{\pgfqpoint{6.180267in}{0.456901in}}%
\pgfpathlineto{\pgfqpoint{6.180267in}{3.632139in}}%
\pgfusepath{stroke}%
\end{pgfscope}%
\begin{pgfscope}%
\pgfsetbuttcap%
\pgfsetroundjoin%
\definecolor{currentfill}{rgb}{1.000000,1.000000,1.000000}%
\pgfsetfillcolor{currentfill}%
\pgfsetlinewidth{0.803000pt}%
\definecolor{currentstroke}{rgb}{0.000000,0.000000,0.000000}%
\pgfsetstrokecolor{currentstroke}%
\pgfsetdash{}{0pt}%
\pgfsys@defobject{currentmarker}{\pgfqpoint{0.000000in}{0.000000in}}{\pgfqpoint{0.000000in}{0.048611in}}{%
\pgfpathmoveto{\pgfqpoint{0.000000in}{0.000000in}}%
\pgfpathlineto{\pgfqpoint{0.000000in}{0.048611in}}%
\pgfusepath{stroke,fill}%
}%
\begin{pgfscope}%
\pgfsys@transformshift{6.180267in}{0.456901in}%
\pgfsys@useobject{currentmarker}{}%
\end{pgfscope}%
\end{pgfscope}%
\begin{pgfscope}%
\definecolor{textcolor}{rgb}{0.000000,0.000000,0.000000}%
\pgfsetstrokecolor{textcolor}%
\pgfsetfillcolor{textcolor}%
\pgftext[x=6.180267in,y=0.408290in,,top]{\color{textcolor}{\rmfamily\fontsize{12.000000}{14.400000}\selectfont\catcode`\^=\active\def^{\ifmmode\sp\else\^{}\fi}\catcode`\%=\active\def%{\%}$\mathdefault{5}$}}%
\end{pgfscope}%
\begin{pgfscope}%
\pgfpathrectangle{\pgfqpoint{0.725873in}{0.456901in}}{\pgfqpoint{5.714127in}{3.175238in}}%
\pgfusepath{clip}%
\pgfsetbuttcap%
\pgfsetroundjoin%
\pgfsetlinewidth{0.501875pt}%
\definecolor{currentstroke}{rgb}{0.501961,0.501961,0.501961}%
\pgfsetstrokecolor{currentstroke}%
\pgfsetstrokeopacity{0.400000}%
\pgfsetdash{{1.850000pt}{0.800000pt}}{0.000000pt}%
\pgfpathmoveto{\pgfqpoint{0.777819in}{0.456901in}}%
\pgfpathlineto{\pgfqpoint{0.777819in}{3.632139in}}%
\pgfusepath{stroke}%
\end{pgfscope}%
\begin{pgfscope}%
\pgfsetbuttcap%
\pgfsetroundjoin%
\definecolor{currentfill}{rgb}{1.000000,1.000000,1.000000}%
\pgfsetfillcolor{currentfill}%
\pgfsetlinewidth{0.602250pt}%
\definecolor{currentstroke}{rgb}{0.000000,0.000000,0.000000}%
\pgfsetstrokecolor{currentstroke}%
\pgfsetdash{}{0pt}%
\pgfsys@defobject{currentmarker}{\pgfqpoint{0.000000in}{0.000000in}}{\pgfqpoint{0.000000in}{0.027778in}}{%
\pgfpathmoveto{\pgfqpoint{0.000000in}{0.000000in}}%
\pgfpathlineto{\pgfqpoint{0.000000in}{0.027778in}}%
\pgfusepath{stroke,fill}%
}%
\begin{pgfscope}%
\pgfsys@transformshift{0.777819in}{0.456901in}%
\pgfsys@useobject{currentmarker}{}%
\end{pgfscope}%
\end{pgfscope}%
\begin{pgfscope}%
\pgfpathrectangle{\pgfqpoint{0.725873in}{0.456901in}}{\pgfqpoint{5.714127in}{3.175238in}}%
\pgfusepath{clip}%
\pgfsetbuttcap%
\pgfsetroundjoin%
\pgfsetlinewidth{0.501875pt}%
\definecolor{currentstroke}{rgb}{0.501961,0.501961,0.501961}%
\pgfsetstrokecolor{currentstroke}%
\pgfsetstrokeopacity{0.400000}%
\pgfsetdash{{1.850000pt}{0.800000pt}}{0.000000pt}%
\pgfpathmoveto{\pgfqpoint{1.193392in}{0.456901in}}%
\pgfpathlineto{\pgfqpoint{1.193392in}{3.632139in}}%
\pgfusepath{stroke}%
\end{pgfscope}%
\begin{pgfscope}%
\pgfsetbuttcap%
\pgfsetroundjoin%
\definecolor{currentfill}{rgb}{1.000000,1.000000,1.000000}%
\pgfsetfillcolor{currentfill}%
\pgfsetlinewidth{0.602250pt}%
\definecolor{currentstroke}{rgb}{0.000000,0.000000,0.000000}%
\pgfsetstrokecolor{currentstroke}%
\pgfsetdash{}{0pt}%
\pgfsys@defobject{currentmarker}{\pgfqpoint{0.000000in}{0.000000in}}{\pgfqpoint{0.000000in}{0.027778in}}{%
\pgfpathmoveto{\pgfqpoint{0.000000in}{0.000000in}}%
\pgfpathlineto{\pgfqpoint{0.000000in}{0.027778in}}%
\pgfusepath{stroke,fill}%
}%
\begin{pgfscope}%
\pgfsys@transformshift{1.193392in}{0.456901in}%
\pgfsys@useobject{currentmarker}{}%
\end{pgfscope}%
\end{pgfscope}%
\begin{pgfscope}%
\pgfpathrectangle{\pgfqpoint{0.725873in}{0.456901in}}{\pgfqpoint{5.714127in}{3.175238in}}%
\pgfusepath{clip}%
\pgfsetbuttcap%
\pgfsetroundjoin%
\pgfsetlinewidth{0.501875pt}%
\definecolor{currentstroke}{rgb}{0.501961,0.501961,0.501961}%
\pgfsetstrokecolor{currentstroke}%
\pgfsetstrokeopacity{0.400000}%
\pgfsetdash{{1.850000pt}{0.800000pt}}{0.000000pt}%
\pgfpathmoveto{\pgfqpoint{1.401179in}{0.456901in}}%
\pgfpathlineto{\pgfqpoint{1.401179in}{3.632139in}}%
\pgfusepath{stroke}%
\end{pgfscope}%
\begin{pgfscope}%
\pgfsetbuttcap%
\pgfsetroundjoin%
\definecolor{currentfill}{rgb}{1.000000,1.000000,1.000000}%
\pgfsetfillcolor{currentfill}%
\pgfsetlinewidth{0.602250pt}%
\definecolor{currentstroke}{rgb}{0.000000,0.000000,0.000000}%
\pgfsetstrokecolor{currentstroke}%
\pgfsetdash{}{0pt}%
\pgfsys@defobject{currentmarker}{\pgfqpoint{0.000000in}{0.000000in}}{\pgfqpoint{0.000000in}{0.027778in}}{%
\pgfpathmoveto{\pgfqpoint{0.000000in}{0.000000in}}%
\pgfpathlineto{\pgfqpoint{0.000000in}{0.027778in}}%
\pgfusepath{stroke,fill}%
}%
\begin{pgfscope}%
\pgfsys@transformshift{1.401179in}{0.456901in}%
\pgfsys@useobject{currentmarker}{}%
\end{pgfscope}%
\end{pgfscope}%
\begin{pgfscope}%
\pgfpathrectangle{\pgfqpoint{0.725873in}{0.456901in}}{\pgfqpoint{5.714127in}{3.175238in}}%
\pgfusepath{clip}%
\pgfsetbuttcap%
\pgfsetroundjoin%
\pgfsetlinewidth{0.501875pt}%
\definecolor{currentstroke}{rgb}{0.501961,0.501961,0.501961}%
\pgfsetstrokecolor{currentstroke}%
\pgfsetstrokeopacity{0.400000}%
\pgfsetdash{{1.850000pt}{0.800000pt}}{0.000000pt}%
\pgfpathmoveto{\pgfqpoint{1.608965in}{0.456901in}}%
\pgfpathlineto{\pgfqpoint{1.608965in}{3.632139in}}%
\pgfusepath{stroke}%
\end{pgfscope}%
\begin{pgfscope}%
\pgfsetbuttcap%
\pgfsetroundjoin%
\definecolor{currentfill}{rgb}{1.000000,1.000000,1.000000}%
\pgfsetfillcolor{currentfill}%
\pgfsetlinewidth{0.602250pt}%
\definecolor{currentstroke}{rgb}{0.000000,0.000000,0.000000}%
\pgfsetstrokecolor{currentstroke}%
\pgfsetdash{}{0pt}%
\pgfsys@defobject{currentmarker}{\pgfqpoint{0.000000in}{0.000000in}}{\pgfqpoint{0.000000in}{0.027778in}}{%
\pgfpathmoveto{\pgfqpoint{0.000000in}{0.000000in}}%
\pgfpathlineto{\pgfqpoint{0.000000in}{0.027778in}}%
\pgfusepath{stroke,fill}%
}%
\begin{pgfscope}%
\pgfsys@transformshift{1.608965in}{0.456901in}%
\pgfsys@useobject{currentmarker}{}%
\end{pgfscope}%
\end{pgfscope}%
\begin{pgfscope}%
\pgfpathrectangle{\pgfqpoint{0.725873in}{0.456901in}}{\pgfqpoint{5.714127in}{3.175238in}}%
\pgfusepath{clip}%
\pgfsetbuttcap%
\pgfsetroundjoin%
\pgfsetlinewidth{0.501875pt}%
\definecolor{currentstroke}{rgb}{0.501961,0.501961,0.501961}%
\pgfsetstrokecolor{currentstroke}%
\pgfsetstrokeopacity{0.400000}%
\pgfsetdash{{1.850000pt}{0.800000pt}}{0.000000pt}%
\pgfpathmoveto{\pgfqpoint{1.816752in}{0.456901in}}%
\pgfpathlineto{\pgfqpoint{1.816752in}{3.632139in}}%
\pgfusepath{stroke}%
\end{pgfscope}%
\begin{pgfscope}%
\pgfsetbuttcap%
\pgfsetroundjoin%
\definecolor{currentfill}{rgb}{1.000000,1.000000,1.000000}%
\pgfsetfillcolor{currentfill}%
\pgfsetlinewidth{0.602250pt}%
\definecolor{currentstroke}{rgb}{0.000000,0.000000,0.000000}%
\pgfsetstrokecolor{currentstroke}%
\pgfsetdash{}{0pt}%
\pgfsys@defobject{currentmarker}{\pgfqpoint{0.000000in}{0.000000in}}{\pgfqpoint{0.000000in}{0.027778in}}{%
\pgfpathmoveto{\pgfqpoint{0.000000in}{0.000000in}}%
\pgfpathlineto{\pgfqpoint{0.000000in}{0.027778in}}%
\pgfusepath{stroke,fill}%
}%
\begin{pgfscope}%
\pgfsys@transformshift{1.816752in}{0.456901in}%
\pgfsys@useobject{currentmarker}{}%
\end{pgfscope}%
\end{pgfscope}%
\begin{pgfscope}%
\pgfpathrectangle{\pgfqpoint{0.725873in}{0.456901in}}{\pgfqpoint{5.714127in}{3.175238in}}%
\pgfusepath{clip}%
\pgfsetbuttcap%
\pgfsetroundjoin%
\pgfsetlinewidth{0.501875pt}%
\definecolor{currentstroke}{rgb}{0.501961,0.501961,0.501961}%
\pgfsetstrokecolor{currentstroke}%
\pgfsetstrokeopacity{0.400000}%
\pgfsetdash{{1.850000pt}{0.800000pt}}{0.000000pt}%
\pgfpathmoveto{\pgfqpoint{2.232325in}{0.456901in}}%
\pgfpathlineto{\pgfqpoint{2.232325in}{3.632139in}}%
\pgfusepath{stroke}%
\end{pgfscope}%
\begin{pgfscope}%
\pgfsetbuttcap%
\pgfsetroundjoin%
\definecolor{currentfill}{rgb}{1.000000,1.000000,1.000000}%
\pgfsetfillcolor{currentfill}%
\pgfsetlinewidth{0.602250pt}%
\definecolor{currentstroke}{rgb}{0.000000,0.000000,0.000000}%
\pgfsetstrokecolor{currentstroke}%
\pgfsetdash{}{0pt}%
\pgfsys@defobject{currentmarker}{\pgfqpoint{0.000000in}{0.000000in}}{\pgfqpoint{0.000000in}{0.027778in}}{%
\pgfpathmoveto{\pgfqpoint{0.000000in}{0.000000in}}%
\pgfpathlineto{\pgfqpoint{0.000000in}{0.027778in}}%
\pgfusepath{stroke,fill}%
}%
\begin{pgfscope}%
\pgfsys@transformshift{2.232325in}{0.456901in}%
\pgfsys@useobject{currentmarker}{}%
\end{pgfscope}%
\end{pgfscope}%
\begin{pgfscope}%
\pgfpathrectangle{\pgfqpoint{0.725873in}{0.456901in}}{\pgfqpoint{5.714127in}{3.175238in}}%
\pgfusepath{clip}%
\pgfsetbuttcap%
\pgfsetroundjoin%
\pgfsetlinewidth{0.501875pt}%
\definecolor{currentstroke}{rgb}{0.501961,0.501961,0.501961}%
\pgfsetstrokecolor{currentstroke}%
\pgfsetstrokeopacity{0.400000}%
\pgfsetdash{{1.850000pt}{0.800000pt}}{0.000000pt}%
\pgfpathmoveto{\pgfqpoint{2.440111in}{0.456901in}}%
\pgfpathlineto{\pgfqpoint{2.440111in}{3.632139in}}%
\pgfusepath{stroke}%
\end{pgfscope}%
\begin{pgfscope}%
\pgfsetbuttcap%
\pgfsetroundjoin%
\definecolor{currentfill}{rgb}{1.000000,1.000000,1.000000}%
\pgfsetfillcolor{currentfill}%
\pgfsetlinewidth{0.602250pt}%
\definecolor{currentstroke}{rgb}{0.000000,0.000000,0.000000}%
\pgfsetstrokecolor{currentstroke}%
\pgfsetdash{}{0pt}%
\pgfsys@defobject{currentmarker}{\pgfqpoint{0.000000in}{0.000000in}}{\pgfqpoint{0.000000in}{0.027778in}}{%
\pgfpathmoveto{\pgfqpoint{0.000000in}{0.000000in}}%
\pgfpathlineto{\pgfqpoint{0.000000in}{0.027778in}}%
\pgfusepath{stroke,fill}%
}%
\begin{pgfscope}%
\pgfsys@transformshift{2.440111in}{0.456901in}%
\pgfsys@useobject{currentmarker}{}%
\end{pgfscope}%
\end{pgfscope}%
\begin{pgfscope}%
\pgfpathrectangle{\pgfqpoint{0.725873in}{0.456901in}}{\pgfqpoint{5.714127in}{3.175238in}}%
\pgfusepath{clip}%
\pgfsetbuttcap%
\pgfsetroundjoin%
\pgfsetlinewidth{0.501875pt}%
\definecolor{currentstroke}{rgb}{0.501961,0.501961,0.501961}%
\pgfsetstrokecolor{currentstroke}%
\pgfsetstrokeopacity{0.400000}%
\pgfsetdash{{1.850000pt}{0.800000pt}}{0.000000pt}%
\pgfpathmoveto{\pgfqpoint{2.647897in}{0.456901in}}%
\pgfpathlineto{\pgfqpoint{2.647897in}{3.632139in}}%
\pgfusepath{stroke}%
\end{pgfscope}%
\begin{pgfscope}%
\pgfsetbuttcap%
\pgfsetroundjoin%
\definecolor{currentfill}{rgb}{1.000000,1.000000,1.000000}%
\pgfsetfillcolor{currentfill}%
\pgfsetlinewidth{0.602250pt}%
\definecolor{currentstroke}{rgb}{0.000000,0.000000,0.000000}%
\pgfsetstrokecolor{currentstroke}%
\pgfsetdash{}{0pt}%
\pgfsys@defobject{currentmarker}{\pgfqpoint{0.000000in}{0.000000in}}{\pgfqpoint{0.000000in}{0.027778in}}{%
\pgfpathmoveto{\pgfqpoint{0.000000in}{0.000000in}}%
\pgfpathlineto{\pgfqpoint{0.000000in}{0.027778in}}%
\pgfusepath{stroke,fill}%
}%
\begin{pgfscope}%
\pgfsys@transformshift{2.647897in}{0.456901in}%
\pgfsys@useobject{currentmarker}{}%
\end{pgfscope}%
\end{pgfscope}%
\begin{pgfscope}%
\pgfpathrectangle{\pgfqpoint{0.725873in}{0.456901in}}{\pgfqpoint{5.714127in}{3.175238in}}%
\pgfusepath{clip}%
\pgfsetbuttcap%
\pgfsetroundjoin%
\pgfsetlinewidth{0.501875pt}%
\definecolor{currentstroke}{rgb}{0.501961,0.501961,0.501961}%
\pgfsetstrokecolor{currentstroke}%
\pgfsetstrokeopacity{0.400000}%
\pgfsetdash{{1.850000pt}{0.800000pt}}{0.000000pt}%
\pgfpathmoveto{\pgfqpoint{2.855684in}{0.456901in}}%
\pgfpathlineto{\pgfqpoint{2.855684in}{3.632139in}}%
\pgfusepath{stroke}%
\end{pgfscope}%
\begin{pgfscope}%
\pgfsetbuttcap%
\pgfsetroundjoin%
\definecolor{currentfill}{rgb}{1.000000,1.000000,1.000000}%
\pgfsetfillcolor{currentfill}%
\pgfsetlinewidth{0.602250pt}%
\definecolor{currentstroke}{rgb}{0.000000,0.000000,0.000000}%
\pgfsetstrokecolor{currentstroke}%
\pgfsetdash{}{0pt}%
\pgfsys@defobject{currentmarker}{\pgfqpoint{0.000000in}{0.000000in}}{\pgfqpoint{0.000000in}{0.027778in}}{%
\pgfpathmoveto{\pgfqpoint{0.000000in}{0.000000in}}%
\pgfpathlineto{\pgfqpoint{0.000000in}{0.027778in}}%
\pgfusepath{stroke,fill}%
}%
\begin{pgfscope}%
\pgfsys@transformshift{2.855684in}{0.456901in}%
\pgfsys@useobject{currentmarker}{}%
\end{pgfscope}%
\end{pgfscope}%
\begin{pgfscope}%
\pgfpathrectangle{\pgfqpoint{0.725873in}{0.456901in}}{\pgfqpoint{5.714127in}{3.175238in}}%
\pgfusepath{clip}%
\pgfsetbuttcap%
\pgfsetroundjoin%
\pgfsetlinewidth{0.501875pt}%
\definecolor{currentstroke}{rgb}{0.501961,0.501961,0.501961}%
\pgfsetstrokecolor{currentstroke}%
\pgfsetstrokeopacity{0.400000}%
\pgfsetdash{{1.850000pt}{0.800000pt}}{0.000000pt}%
\pgfpathmoveto{\pgfqpoint{3.063470in}{0.456901in}}%
\pgfpathlineto{\pgfqpoint{3.063470in}{3.632139in}}%
\pgfusepath{stroke}%
\end{pgfscope}%
\begin{pgfscope}%
\pgfsetbuttcap%
\pgfsetroundjoin%
\definecolor{currentfill}{rgb}{1.000000,1.000000,1.000000}%
\pgfsetfillcolor{currentfill}%
\pgfsetlinewidth{0.602250pt}%
\definecolor{currentstroke}{rgb}{0.000000,0.000000,0.000000}%
\pgfsetstrokecolor{currentstroke}%
\pgfsetdash{}{0pt}%
\pgfsys@defobject{currentmarker}{\pgfqpoint{0.000000in}{0.000000in}}{\pgfqpoint{0.000000in}{0.027778in}}{%
\pgfpathmoveto{\pgfqpoint{0.000000in}{0.000000in}}%
\pgfpathlineto{\pgfqpoint{0.000000in}{0.027778in}}%
\pgfusepath{stroke,fill}%
}%
\begin{pgfscope}%
\pgfsys@transformshift{3.063470in}{0.456901in}%
\pgfsys@useobject{currentmarker}{}%
\end{pgfscope}%
\end{pgfscope}%
\begin{pgfscope}%
\pgfpathrectangle{\pgfqpoint{0.725873in}{0.456901in}}{\pgfqpoint{5.714127in}{3.175238in}}%
\pgfusepath{clip}%
\pgfsetbuttcap%
\pgfsetroundjoin%
\pgfsetlinewidth{0.501875pt}%
\definecolor{currentstroke}{rgb}{0.501961,0.501961,0.501961}%
\pgfsetstrokecolor{currentstroke}%
\pgfsetstrokeopacity{0.400000}%
\pgfsetdash{{1.850000pt}{0.800000pt}}{0.000000pt}%
\pgfpathmoveto{\pgfqpoint{3.271257in}{0.456901in}}%
\pgfpathlineto{\pgfqpoint{3.271257in}{3.632139in}}%
\pgfusepath{stroke}%
\end{pgfscope}%
\begin{pgfscope}%
\pgfsetbuttcap%
\pgfsetroundjoin%
\definecolor{currentfill}{rgb}{1.000000,1.000000,1.000000}%
\pgfsetfillcolor{currentfill}%
\pgfsetlinewidth{0.602250pt}%
\definecolor{currentstroke}{rgb}{0.000000,0.000000,0.000000}%
\pgfsetstrokecolor{currentstroke}%
\pgfsetdash{}{0pt}%
\pgfsys@defobject{currentmarker}{\pgfqpoint{0.000000in}{0.000000in}}{\pgfqpoint{0.000000in}{0.027778in}}{%
\pgfpathmoveto{\pgfqpoint{0.000000in}{0.000000in}}%
\pgfpathlineto{\pgfqpoint{0.000000in}{0.027778in}}%
\pgfusepath{stroke,fill}%
}%
\begin{pgfscope}%
\pgfsys@transformshift{3.271257in}{0.456901in}%
\pgfsys@useobject{currentmarker}{}%
\end{pgfscope}%
\end{pgfscope}%
\begin{pgfscope}%
\pgfpathrectangle{\pgfqpoint{0.725873in}{0.456901in}}{\pgfqpoint{5.714127in}{3.175238in}}%
\pgfusepath{clip}%
\pgfsetbuttcap%
\pgfsetroundjoin%
\pgfsetlinewidth{0.501875pt}%
\definecolor{currentstroke}{rgb}{0.501961,0.501961,0.501961}%
\pgfsetstrokecolor{currentstroke}%
\pgfsetstrokeopacity{0.400000}%
\pgfsetdash{{1.850000pt}{0.800000pt}}{0.000000pt}%
\pgfpathmoveto{\pgfqpoint{3.479043in}{0.456901in}}%
\pgfpathlineto{\pgfqpoint{3.479043in}{3.632139in}}%
\pgfusepath{stroke}%
\end{pgfscope}%
\begin{pgfscope}%
\pgfsetbuttcap%
\pgfsetroundjoin%
\definecolor{currentfill}{rgb}{1.000000,1.000000,1.000000}%
\pgfsetfillcolor{currentfill}%
\pgfsetlinewidth{0.602250pt}%
\definecolor{currentstroke}{rgb}{0.000000,0.000000,0.000000}%
\pgfsetstrokecolor{currentstroke}%
\pgfsetdash{}{0pt}%
\pgfsys@defobject{currentmarker}{\pgfqpoint{0.000000in}{0.000000in}}{\pgfqpoint{0.000000in}{0.027778in}}{%
\pgfpathmoveto{\pgfqpoint{0.000000in}{0.000000in}}%
\pgfpathlineto{\pgfqpoint{0.000000in}{0.027778in}}%
\pgfusepath{stroke,fill}%
}%
\begin{pgfscope}%
\pgfsys@transformshift{3.479043in}{0.456901in}%
\pgfsys@useobject{currentmarker}{}%
\end{pgfscope}%
\end{pgfscope}%
\begin{pgfscope}%
\pgfpathrectangle{\pgfqpoint{0.725873in}{0.456901in}}{\pgfqpoint{5.714127in}{3.175238in}}%
\pgfusepath{clip}%
\pgfsetbuttcap%
\pgfsetroundjoin%
\pgfsetlinewidth{0.501875pt}%
\definecolor{currentstroke}{rgb}{0.501961,0.501961,0.501961}%
\pgfsetstrokecolor{currentstroke}%
\pgfsetstrokeopacity{0.400000}%
\pgfsetdash{{1.850000pt}{0.800000pt}}{0.000000pt}%
\pgfpathmoveto{\pgfqpoint{3.686830in}{0.456901in}}%
\pgfpathlineto{\pgfqpoint{3.686830in}{3.632139in}}%
\pgfusepath{stroke}%
\end{pgfscope}%
\begin{pgfscope}%
\pgfsetbuttcap%
\pgfsetroundjoin%
\definecolor{currentfill}{rgb}{1.000000,1.000000,1.000000}%
\pgfsetfillcolor{currentfill}%
\pgfsetlinewidth{0.602250pt}%
\definecolor{currentstroke}{rgb}{0.000000,0.000000,0.000000}%
\pgfsetstrokecolor{currentstroke}%
\pgfsetdash{}{0pt}%
\pgfsys@defobject{currentmarker}{\pgfqpoint{0.000000in}{0.000000in}}{\pgfqpoint{0.000000in}{0.027778in}}{%
\pgfpathmoveto{\pgfqpoint{0.000000in}{0.000000in}}%
\pgfpathlineto{\pgfqpoint{0.000000in}{0.027778in}}%
\pgfusepath{stroke,fill}%
}%
\begin{pgfscope}%
\pgfsys@transformshift{3.686830in}{0.456901in}%
\pgfsys@useobject{currentmarker}{}%
\end{pgfscope}%
\end{pgfscope}%
\begin{pgfscope}%
\pgfpathrectangle{\pgfqpoint{0.725873in}{0.456901in}}{\pgfqpoint{5.714127in}{3.175238in}}%
\pgfusepath{clip}%
\pgfsetbuttcap%
\pgfsetroundjoin%
\pgfsetlinewidth{0.501875pt}%
\definecolor{currentstroke}{rgb}{0.501961,0.501961,0.501961}%
\pgfsetstrokecolor{currentstroke}%
\pgfsetstrokeopacity{0.400000}%
\pgfsetdash{{1.850000pt}{0.800000pt}}{0.000000pt}%
\pgfpathmoveto{\pgfqpoint{3.894616in}{0.456901in}}%
\pgfpathlineto{\pgfqpoint{3.894616in}{3.632139in}}%
\pgfusepath{stroke}%
\end{pgfscope}%
\begin{pgfscope}%
\pgfsetbuttcap%
\pgfsetroundjoin%
\definecolor{currentfill}{rgb}{1.000000,1.000000,1.000000}%
\pgfsetfillcolor{currentfill}%
\pgfsetlinewidth{0.602250pt}%
\definecolor{currentstroke}{rgb}{0.000000,0.000000,0.000000}%
\pgfsetstrokecolor{currentstroke}%
\pgfsetdash{}{0pt}%
\pgfsys@defobject{currentmarker}{\pgfqpoint{0.000000in}{0.000000in}}{\pgfqpoint{0.000000in}{0.027778in}}{%
\pgfpathmoveto{\pgfqpoint{0.000000in}{0.000000in}}%
\pgfpathlineto{\pgfqpoint{0.000000in}{0.027778in}}%
\pgfusepath{stroke,fill}%
}%
\begin{pgfscope}%
\pgfsys@transformshift{3.894616in}{0.456901in}%
\pgfsys@useobject{currentmarker}{}%
\end{pgfscope}%
\end{pgfscope}%
\begin{pgfscope}%
\pgfpathrectangle{\pgfqpoint{0.725873in}{0.456901in}}{\pgfqpoint{5.714127in}{3.175238in}}%
\pgfusepath{clip}%
\pgfsetbuttcap%
\pgfsetroundjoin%
\pgfsetlinewidth{0.501875pt}%
\definecolor{currentstroke}{rgb}{0.501961,0.501961,0.501961}%
\pgfsetstrokecolor{currentstroke}%
\pgfsetstrokeopacity{0.400000}%
\pgfsetdash{{1.850000pt}{0.800000pt}}{0.000000pt}%
\pgfpathmoveto{\pgfqpoint{4.310189in}{0.456901in}}%
\pgfpathlineto{\pgfqpoint{4.310189in}{3.632139in}}%
\pgfusepath{stroke}%
\end{pgfscope}%
\begin{pgfscope}%
\pgfsetbuttcap%
\pgfsetroundjoin%
\definecolor{currentfill}{rgb}{1.000000,1.000000,1.000000}%
\pgfsetfillcolor{currentfill}%
\pgfsetlinewidth{0.602250pt}%
\definecolor{currentstroke}{rgb}{0.000000,0.000000,0.000000}%
\pgfsetstrokecolor{currentstroke}%
\pgfsetdash{}{0pt}%
\pgfsys@defobject{currentmarker}{\pgfqpoint{0.000000in}{0.000000in}}{\pgfqpoint{0.000000in}{0.027778in}}{%
\pgfpathmoveto{\pgfqpoint{0.000000in}{0.000000in}}%
\pgfpathlineto{\pgfqpoint{0.000000in}{0.027778in}}%
\pgfusepath{stroke,fill}%
}%
\begin{pgfscope}%
\pgfsys@transformshift{4.310189in}{0.456901in}%
\pgfsys@useobject{currentmarker}{}%
\end{pgfscope}%
\end{pgfscope}%
\begin{pgfscope}%
\pgfpathrectangle{\pgfqpoint{0.725873in}{0.456901in}}{\pgfqpoint{5.714127in}{3.175238in}}%
\pgfusepath{clip}%
\pgfsetbuttcap%
\pgfsetroundjoin%
\pgfsetlinewidth{0.501875pt}%
\definecolor{currentstroke}{rgb}{0.501961,0.501961,0.501961}%
\pgfsetstrokecolor{currentstroke}%
\pgfsetstrokeopacity{0.400000}%
\pgfsetdash{{1.850000pt}{0.800000pt}}{0.000000pt}%
\pgfpathmoveto{\pgfqpoint{4.517975in}{0.456901in}}%
\pgfpathlineto{\pgfqpoint{4.517975in}{3.632139in}}%
\pgfusepath{stroke}%
\end{pgfscope}%
\begin{pgfscope}%
\pgfsetbuttcap%
\pgfsetroundjoin%
\definecolor{currentfill}{rgb}{1.000000,1.000000,1.000000}%
\pgfsetfillcolor{currentfill}%
\pgfsetlinewidth{0.602250pt}%
\definecolor{currentstroke}{rgb}{0.000000,0.000000,0.000000}%
\pgfsetstrokecolor{currentstroke}%
\pgfsetdash{}{0pt}%
\pgfsys@defobject{currentmarker}{\pgfqpoint{0.000000in}{0.000000in}}{\pgfqpoint{0.000000in}{0.027778in}}{%
\pgfpathmoveto{\pgfqpoint{0.000000in}{0.000000in}}%
\pgfpathlineto{\pgfqpoint{0.000000in}{0.027778in}}%
\pgfusepath{stroke,fill}%
}%
\begin{pgfscope}%
\pgfsys@transformshift{4.517975in}{0.456901in}%
\pgfsys@useobject{currentmarker}{}%
\end{pgfscope}%
\end{pgfscope}%
\begin{pgfscope}%
\pgfpathrectangle{\pgfqpoint{0.725873in}{0.456901in}}{\pgfqpoint{5.714127in}{3.175238in}}%
\pgfusepath{clip}%
\pgfsetbuttcap%
\pgfsetroundjoin%
\pgfsetlinewidth{0.501875pt}%
\definecolor{currentstroke}{rgb}{0.501961,0.501961,0.501961}%
\pgfsetstrokecolor{currentstroke}%
\pgfsetstrokeopacity{0.400000}%
\pgfsetdash{{1.850000pt}{0.800000pt}}{0.000000pt}%
\pgfpathmoveto{\pgfqpoint{4.725762in}{0.456901in}}%
\pgfpathlineto{\pgfqpoint{4.725762in}{3.632139in}}%
\pgfusepath{stroke}%
\end{pgfscope}%
\begin{pgfscope}%
\pgfsetbuttcap%
\pgfsetroundjoin%
\definecolor{currentfill}{rgb}{1.000000,1.000000,1.000000}%
\pgfsetfillcolor{currentfill}%
\pgfsetlinewidth{0.602250pt}%
\definecolor{currentstroke}{rgb}{0.000000,0.000000,0.000000}%
\pgfsetstrokecolor{currentstroke}%
\pgfsetdash{}{0pt}%
\pgfsys@defobject{currentmarker}{\pgfqpoint{0.000000in}{0.000000in}}{\pgfqpoint{0.000000in}{0.027778in}}{%
\pgfpathmoveto{\pgfqpoint{0.000000in}{0.000000in}}%
\pgfpathlineto{\pgfqpoint{0.000000in}{0.027778in}}%
\pgfusepath{stroke,fill}%
}%
\begin{pgfscope}%
\pgfsys@transformshift{4.725762in}{0.456901in}%
\pgfsys@useobject{currentmarker}{}%
\end{pgfscope}%
\end{pgfscope}%
\begin{pgfscope}%
\pgfpathrectangle{\pgfqpoint{0.725873in}{0.456901in}}{\pgfqpoint{5.714127in}{3.175238in}}%
\pgfusepath{clip}%
\pgfsetbuttcap%
\pgfsetroundjoin%
\pgfsetlinewidth{0.501875pt}%
\definecolor{currentstroke}{rgb}{0.501961,0.501961,0.501961}%
\pgfsetstrokecolor{currentstroke}%
\pgfsetstrokeopacity{0.400000}%
\pgfsetdash{{1.850000pt}{0.800000pt}}{0.000000pt}%
\pgfpathmoveto{\pgfqpoint{4.933548in}{0.456901in}}%
\pgfpathlineto{\pgfqpoint{4.933548in}{3.632139in}}%
\pgfusepath{stroke}%
\end{pgfscope}%
\begin{pgfscope}%
\pgfsetbuttcap%
\pgfsetroundjoin%
\definecolor{currentfill}{rgb}{1.000000,1.000000,1.000000}%
\pgfsetfillcolor{currentfill}%
\pgfsetlinewidth{0.602250pt}%
\definecolor{currentstroke}{rgb}{0.000000,0.000000,0.000000}%
\pgfsetstrokecolor{currentstroke}%
\pgfsetdash{}{0pt}%
\pgfsys@defobject{currentmarker}{\pgfqpoint{0.000000in}{0.000000in}}{\pgfqpoint{0.000000in}{0.027778in}}{%
\pgfpathmoveto{\pgfqpoint{0.000000in}{0.000000in}}%
\pgfpathlineto{\pgfqpoint{0.000000in}{0.027778in}}%
\pgfusepath{stroke,fill}%
}%
\begin{pgfscope}%
\pgfsys@transformshift{4.933548in}{0.456901in}%
\pgfsys@useobject{currentmarker}{}%
\end{pgfscope}%
\end{pgfscope}%
\begin{pgfscope}%
\pgfpathrectangle{\pgfqpoint{0.725873in}{0.456901in}}{\pgfqpoint{5.714127in}{3.175238in}}%
\pgfusepath{clip}%
\pgfsetbuttcap%
\pgfsetroundjoin%
\pgfsetlinewidth{0.501875pt}%
\definecolor{currentstroke}{rgb}{0.501961,0.501961,0.501961}%
\pgfsetstrokecolor{currentstroke}%
\pgfsetstrokeopacity{0.400000}%
\pgfsetdash{{1.850000pt}{0.800000pt}}{0.000000pt}%
\pgfpathmoveto{\pgfqpoint{5.141335in}{0.456901in}}%
\pgfpathlineto{\pgfqpoint{5.141335in}{3.632139in}}%
\pgfusepath{stroke}%
\end{pgfscope}%
\begin{pgfscope}%
\pgfsetbuttcap%
\pgfsetroundjoin%
\definecolor{currentfill}{rgb}{1.000000,1.000000,1.000000}%
\pgfsetfillcolor{currentfill}%
\pgfsetlinewidth{0.602250pt}%
\definecolor{currentstroke}{rgb}{0.000000,0.000000,0.000000}%
\pgfsetstrokecolor{currentstroke}%
\pgfsetdash{}{0pt}%
\pgfsys@defobject{currentmarker}{\pgfqpoint{0.000000in}{0.000000in}}{\pgfqpoint{0.000000in}{0.027778in}}{%
\pgfpathmoveto{\pgfqpoint{0.000000in}{0.000000in}}%
\pgfpathlineto{\pgfqpoint{0.000000in}{0.027778in}}%
\pgfusepath{stroke,fill}%
}%
\begin{pgfscope}%
\pgfsys@transformshift{5.141335in}{0.456901in}%
\pgfsys@useobject{currentmarker}{}%
\end{pgfscope}%
\end{pgfscope}%
\begin{pgfscope}%
\pgfpathrectangle{\pgfqpoint{0.725873in}{0.456901in}}{\pgfqpoint{5.714127in}{3.175238in}}%
\pgfusepath{clip}%
\pgfsetbuttcap%
\pgfsetroundjoin%
\pgfsetlinewidth{0.501875pt}%
\definecolor{currentstroke}{rgb}{0.501961,0.501961,0.501961}%
\pgfsetstrokecolor{currentstroke}%
\pgfsetstrokeopacity{0.400000}%
\pgfsetdash{{1.850000pt}{0.800000pt}}{0.000000pt}%
\pgfpathmoveto{\pgfqpoint{5.349121in}{0.456901in}}%
\pgfpathlineto{\pgfqpoint{5.349121in}{3.632139in}}%
\pgfusepath{stroke}%
\end{pgfscope}%
\begin{pgfscope}%
\pgfsetbuttcap%
\pgfsetroundjoin%
\definecolor{currentfill}{rgb}{1.000000,1.000000,1.000000}%
\pgfsetfillcolor{currentfill}%
\pgfsetlinewidth{0.602250pt}%
\definecolor{currentstroke}{rgb}{0.000000,0.000000,0.000000}%
\pgfsetstrokecolor{currentstroke}%
\pgfsetdash{}{0pt}%
\pgfsys@defobject{currentmarker}{\pgfqpoint{0.000000in}{0.000000in}}{\pgfqpoint{0.000000in}{0.027778in}}{%
\pgfpathmoveto{\pgfqpoint{0.000000in}{0.000000in}}%
\pgfpathlineto{\pgfqpoint{0.000000in}{0.027778in}}%
\pgfusepath{stroke,fill}%
}%
\begin{pgfscope}%
\pgfsys@transformshift{5.349121in}{0.456901in}%
\pgfsys@useobject{currentmarker}{}%
\end{pgfscope}%
\end{pgfscope}%
\begin{pgfscope}%
\pgfpathrectangle{\pgfqpoint{0.725873in}{0.456901in}}{\pgfqpoint{5.714127in}{3.175238in}}%
\pgfusepath{clip}%
\pgfsetbuttcap%
\pgfsetroundjoin%
\pgfsetlinewidth{0.501875pt}%
\definecolor{currentstroke}{rgb}{0.501961,0.501961,0.501961}%
\pgfsetstrokecolor{currentstroke}%
\pgfsetstrokeopacity{0.400000}%
\pgfsetdash{{1.850000pt}{0.800000pt}}{0.000000pt}%
\pgfpathmoveto{\pgfqpoint{5.556908in}{0.456901in}}%
\pgfpathlineto{\pgfqpoint{5.556908in}{3.632139in}}%
\pgfusepath{stroke}%
\end{pgfscope}%
\begin{pgfscope}%
\pgfsetbuttcap%
\pgfsetroundjoin%
\definecolor{currentfill}{rgb}{1.000000,1.000000,1.000000}%
\pgfsetfillcolor{currentfill}%
\pgfsetlinewidth{0.602250pt}%
\definecolor{currentstroke}{rgb}{0.000000,0.000000,0.000000}%
\pgfsetstrokecolor{currentstroke}%
\pgfsetdash{}{0pt}%
\pgfsys@defobject{currentmarker}{\pgfqpoint{0.000000in}{0.000000in}}{\pgfqpoint{0.000000in}{0.027778in}}{%
\pgfpathmoveto{\pgfqpoint{0.000000in}{0.000000in}}%
\pgfpathlineto{\pgfqpoint{0.000000in}{0.027778in}}%
\pgfusepath{stroke,fill}%
}%
\begin{pgfscope}%
\pgfsys@transformshift{5.556908in}{0.456901in}%
\pgfsys@useobject{currentmarker}{}%
\end{pgfscope}%
\end{pgfscope}%
\begin{pgfscope}%
\pgfpathrectangle{\pgfqpoint{0.725873in}{0.456901in}}{\pgfqpoint{5.714127in}{3.175238in}}%
\pgfusepath{clip}%
\pgfsetbuttcap%
\pgfsetroundjoin%
\pgfsetlinewidth{0.501875pt}%
\definecolor{currentstroke}{rgb}{0.501961,0.501961,0.501961}%
\pgfsetstrokecolor{currentstroke}%
\pgfsetstrokeopacity{0.400000}%
\pgfsetdash{{1.850000pt}{0.800000pt}}{0.000000pt}%
\pgfpathmoveto{\pgfqpoint{5.764694in}{0.456901in}}%
\pgfpathlineto{\pgfqpoint{5.764694in}{3.632139in}}%
\pgfusepath{stroke}%
\end{pgfscope}%
\begin{pgfscope}%
\pgfsetbuttcap%
\pgfsetroundjoin%
\definecolor{currentfill}{rgb}{1.000000,1.000000,1.000000}%
\pgfsetfillcolor{currentfill}%
\pgfsetlinewidth{0.602250pt}%
\definecolor{currentstroke}{rgb}{0.000000,0.000000,0.000000}%
\pgfsetstrokecolor{currentstroke}%
\pgfsetdash{}{0pt}%
\pgfsys@defobject{currentmarker}{\pgfqpoint{0.000000in}{0.000000in}}{\pgfqpoint{0.000000in}{0.027778in}}{%
\pgfpathmoveto{\pgfqpoint{0.000000in}{0.000000in}}%
\pgfpathlineto{\pgfqpoint{0.000000in}{0.027778in}}%
\pgfusepath{stroke,fill}%
}%
\begin{pgfscope}%
\pgfsys@transformshift{5.764694in}{0.456901in}%
\pgfsys@useobject{currentmarker}{}%
\end{pgfscope}%
\end{pgfscope}%
\begin{pgfscope}%
\pgfpathrectangle{\pgfqpoint{0.725873in}{0.456901in}}{\pgfqpoint{5.714127in}{3.175238in}}%
\pgfusepath{clip}%
\pgfsetbuttcap%
\pgfsetroundjoin%
\pgfsetlinewidth{0.501875pt}%
\definecolor{currentstroke}{rgb}{0.501961,0.501961,0.501961}%
\pgfsetstrokecolor{currentstroke}%
\pgfsetstrokeopacity{0.400000}%
\pgfsetdash{{1.850000pt}{0.800000pt}}{0.000000pt}%
\pgfpathmoveto{\pgfqpoint{5.972481in}{0.456901in}}%
\pgfpathlineto{\pgfqpoint{5.972481in}{3.632139in}}%
\pgfusepath{stroke}%
\end{pgfscope}%
\begin{pgfscope}%
\pgfsetbuttcap%
\pgfsetroundjoin%
\definecolor{currentfill}{rgb}{1.000000,1.000000,1.000000}%
\pgfsetfillcolor{currentfill}%
\pgfsetlinewidth{0.602250pt}%
\definecolor{currentstroke}{rgb}{0.000000,0.000000,0.000000}%
\pgfsetstrokecolor{currentstroke}%
\pgfsetdash{}{0pt}%
\pgfsys@defobject{currentmarker}{\pgfqpoint{0.000000in}{0.000000in}}{\pgfqpoint{0.000000in}{0.027778in}}{%
\pgfpathmoveto{\pgfqpoint{0.000000in}{0.000000in}}%
\pgfpathlineto{\pgfqpoint{0.000000in}{0.027778in}}%
\pgfusepath{stroke,fill}%
}%
\begin{pgfscope}%
\pgfsys@transformshift{5.972481in}{0.456901in}%
\pgfsys@useobject{currentmarker}{}%
\end{pgfscope}%
\end{pgfscope}%
\begin{pgfscope}%
\pgfpathrectangle{\pgfqpoint{0.725873in}{0.456901in}}{\pgfqpoint{5.714127in}{3.175238in}}%
\pgfusepath{clip}%
\pgfsetbuttcap%
\pgfsetroundjoin%
\pgfsetlinewidth{0.501875pt}%
\definecolor{currentstroke}{rgb}{0.501961,0.501961,0.501961}%
\pgfsetstrokecolor{currentstroke}%
\pgfsetstrokeopacity{0.400000}%
\pgfsetdash{{1.850000pt}{0.800000pt}}{0.000000pt}%
\pgfpathmoveto{\pgfqpoint{6.388053in}{0.456901in}}%
\pgfpathlineto{\pgfqpoint{6.388053in}{3.632139in}}%
\pgfusepath{stroke}%
\end{pgfscope}%
\begin{pgfscope}%
\pgfsetbuttcap%
\pgfsetroundjoin%
\definecolor{currentfill}{rgb}{1.000000,1.000000,1.000000}%
\pgfsetfillcolor{currentfill}%
\pgfsetlinewidth{0.602250pt}%
\definecolor{currentstroke}{rgb}{0.000000,0.000000,0.000000}%
\pgfsetstrokecolor{currentstroke}%
\pgfsetdash{}{0pt}%
\pgfsys@defobject{currentmarker}{\pgfqpoint{0.000000in}{0.000000in}}{\pgfqpoint{0.000000in}{0.027778in}}{%
\pgfpathmoveto{\pgfqpoint{0.000000in}{0.000000in}}%
\pgfpathlineto{\pgfqpoint{0.000000in}{0.027778in}}%
\pgfusepath{stroke,fill}%
}%
\begin{pgfscope}%
\pgfsys@transformshift{6.388053in}{0.456901in}%
\pgfsys@useobject{currentmarker}{}%
\end{pgfscope}%
\end{pgfscope}%
\begin{pgfscope}%
\definecolor{textcolor}{rgb}{0.000000,0.000000,0.000000}%
\pgfsetstrokecolor{textcolor}%
\pgfsetfillcolor{textcolor}%
\pgftext[x=3.582936in,y=0.201367in,,top]{\color{textcolor}{\rmfamily\fontsize{12.000000}{14.400000}\selectfont\catcode`\^=\active\def^{\ifmmode\sp\else\^{}\fi}\catcode`\%=\active\def%{\%}$k$}}%
\end{pgfscope}%
\begin{pgfscope}%
\pgfpathrectangle{\pgfqpoint{0.725873in}{0.456901in}}{\pgfqpoint{5.714127in}{3.175238in}}%
\pgfusepath{clip}%
\pgfsetbuttcap%
\pgfsetroundjoin%
\pgfsetlinewidth{0.501875pt}%
\definecolor{currentstroke}{rgb}{0.501961,0.501961,0.501961}%
\pgfsetstrokecolor{currentstroke}%
\pgfsetstrokeopacity{0.400000}%
\pgfsetdash{{1.850000pt}{0.800000pt}}{0.000000pt}%
\pgfpathmoveto{\pgfqpoint{0.725873in}{0.456901in}}%
\pgfpathlineto{\pgfqpoint{6.440000in}{0.456901in}}%
\pgfusepath{stroke}%
\end{pgfscope}%
\begin{pgfscope}%
\pgfsetbuttcap%
\pgfsetroundjoin%
\definecolor{currentfill}{rgb}{1.000000,1.000000,1.000000}%
\pgfsetfillcolor{currentfill}%
\pgfsetlinewidth{0.803000pt}%
\definecolor{currentstroke}{rgb}{0.000000,0.000000,0.000000}%
\pgfsetstrokecolor{currentstroke}%
\pgfsetdash{}{0pt}%
\pgfsys@defobject{currentmarker}{\pgfqpoint{0.000000in}{0.000000in}}{\pgfqpoint{0.048611in}{0.000000in}}{%
\pgfpathmoveto{\pgfqpoint{0.000000in}{0.000000in}}%
\pgfpathlineto{\pgfqpoint{0.048611in}{0.000000in}}%
\pgfusepath{stroke,fill}%
}%
\begin{pgfscope}%
\pgfsys@transformshift{0.725873in}{0.456901in}%
\pgfsys@useobject{currentmarker}{}%
\end{pgfscope}%
\end{pgfscope}%
\begin{pgfscope}%
\definecolor{textcolor}{rgb}{0.000000,0.000000,0.000000}%
\pgfsetstrokecolor{textcolor}%
\pgfsetfillcolor{textcolor}%
\pgftext[x=0.257511in, y=0.399040in, left, base]{\color{textcolor}{\rmfamily\fontsize{12.000000}{14.400000}\selectfont\catcode`\^=\active\def^{\ifmmode\sp\else\^{}\fi}\catcode`\%=\active\def%{\%}$\mathdefault{\ensuremath{-}10.0}$}}%
\end{pgfscope}%
\begin{pgfscope}%
\pgfpathrectangle{\pgfqpoint{0.725873in}{0.456901in}}{\pgfqpoint{5.714127in}{3.175238in}}%
\pgfusepath{clip}%
\pgfsetbuttcap%
\pgfsetroundjoin%
\pgfsetlinewidth{0.501875pt}%
\definecolor{currentstroke}{rgb}{0.501961,0.501961,0.501961}%
\pgfsetstrokecolor{currentstroke}%
\pgfsetstrokeopacity{0.400000}%
\pgfsetdash{{1.850000pt}{0.800000pt}}{0.000000pt}%
\pgfpathmoveto{\pgfqpoint{0.725873in}{0.853806in}}%
\pgfpathlineto{\pgfqpoint{6.440000in}{0.853806in}}%
\pgfusepath{stroke}%
\end{pgfscope}%
\begin{pgfscope}%
\pgfsetbuttcap%
\pgfsetroundjoin%
\definecolor{currentfill}{rgb}{1.000000,1.000000,1.000000}%
\pgfsetfillcolor{currentfill}%
\pgfsetlinewidth{0.803000pt}%
\definecolor{currentstroke}{rgb}{0.000000,0.000000,0.000000}%
\pgfsetstrokecolor{currentstroke}%
\pgfsetdash{}{0pt}%
\pgfsys@defobject{currentmarker}{\pgfqpoint{0.000000in}{0.000000in}}{\pgfqpoint{0.048611in}{0.000000in}}{%
\pgfpathmoveto{\pgfqpoint{0.000000in}{0.000000in}}%
\pgfpathlineto{\pgfqpoint{0.048611in}{0.000000in}}%
\pgfusepath{stroke,fill}%
}%
\begin{pgfscope}%
\pgfsys@transformshift{0.725873in}{0.853806in}%
\pgfsys@useobject{currentmarker}{}%
\end{pgfscope}%
\end{pgfscope}%
\begin{pgfscope}%
\definecolor{textcolor}{rgb}{0.000000,0.000000,0.000000}%
\pgfsetstrokecolor{textcolor}%
\pgfsetfillcolor{textcolor}%
\pgftext[x=0.339108in, y=0.795944in, left, base]{\color{textcolor}{\rmfamily\fontsize{12.000000}{14.400000}\selectfont\catcode`\^=\active\def^{\ifmmode\sp\else\^{}\fi}\catcode`\%=\active\def%{\%}$\mathdefault{\ensuremath{-}7.5}$}}%
\end{pgfscope}%
\begin{pgfscope}%
\pgfpathrectangle{\pgfqpoint{0.725873in}{0.456901in}}{\pgfqpoint{5.714127in}{3.175238in}}%
\pgfusepath{clip}%
\pgfsetbuttcap%
\pgfsetroundjoin%
\pgfsetlinewidth{0.501875pt}%
\definecolor{currentstroke}{rgb}{0.501961,0.501961,0.501961}%
\pgfsetstrokecolor{currentstroke}%
\pgfsetstrokeopacity{0.400000}%
\pgfsetdash{{1.850000pt}{0.800000pt}}{0.000000pt}%
\pgfpathmoveto{\pgfqpoint{0.725873in}{1.250710in}}%
\pgfpathlineto{\pgfqpoint{6.440000in}{1.250710in}}%
\pgfusepath{stroke}%
\end{pgfscope}%
\begin{pgfscope}%
\pgfsetbuttcap%
\pgfsetroundjoin%
\definecolor{currentfill}{rgb}{1.000000,1.000000,1.000000}%
\pgfsetfillcolor{currentfill}%
\pgfsetlinewidth{0.803000pt}%
\definecolor{currentstroke}{rgb}{0.000000,0.000000,0.000000}%
\pgfsetstrokecolor{currentstroke}%
\pgfsetdash{}{0pt}%
\pgfsys@defobject{currentmarker}{\pgfqpoint{0.000000in}{0.000000in}}{\pgfqpoint{0.048611in}{0.000000in}}{%
\pgfpathmoveto{\pgfqpoint{0.000000in}{0.000000in}}%
\pgfpathlineto{\pgfqpoint{0.048611in}{0.000000in}}%
\pgfusepath{stroke,fill}%
}%
\begin{pgfscope}%
\pgfsys@transformshift{0.725873in}{1.250710in}%
\pgfsys@useobject{currentmarker}{}%
\end{pgfscope}%
\end{pgfscope}%
\begin{pgfscope}%
\definecolor{textcolor}{rgb}{0.000000,0.000000,0.000000}%
\pgfsetstrokecolor{textcolor}%
\pgfsetfillcolor{textcolor}%
\pgftext[x=0.339108in, y=1.192849in, left, base]{\color{textcolor}{\rmfamily\fontsize{12.000000}{14.400000}\selectfont\catcode`\^=\active\def^{\ifmmode\sp\else\^{}\fi}\catcode`\%=\active\def%{\%}$\mathdefault{\ensuremath{-}5.0}$}}%
\end{pgfscope}%
\begin{pgfscope}%
\pgfpathrectangle{\pgfqpoint{0.725873in}{0.456901in}}{\pgfqpoint{5.714127in}{3.175238in}}%
\pgfusepath{clip}%
\pgfsetbuttcap%
\pgfsetroundjoin%
\pgfsetlinewidth{0.501875pt}%
\definecolor{currentstroke}{rgb}{0.501961,0.501961,0.501961}%
\pgfsetstrokecolor{currentstroke}%
\pgfsetstrokeopacity{0.400000}%
\pgfsetdash{{1.850000pt}{0.800000pt}}{0.000000pt}%
\pgfpathmoveto{\pgfqpoint{0.725873in}{1.647615in}}%
\pgfpathlineto{\pgfqpoint{6.440000in}{1.647615in}}%
\pgfusepath{stroke}%
\end{pgfscope}%
\begin{pgfscope}%
\pgfsetbuttcap%
\pgfsetroundjoin%
\definecolor{currentfill}{rgb}{1.000000,1.000000,1.000000}%
\pgfsetfillcolor{currentfill}%
\pgfsetlinewidth{0.803000pt}%
\definecolor{currentstroke}{rgb}{0.000000,0.000000,0.000000}%
\pgfsetstrokecolor{currentstroke}%
\pgfsetdash{}{0pt}%
\pgfsys@defobject{currentmarker}{\pgfqpoint{0.000000in}{0.000000in}}{\pgfqpoint{0.048611in}{0.000000in}}{%
\pgfpathmoveto{\pgfqpoint{0.000000in}{0.000000in}}%
\pgfpathlineto{\pgfqpoint{0.048611in}{0.000000in}}%
\pgfusepath{stroke,fill}%
}%
\begin{pgfscope}%
\pgfsys@transformshift{0.725873in}{1.647615in}%
\pgfsys@useobject{currentmarker}{}%
\end{pgfscope}%
\end{pgfscope}%
\begin{pgfscope}%
\definecolor{textcolor}{rgb}{0.000000,0.000000,0.000000}%
\pgfsetstrokecolor{textcolor}%
\pgfsetfillcolor{textcolor}%
\pgftext[x=0.339108in, y=1.589754in, left, base]{\color{textcolor}{\rmfamily\fontsize{12.000000}{14.400000}\selectfont\catcode`\^=\active\def^{\ifmmode\sp\else\^{}\fi}\catcode`\%=\active\def%{\%}$\mathdefault{\ensuremath{-}2.5}$}}%
\end{pgfscope}%
\begin{pgfscope}%
\pgfpathrectangle{\pgfqpoint{0.725873in}{0.456901in}}{\pgfqpoint{5.714127in}{3.175238in}}%
\pgfusepath{clip}%
\pgfsetbuttcap%
\pgfsetroundjoin%
\pgfsetlinewidth{0.501875pt}%
\definecolor{currentstroke}{rgb}{0.501961,0.501961,0.501961}%
\pgfsetstrokecolor{currentstroke}%
\pgfsetstrokeopacity{0.400000}%
\pgfsetdash{{1.850000pt}{0.800000pt}}{0.000000pt}%
\pgfpathmoveto{\pgfqpoint{0.725873in}{2.044520in}}%
\pgfpathlineto{\pgfqpoint{6.440000in}{2.044520in}}%
\pgfusepath{stroke}%
\end{pgfscope}%
\begin{pgfscope}%
\pgfsetbuttcap%
\pgfsetroundjoin%
\definecolor{currentfill}{rgb}{1.000000,1.000000,1.000000}%
\pgfsetfillcolor{currentfill}%
\pgfsetlinewidth{0.803000pt}%
\definecolor{currentstroke}{rgb}{0.000000,0.000000,0.000000}%
\pgfsetstrokecolor{currentstroke}%
\pgfsetdash{}{0pt}%
\pgfsys@defobject{currentmarker}{\pgfqpoint{0.000000in}{0.000000in}}{\pgfqpoint{0.048611in}{0.000000in}}{%
\pgfpathmoveto{\pgfqpoint{0.000000in}{0.000000in}}%
\pgfpathlineto{\pgfqpoint{0.048611in}{0.000000in}}%
\pgfusepath{stroke,fill}%
}%
\begin{pgfscope}%
\pgfsys@transformshift{0.725873in}{2.044520in}%
\pgfsys@useobject{currentmarker}{}%
\end{pgfscope}%
\end{pgfscope}%
\begin{pgfscope}%
\definecolor{textcolor}{rgb}{0.000000,0.000000,0.000000}%
\pgfsetstrokecolor{textcolor}%
\pgfsetfillcolor{textcolor}%
\pgftext[x=0.468738in, y=1.986658in, left, base]{\color{textcolor}{\rmfamily\fontsize{12.000000}{14.400000}\selectfont\catcode`\^=\active\def^{\ifmmode\sp\else\^{}\fi}\catcode`\%=\active\def%{\%}$\mathdefault{0.0}$}}%
\end{pgfscope}%
\begin{pgfscope}%
\pgfpathrectangle{\pgfqpoint{0.725873in}{0.456901in}}{\pgfqpoint{5.714127in}{3.175238in}}%
\pgfusepath{clip}%
\pgfsetbuttcap%
\pgfsetroundjoin%
\pgfsetlinewidth{0.501875pt}%
\definecolor{currentstroke}{rgb}{0.501961,0.501961,0.501961}%
\pgfsetstrokecolor{currentstroke}%
\pgfsetstrokeopacity{0.400000}%
\pgfsetdash{{1.850000pt}{0.800000pt}}{0.000000pt}%
\pgfpathmoveto{\pgfqpoint{0.725873in}{2.441424in}}%
\pgfpathlineto{\pgfqpoint{6.440000in}{2.441424in}}%
\pgfusepath{stroke}%
\end{pgfscope}%
\begin{pgfscope}%
\pgfsetbuttcap%
\pgfsetroundjoin%
\definecolor{currentfill}{rgb}{1.000000,1.000000,1.000000}%
\pgfsetfillcolor{currentfill}%
\pgfsetlinewidth{0.803000pt}%
\definecolor{currentstroke}{rgb}{0.000000,0.000000,0.000000}%
\pgfsetstrokecolor{currentstroke}%
\pgfsetdash{}{0pt}%
\pgfsys@defobject{currentmarker}{\pgfqpoint{0.000000in}{0.000000in}}{\pgfqpoint{0.048611in}{0.000000in}}{%
\pgfpathmoveto{\pgfqpoint{0.000000in}{0.000000in}}%
\pgfpathlineto{\pgfqpoint{0.048611in}{0.000000in}}%
\pgfusepath{stroke,fill}%
}%
\begin{pgfscope}%
\pgfsys@transformshift{0.725873in}{2.441424in}%
\pgfsys@useobject{currentmarker}{}%
\end{pgfscope}%
\end{pgfscope}%
\begin{pgfscope}%
\definecolor{textcolor}{rgb}{0.000000,0.000000,0.000000}%
\pgfsetstrokecolor{textcolor}%
\pgfsetfillcolor{textcolor}%
\pgftext[x=0.468738in, y=2.383563in, left, base]{\color{textcolor}{\rmfamily\fontsize{12.000000}{14.400000}\selectfont\catcode`\^=\active\def^{\ifmmode\sp\else\^{}\fi}\catcode`\%=\active\def%{\%}$\mathdefault{2.5}$}}%
\end{pgfscope}%
\begin{pgfscope}%
\pgfpathrectangle{\pgfqpoint{0.725873in}{0.456901in}}{\pgfqpoint{5.714127in}{3.175238in}}%
\pgfusepath{clip}%
\pgfsetbuttcap%
\pgfsetroundjoin%
\pgfsetlinewidth{0.501875pt}%
\definecolor{currentstroke}{rgb}{0.501961,0.501961,0.501961}%
\pgfsetstrokecolor{currentstroke}%
\pgfsetstrokeopacity{0.400000}%
\pgfsetdash{{1.850000pt}{0.800000pt}}{0.000000pt}%
\pgfpathmoveto{\pgfqpoint{0.725873in}{2.838329in}}%
\pgfpathlineto{\pgfqpoint{6.440000in}{2.838329in}}%
\pgfusepath{stroke}%
\end{pgfscope}%
\begin{pgfscope}%
\pgfsetbuttcap%
\pgfsetroundjoin%
\definecolor{currentfill}{rgb}{1.000000,1.000000,1.000000}%
\pgfsetfillcolor{currentfill}%
\pgfsetlinewidth{0.803000pt}%
\definecolor{currentstroke}{rgb}{0.000000,0.000000,0.000000}%
\pgfsetstrokecolor{currentstroke}%
\pgfsetdash{}{0pt}%
\pgfsys@defobject{currentmarker}{\pgfqpoint{0.000000in}{0.000000in}}{\pgfqpoint{0.048611in}{0.000000in}}{%
\pgfpathmoveto{\pgfqpoint{0.000000in}{0.000000in}}%
\pgfpathlineto{\pgfqpoint{0.048611in}{0.000000in}}%
\pgfusepath{stroke,fill}%
}%
\begin{pgfscope}%
\pgfsys@transformshift{0.725873in}{2.838329in}%
\pgfsys@useobject{currentmarker}{}%
\end{pgfscope}%
\end{pgfscope}%
\begin{pgfscope}%
\definecolor{textcolor}{rgb}{0.000000,0.000000,0.000000}%
\pgfsetstrokecolor{textcolor}%
\pgfsetfillcolor{textcolor}%
\pgftext[x=0.468738in, y=2.780468in, left, base]{\color{textcolor}{\rmfamily\fontsize{12.000000}{14.400000}\selectfont\catcode`\^=\active\def^{\ifmmode\sp\else\^{}\fi}\catcode`\%=\active\def%{\%}$\mathdefault{5.0}$}}%
\end{pgfscope}%
\begin{pgfscope}%
\pgfpathrectangle{\pgfqpoint{0.725873in}{0.456901in}}{\pgfqpoint{5.714127in}{3.175238in}}%
\pgfusepath{clip}%
\pgfsetbuttcap%
\pgfsetroundjoin%
\pgfsetlinewidth{0.501875pt}%
\definecolor{currentstroke}{rgb}{0.501961,0.501961,0.501961}%
\pgfsetstrokecolor{currentstroke}%
\pgfsetstrokeopacity{0.400000}%
\pgfsetdash{{1.850000pt}{0.800000pt}}{0.000000pt}%
\pgfpathmoveto{\pgfqpoint{0.725873in}{3.235234in}}%
\pgfpathlineto{\pgfqpoint{6.440000in}{3.235234in}}%
\pgfusepath{stroke}%
\end{pgfscope}%
\begin{pgfscope}%
\pgfsetbuttcap%
\pgfsetroundjoin%
\definecolor{currentfill}{rgb}{1.000000,1.000000,1.000000}%
\pgfsetfillcolor{currentfill}%
\pgfsetlinewidth{0.803000pt}%
\definecolor{currentstroke}{rgb}{0.000000,0.000000,0.000000}%
\pgfsetstrokecolor{currentstroke}%
\pgfsetdash{}{0pt}%
\pgfsys@defobject{currentmarker}{\pgfqpoint{0.000000in}{0.000000in}}{\pgfqpoint{0.048611in}{0.000000in}}{%
\pgfpathmoveto{\pgfqpoint{0.000000in}{0.000000in}}%
\pgfpathlineto{\pgfqpoint{0.048611in}{0.000000in}}%
\pgfusepath{stroke,fill}%
}%
\begin{pgfscope}%
\pgfsys@transformshift{0.725873in}{3.235234in}%
\pgfsys@useobject{currentmarker}{}%
\end{pgfscope}%
\end{pgfscope}%
\begin{pgfscope}%
\definecolor{textcolor}{rgb}{0.000000,0.000000,0.000000}%
\pgfsetstrokecolor{textcolor}%
\pgfsetfillcolor{textcolor}%
\pgftext[x=0.468738in, y=3.177373in, left, base]{\color{textcolor}{\rmfamily\fontsize{12.000000}{14.400000}\selectfont\catcode`\^=\active\def^{\ifmmode\sp\else\^{}\fi}\catcode`\%=\active\def%{\%}$\mathdefault{7.5}$}}%
\end{pgfscope}%
\begin{pgfscope}%
\pgfpathrectangle{\pgfqpoint{0.725873in}{0.456901in}}{\pgfqpoint{5.714127in}{3.175238in}}%
\pgfusepath{clip}%
\pgfsetbuttcap%
\pgfsetroundjoin%
\pgfsetlinewidth{0.501875pt}%
\definecolor{currentstroke}{rgb}{0.501961,0.501961,0.501961}%
\pgfsetstrokecolor{currentstroke}%
\pgfsetstrokeopacity{0.400000}%
\pgfsetdash{{1.850000pt}{0.800000pt}}{0.000000pt}%
\pgfpathmoveto{\pgfqpoint{0.725873in}{3.632139in}}%
\pgfpathlineto{\pgfqpoint{6.440000in}{3.632139in}}%
\pgfusepath{stroke}%
\end{pgfscope}%
\begin{pgfscope}%
\pgfsetbuttcap%
\pgfsetroundjoin%
\definecolor{currentfill}{rgb}{1.000000,1.000000,1.000000}%
\pgfsetfillcolor{currentfill}%
\pgfsetlinewidth{0.803000pt}%
\definecolor{currentstroke}{rgb}{0.000000,0.000000,0.000000}%
\pgfsetstrokecolor{currentstroke}%
\pgfsetdash{}{0pt}%
\pgfsys@defobject{currentmarker}{\pgfqpoint{0.000000in}{0.000000in}}{\pgfqpoint{0.048611in}{0.000000in}}{%
\pgfpathmoveto{\pgfqpoint{0.000000in}{0.000000in}}%
\pgfpathlineto{\pgfqpoint{0.048611in}{0.000000in}}%
\pgfusepath{stroke,fill}%
}%
\begin{pgfscope}%
\pgfsys@transformshift{0.725873in}{3.632139in}%
\pgfsys@useobject{currentmarker}{}%
\end{pgfscope}%
\end{pgfscope}%
\begin{pgfscope}%
\definecolor{textcolor}{rgb}{0.000000,0.000000,0.000000}%
\pgfsetstrokecolor{textcolor}%
\pgfsetfillcolor{textcolor}%
\pgftext[x=0.387141in, y=3.574277in, left, base]{\color{textcolor}{\rmfamily\fontsize{12.000000}{14.400000}\selectfont\catcode`\^=\active\def^{\ifmmode\sp\else\^{}\fi}\catcode`\%=\active\def%{\%}$\mathdefault{10.0}$}}%
\end{pgfscope}%
\begin{pgfscope}%
\definecolor{textcolor}{rgb}{0.000000,0.000000,0.000000}%
\pgfsetstrokecolor{textcolor}%
\pgfsetfillcolor{textcolor}%
\pgftext[x=0.201956in,y=2.044520in,,bottom,rotate=90.000000]{\color{textcolor}{\rmfamily\fontsize{12.000000}{14.400000}\selectfont\catcode`\^=\active\def^{\ifmmode\sp\else\^{}\fi}\catcode`\%=\active\def%{\%}$\Phi_{1k}, ^\circ$}}%
\end{pgfscope}%
\begin{pgfscope}%
\pgfpathrectangle{\pgfqpoint{0.725873in}{0.456901in}}{\pgfqpoint{5.714127in}{3.175238in}}%
\pgfusepath{clip}%
\pgfsetbuttcap%
\pgfsetroundjoin%
\pgfsetlinewidth{2.007500pt}%
\definecolor{currentstroke}{rgb}{0.000000,0.000000,0.000000}%
\pgfsetstrokecolor{currentstroke}%
\pgfsetdash{{7.400000pt}{3.200000pt}}{0.000000pt}%
\pgfpathmoveto{\pgfqpoint{0.985606in}{2.044520in}}%
\pgfpathlineto{\pgfqpoint{0.985606in}{2.044520in}}%
\pgfusepath{stroke}%
\end{pgfscope}%
\begin{pgfscope}%
\pgfpathrectangle{\pgfqpoint{0.725873in}{0.456901in}}{\pgfqpoint{5.714127in}{3.175238in}}%
\pgfusepath{clip}%
\pgfsetbuttcap%
\pgfsetroundjoin%
\pgfsetlinewidth{2.007500pt}%
\definecolor{currentstroke}{rgb}{0.000000,0.000000,0.000000}%
\pgfsetstrokecolor{currentstroke}%
\pgfsetdash{{7.400000pt}{3.200000pt}}{0.000000pt}%
\pgfpathmoveto{\pgfqpoint{2.024538in}{2.044520in}}%
\pgfpathlineto{\pgfqpoint{2.024538in}{2.044520in}}%
\pgfusepath{stroke}%
\end{pgfscope}%
\begin{pgfscope}%
\pgfpathrectangle{\pgfqpoint{0.725873in}{0.456901in}}{\pgfqpoint{5.714127in}{3.175238in}}%
\pgfusepath{clip}%
\pgfsetbuttcap%
\pgfsetroundjoin%
\pgfsetlinewidth{2.007500pt}%
\definecolor{currentstroke}{rgb}{0.000000,0.000000,0.000000}%
\pgfsetstrokecolor{currentstroke}%
\pgfsetdash{{7.400000pt}{3.200000pt}}{0.000000pt}%
\pgfpathmoveto{\pgfqpoint{4.102403in}{2.044520in}}%
\pgfpathlineto{\pgfqpoint{4.102403in}{2.044520in}}%
\pgfusepath{stroke}%
\end{pgfscope}%
\begin{pgfscope}%
\pgfpathrectangle{\pgfqpoint{0.725873in}{0.456901in}}{\pgfqpoint{5.714127in}{3.175238in}}%
\pgfusepath{clip}%
\pgfsetbuttcap%
\pgfsetroundjoin%
\pgfsetlinewidth{2.007500pt}%
\definecolor{currentstroke}{rgb}{0.000000,0.000000,0.000000}%
\pgfsetstrokecolor{currentstroke}%
\pgfsetdash{{7.400000pt}{3.200000pt}}{0.000000pt}%
\pgfpathmoveto{\pgfqpoint{6.180267in}{2.044520in}}%
\pgfpathlineto{\pgfqpoint{6.180267in}{2.044520in}}%
\pgfusepath{stroke}%
\end{pgfscope}%
\begin{pgfscope}%
\pgfpathrectangle{\pgfqpoint{0.725873in}{0.456901in}}{\pgfqpoint{5.714127in}{3.175238in}}%
\pgfusepath{clip}%
\pgfsetbuttcap%
\pgfsetroundjoin%
\definecolor{currentfill}{rgb}{1.000000,1.000000,1.000000}%
\pgfsetfillcolor{currentfill}%
\pgfsetlinewidth{1.505625pt}%
\definecolor{currentstroke}{rgb}{0.000000,0.000000,0.000000}%
\pgfsetstrokecolor{currentstroke}%
\pgfsetdash{}{0pt}%
\pgfsys@defobject{currentmarker}{\pgfqpoint{-0.027778in}{-0.027778in}}{\pgfqpoint{0.027778in}{0.027778in}}{%
\pgfpathmoveto{\pgfqpoint{0.000000in}{-0.027778in}}%
\pgfpathcurveto{\pgfqpoint{0.007367in}{-0.027778in}}{\pgfqpoint{0.014433in}{-0.024851in}}{\pgfqpoint{0.019642in}{-0.019642in}}%
\pgfpathcurveto{\pgfqpoint{0.024851in}{-0.014433in}}{\pgfqpoint{0.027778in}{-0.007367in}}{\pgfqpoint{0.027778in}{0.000000in}}%
\pgfpathcurveto{\pgfqpoint{0.027778in}{0.007367in}}{\pgfqpoint{0.024851in}{0.014433in}}{\pgfqpoint{0.019642in}{0.019642in}}%
\pgfpathcurveto{\pgfqpoint{0.014433in}{0.024851in}}{\pgfqpoint{0.007367in}{0.027778in}}{\pgfqpoint{0.000000in}{0.027778in}}%
\pgfpathcurveto{\pgfqpoint{-0.007367in}{0.027778in}}{\pgfqpoint{-0.014433in}{0.024851in}}{\pgfqpoint{-0.019642in}{0.019642in}}%
\pgfpathcurveto{\pgfqpoint{-0.024851in}{0.014433in}}{\pgfqpoint{-0.027778in}{0.007367in}}{\pgfqpoint{-0.027778in}{0.000000in}}%
\pgfpathcurveto{\pgfqpoint{-0.027778in}{-0.007367in}}{\pgfqpoint{-0.024851in}{-0.014433in}}{\pgfqpoint{-0.019642in}{-0.019642in}}%
\pgfpathcurveto{\pgfqpoint{-0.014433in}{-0.024851in}}{\pgfqpoint{-0.007367in}{-0.027778in}}{\pgfqpoint{0.000000in}{-0.027778in}}%
\pgfpathlineto{\pgfqpoint{0.000000in}{-0.027778in}}%
\pgfpathclose%
\pgfusepath{stroke,fill}%
}%
\begin{pgfscope}%
\pgfsys@transformshift{0.985606in}{2.044520in}%
\pgfsys@useobject{currentmarker}{}%
\end{pgfscope}%
\begin{pgfscope}%
\pgfsys@transformshift{2.024538in}{2.044520in}%
\pgfsys@useobject{currentmarker}{}%
\end{pgfscope}%
\begin{pgfscope}%
\pgfsys@transformshift{4.102403in}{2.044520in}%
\pgfsys@useobject{currentmarker}{}%
\end{pgfscope}%
\begin{pgfscope}%
\pgfsys@transformshift{6.180267in}{2.044520in}%
\pgfsys@useobject{currentmarker}{}%
\end{pgfscope}%
\end{pgfscope}%
\begin{pgfscope}%
\pgfpathrectangle{\pgfqpoint{0.725873in}{0.456901in}}{\pgfqpoint{5.714127in}{3.175238in}}%
\pgfusepath{clip}%
\pgfsetrectcap%
\pgfsetroundjoin%
\pgfsetlinewidth{2.007500pt}%
\definecolor{currentstroke}{rgb}{0.000000,0.000000,0.000000}%
\pgfsetstrokecolor{currentstroke}%
\pgfsetdash{}{0pt}%
\pgfpathmoveto{\pgfqpoint{0.985606in}{2.044520in}}%
\pgfpathlineto{\pgfqpoint{6.180267in}{2.044520in}}%
\pgfusepath{stroke}%
\end{pgfscope}%
\begin{pgfscope}%
\pgfsetrectcap%
\pgfsetmiterjoin%
\pgfsetlinewidth{0.602250pt}%
\definecolor{currentstroke}{rgb}{0.000000,0.000000,0.000000}%
\pgfsetstrokecolor{currentstroke}%
\pgfsetdash{}{0pt}%
\pgfpathmoveto{\pgfqpoint{0.725873in}{0.456901in}}%
\pgfpathlineto{\pgfqpoint{0.725873in}{3.632139in}}%
\pgfusepath{stroke}%
\end{pgfscope}%
\begin{pgfscope}%
\pgfsetrectcap%
\pgfsetmiterjoin%
\pgfsetlinewidth{0.602250pt}%
\definecolor{currentstroke}{rgb}{0.000000,0.000000,0.000000}%
\pgfsetstrokecolor{currentstroke}%
\pgfsetdash{}{0pt}%
\pgfpathmoveto{\pgfqpoint{0.725873in}{0.456901in}}%
\pgfpathlineto{\pgfqpoint{6.440000in}{0.456901in}}%
\pgfusepath{stroke}%
\end{pgfscope}%
\end{pgfpicture}%
\makeatother%
\endgroup%

    \caption{Дискретный фазовый спектр входного сигнала}
    \label{fig:plot_disc_spec_Phi}
\end{figure}

% Отрезок ряда Фурье, аппроксимирующий входной периодический сигнал 
при $N=\pv[.0]{ex10['N']}$ (учитываются гармоники $k=1, 3, 5$)
% , имеет вид:
\begin{equation}
    i_1(t) \approx \pv[.3]{ex10['A1k'][1]} \cos(\pv[.3]{ex10['w1']}t) + \pv[.3]{ex10['A1k'][2]} \cos(\pv[.3]{3*ex10['w1']}t) + \pv[.3]{ex10['A1k'][3]} \cos(\pv[.3]{5*ex10['w1']}t)
\end{equation}

% На \cref{fig:plot_fourier_approx} представлен график исходного периодического сигнала и его аппроксимация отрезком ряда Фурье. Тонкими линиями показаны отдельные гармоники, составляющие аппроксимацию.

\begin{figure}[H]
    \centering
    %% Creator: Matplotlib, PGF backend
%%
%% To include the figure in your LaTeX document, write
%%   \input{<filename>.pgf}
%%
%% Make sure the required packages are loaded in your preamble
%%   \usepackage{pgf}
%%
%% Also ensure that all the required font packages are loaded; for instance,
%% the lmodern package is sometimes necessary when using math font.
%%   \usepackage{lmodern}
%%
%% Figures using additional raster images can only be included by \input if
%% they are in the same directory as the main LaTeX file. For loading figures
%% from other directories you can use the `import` package
%%   \usepackage{import}
%%
%% and then include the figures with
%%   \import{<path to file>}{<filename>.pgf}
%%
%% Matplotlib used the following preamble
%%   \def\mathdefault#1{#1}
%%   \everymath=\expandafter{\the\everymath\displaystyle}
%%   \IfFileExists{scrextend.sty}{
%%     \usepackage[fontsize=12.000000pt]{scrextend}
%%   }{
%%     \renewcommand{\normalsize}{\fontsize{12.000000}{14.400000}\selectfont}
%%     \normalsize
%%   }
%%   \usepackage{fontspec}\setmainfont{Times New Roman}\usepackage[english,russian]{babel}\usepackage{amsmath}
%%   \ifdefined\pdftexversion\else  % non-pdftex case.
%%     \usepackage{fontspec}
%%     \setmainfont{times.ttf}[Path=\detokenize{/usr/local/share/fonts/truetype/times-new-roman/}]
%%     \setsansfont{DejaVuSans.ttf}[Path=\detokenize{/usr/share/matplotlib/mpl-data/fonts/ttf/}]
%%     \setmonofont{DejaVuSansMono.ttf}[Path=\detokenize{/usr/share/matplotlib/mpl-data/fonts/ttf/}]
%%   \fi
%%   \makeatletter\@ifpackageloaded{underscore}{}{\usepackage[strings]{underscore}}\makeatother
%%
\begingroup%
\makeatletter%
\begin{pgfpicture}%
\pgfpathrectangle{\pgfpointorigin}{\pgfqpoint{6.488011in}{3.740000in}}%
\pgfusepath{use as bounding box, clip}%
\begin{pgfscope}%
\pgfsetbuttcap%
\pgfsetmiterjoin%
\definecolor{currentfill}{rgb}{1.000000,1.000000,1.000000}%
\pgfsetfillcolor{currentfill}%
\pgfsetlinewidth{0.000000pt}%
\definecolor{currentstroke}{rgb}{1.000000,1.000000,1.000000}%
\pgfsetstrokecolor{currentstroke}%
\pgfsetdash{}{0pt}%
\pgfpathmoveto{\pgfqpoint{0.000000in}{0.000000in}}%
\pgfpathlineto{\pgfqpoint{6.488011in}{0.000000in}}%
\pgfpathlineto{\pgfqpoint{6.488011in}{3.740000in}}%
\pgfpathlineto{\pgfqpoint{0.000000in}{3.740000in}}%
\pgfpathlineto{\pgfqpoint{0.000000in}{0.000000in}}%
\pgfpathclose%
\pgfusepath{fill}%
\end{pgfscope}%
\begin{pgfscope}%
\pgfsetbuttcap%
\pgfsetmiterjoin%
\definecolor{currentfill}{rgb}{1.000000,1.000000,1.000000}%
\pgfsetfillcolor{currentfill}%
\pgfsetlinewidth{0.000000pt}%
\definecolor{currentstroke}{rgb}{0.000000,0.000000,0.000000}%
\pgfsetstrokecolor{currentstroke}%
\pgfsetstrokeopacity{0.000000}%
\pgfsetdash{}{0pt}%
\pgfpathmoveto{\pgfqpoint{0.532060in}{0.456901in}}%
\pgfpathlineto{\pgfqpoint{6.335780in}{0.456901in}}%
\pgfpathlineto{\pgfqpoint{6.335780in}{3.273333in}}%
\pgfpathlineto{\pgfqpoint{0.532060in}{3.273333in}}%
\pgfpathlineto{\pgfqpoint{0.532060in}{0.456901in}}%
\pgfpathclose%
\pgfusepath{fill}%
\end{pgfscope}%
\begin{pgfscope}%
\pgfpathrectangle{\pgfqpoint{0.532060in}{0.456901in}}{\pgfqpoint{5.803720in}{2.816433in}}%
\pgfusepath{clip}%
\pgfsetbuttcap%
\pgfsetroundjoin%
\pgfsetlinewidth{0.501875pt}%
\definecolor{currentstroke}{rgb}{0.501961,0.501961,0.501961}%
\pgfsetstrokecolor{currentstroke}%
\pgfsetstrokeopacity{0.400000}%
\pgfsetdash{{1.850000pt}{0.800000pt}}{0.000000pt}%
\pgfpathmoveto{\pgfqpoint{0.714651in}{0.456901in}}%
\pgfpathlineto{\pgfqpoint{0.714651in}{3.273333in}}%
\pgfusepath{stroke}%
\end{pgfscope}%
\begin{pgfscope}%
\pgfsetbuttcap%
\pgfsetroundjoin%
\definecolor{currentfill}{rgb}{1.000000,1.000000,1.000000}%
\pgfsetfillcolor{currentfill}%
\pgfsetlinewidth{0.803000pt}%
\definecolor{currentstroke}{rgb}{0.000000,0.000000,0.000000}%
\pgfsetstrokecolor{currentstroke}%
\pgfsetdash{}{0pt}%
\pgfsys@defobject{currentmarker}{\pgfqpoint{0.000000in}{0.000000in}}{\pgfqpoint{0.000000in}{0.048611in}}{%
\pgfpathmoveto{\pgfqpoint{0.000000in}{0.000000in}}%
\pgfpathlineto{\pgfqpoint{0.000000in}{0.048611in}}%
\pgfusepath{stroke,fill}%
}%
\begin{pgfscope}%
\pgfsys@transformshift{0.714651in}{0.456901in}%
\pgfsys@useobject{currentmarker}{}%
\end{pgfscope}%
\end{pgfscope}%
\begin{pgfscope}%
\definecolor{textcolor}{rgb}{0.000000,0.000000,0.000000}%
\pgfsetstrokecolor{textcolor}%
\pgfsetfillcolor{textcolor}%
\pgftext[x=0.714651in,y=0.408290in,,top]{\color{textcolor}{\rmfamily\fontsize{12.000000}{14.400000}\selectfont\catcode`\^=\active\def^{\ifmmode\sp\else\^{}\fi}\catcode`\%=\active\def%{\%}$\mathdefault{\ensuremath{-}50}$}}%
\end{pgfscope}%
\begin{pgfscope}%
\pgfpathrectangle{\pgfqpoint{0.532060in}{0.456901in}}{\pgfqpoint{5.803720in}{2.816433in}}%
\pgfusepath{clip}%
\pgfsetbuttcap%
\pgfsetroundjoin%
\pgfsetlinewidth{0.501875pt}%
\definecolor{currentstroke}{rgb}{0.501961,0.501961,0.501961}%
\pgfsetstrokecolor{currentstroke}%
\pgfsetstrokeopacity{0.400000}%
\pgfsetdash{{1.850000pt}{0.800000pt}}{0.000000pt}%
\pgfpathmoveto{\pgfqpoint{1.414772in}{0.456901in}}%
\pgfpathlineto{\pgfqpoint{1.414772in}{3.273333in}}%
\pgfusepath{stroke}%
\end{pgfscope}%
\begin{pgfscope}%
\pgfsetbuttcap%
\pgfsetroundjoin%
\definecolor{currentfill}{rgb}{1.000000,1.000000,1.000000}%
\pgfsetfillcolor{currentfill}%
\pgfsetlinewidth{0.803000pt}%
\definecolor{currentstroke}{rgb}{0.000000,0.000000,0.000000}%
\pgfsetstrokecolor{currentstroke}%
\pgfsetdash{}{0pt}%
\pgfsys@defobject{currentmarker}{\pgfqpoint{0.000000in}{0.000000in}}{\pgfqpoint{0.000000in}{0.048611in}}{%
\pgfpathmoveto{\pgfqpoint{0.000000in}{0.000000in}}%
\pgfpathlineto{\pgfqpoint{0.000000in}{0.048611in}}%
\pgfusepath{stroke,fill}%
}%
\begin{pgfscope}%
\pgfsys@transformshift{1.414772in}{0.456901in}%
\pgfsys@useobject{currentmarker}{}%
\end{pgfscope}%
\end{pgfscope}%
\begin{pgfscope}%
\definecolor{textcolor}{rgb}{0.000000,0.000000,0.000000}%
\pgfsetstrokecolor{textcolor}%
\pgfsetfillcolor{textcolor}%
\pgftext[x=1.414772in,y=0.408290in,,top]{\color{textcolor}{\rmfamily\fontsize{12.000000}{14.400000}\selectfont\catcode`\^=\active\def^{\ifmmode\sp\else\^{}\fi}\catcode`\%=\active\def%{\%}$\mathdefault{\ensuremath{-}25}$}}%
\end{pgfscope}%
\begin{pgfscope}%
\pgfpathrectangle{\pgfqpoint{0.532060in}{0.456901in}}{\pgfqpoint{5.803720in}{2.816433in}}%
\pgfusepath{clip}%
\pgfsetbuttcap%
\pgfsetroundjoin%
\pgfsetlinewidth{0.501875pt}%
\definecolor{currentstroke}{rgb}{0.501961,0.501961,0.501961}%
\pgfsetstrokecolor{currentstroke}%
\pgfsetstrokeopacity{0.400000}%
\pgfsetdash{{1.850000pt}{0.800000pt}}{0.000000pt}%
\pgfpathmoveto{\pgfqpoint{2.114893in}{0.456901in}}%
\pgfpathlineto{\pgfqpoint{2.114893in}{3.273333in}}%
\pgfusepath{stroke}%
\end{pgfscope}%
\begin{pgfscope}%
\pgfsetbuttcap%
\pgfsetroundjoin%
\definecolor{currentfill}{rgb}{1.000000,1.000000,1.000000}%
\pgfsetfillcolor{currentfill}%
\pgfsetlinewidth{0.803000pt}%
\definecolor{currentstroke}{rgb}{0.000000,0.000000,0.000000}%
\pgfsetstrokecolor{currentstroke}%
\pgfsetdash{}{0pt}%
\pgfsys@defobject{currentmarker}{\pgfqpoint{0.000000in}{0.000000in}}{\pgfqpoint{0.000000in}{0.048611in}}{%
\pgfpathmoveto{\pgfqpoint{0.000000in}{0.000000in}}%
\pgfpathlineto{\pgfqpoint{0.000000in}{0.048611in}}%
\pgfusepath{stroke,fill}%
}%
\begin{pgfscope}%
\pgfsys@transformshift{2.114893in}{0.456901in}%
\pgfsys@useobject{currentmarker}{}%
\end{pgfscope}%
\end{pgfscope}%
\begin{pgfscope}%
\definecolor{textcolor}{rgb}{0.000000,0.000000,0.000000}%
\pgfsetstrokecolor{textcolor}%
\pgfsetfillcolor{textcolor}%
\pgftext[x=2.114893in,y=0.408290in,,top]{\color{textcolor}{\rmfamily\fontsize{12.000000}{14.400000}\selectfont\catcode`\^=\active\def^{\ifmmode\sp\else\^{}\fi}\catcode`\%=\active\def%{\%}$\mathdefault{0}$}}%
\end{pgfscope}%
\begin{pgfscope}%
\pgfpathrectangle{\pgfqpoint{0.532060in}{0.456901in}}{\pgfqpoint{5.803720in}{2.816433in}}%
\pgfusepath{clip}%
\pgfsetbuttcap%
\pgfsetroundjoin%
\pgfsetlinewidth{0.501875pt}%
\definecolor{currentstroke}{rgb}{0.501961,0.501961,0.501961}%
\pgfsetstrokecolor{currentstroke}%
\pgfsetstrokeopacity{0.400000}%
\pgfsetdash{{1.850000pt}{0.800000pt}}{0.000000pt}%
\pgfpathmoveto{\pgfqpoint{2.815013in}{0.456901in}}%
\pgfpathlineto{\pgfqpoint{2.815013in}{3.273333in}}%
\pgfusepath{stroke}%
\end{pgfscope}%
\begin{pgfscope}%
\pgfsetbuttcap%
\pgfsetroundjoin%
\definecolor{currentfill}{rgb}{1.000000,1.000000,1.000000}%
\pgfsetfillcolor{currentfill}%
\pgfsetlinewidth{0.803000pt}%
\definecolor{currentstroke}{rgb}{0.000000,0.000000,0.000000}%
\pgfsetstrokecolor{currentstroke}%
\pgfsetdash{}{0pt}%
\pgfsys@defobject{currentmarker}{\pgfqpoint{0.000000in}{0.000000in}}{\pgfqpoint{0.000000in}{0.048611in}}{%
\pgfpathmoveto{\pgfqpoint{0.000000in}{0.000000in}}%
\pgfpathlineto{\pgfqpoint{0.000000in}{0.048611in}}%
\pgfusepath{stroke,fill}%
}%
\begin{pgfscope}%
\pgfsys@transformshift{2.815013in}{0.456901in}%
\pgfsys@useobject{currentmarker}{}%
\end{pgfscope}%
\end{pgfscope}%
\begin{pgfscope}%
\definecolor{textcolor}{rgb}{0.000000,0.000000,0.000000}%
\pgfsetstrokecolor{textcolor}%
\pgfsetfillcolor{textcolor}%
\pgftext[x=2.815013in,y=0.408290in,,top]{\color{textcolor}{\rmfamily\fontsize{12.000000}{14.400000}\selectfont\catcode`\^=\active\def^{\ifmmode\sp\else\^{}\fi}\catcode`\%=\active\def%{\%}$\mathdefault{25}$}}%
\end{pgfscope}%
\begin{pgfscope}%
\pgfpathrectangle{\pgfqpoint{0.532060in}{0.456901in}}{\pgfqpoint{5.803720in}{2.816433in}}%
\pgfusepath{clip}%
\pgfsetbuttcap%
\pgfsetroundjoin%
\pgfsetlinewidth{0.501875pt}%
\definecolor{currentstroke}{rgb}{0.501961,0.501961,0.501961}%
\pgfsetstrokecolor{currentstroke}%
\pgfsetstrokeopacity{0.400000}%
\pgfsetdash{{1.850000pt}{0.800000pt}}{0.000000pt}%
\pgfpathmoveto{\pgfqpoint{3.515134in}{0.456901in}}%
\pgfpathlineto{\pgfqpoint{3.515134in}{3.273333in}}%
\pgfusepath{stroke}%
\end{pgfscope}%
\begin{pgfscope}%
\pgfsetbuttcap%
\pgfsetroundjoin%
\definecolor{currentfill}{rgb}{1.000000,1.000000,1.000000}%
\pgfsetfillcolor{currentfill}%
\pgfsetlinewidth{0.803000pt}%
\definecolor{currentstroke}{rgb}{0.000000,0.000000,0.000000}%
\pgfsetstrokecolor{currentstroke}%
\pgfsetdash{}{0pt}%
\pgfsys@defobject{currentmarker}{\pgfqpoint{0.000000in}{0.000000in}}{\pgfqpoint{0.000000in}{0.048611in}}{%
\pgfpathmoveto{\pgfqpoint{0.000000in}{0.000000in}}%
\pgfpathlineto{\pgfqpoint{0.000000in}{0.048611in}}%
\pgfusepath{stroke,fill}%
}%
\begin{pgfscope}%
\pgfsys@transformshift{3.515134in}{0.456901in}%
\pgfsys@useobject{currentmarker}{}%
\end{pgfscope}%
\end{pgfscope}%
\begin{pgfscope}%
\definecolor{textcolor}{rgb}{0.000000,0.000000,0.000000}%
\pgfsetstrokecolor{textcolor}%
\pgfsetfillcolor{textcolor}%
\pgftext[x=3.515134in,y=0.408290in,,top]{\color{textcolor}{\rmfamily\fontsize{12.000000}{14.400000}\selectfont\catcode`\^=\active\def^{\ifmmode\sp\else\^{}\fi}\catcode`\%=\active\def%{\%}$\mathdefault{50}$}}%
\end{pgfscope}%
\begin{pgfscope}%
\pgfpathrectangle{\pgfqpoint{0.532060in}{0.456901in}}{\pgfqpoint{5.803720in}{2.816433in}}%
\pgfusepath{clip}%
\pgfsetbuttcap%
\pgfsetroundjoin%
\pgfsetlinewidth{0.501875pt}%
\definecolor{currentstroke}{rgb}{0.501961,0.501961,0.501961}%
\pgfsetstrokecolor{currentstroke}%
\pgfsetstrokeopacity{0.400000}%
\pgfsetdash{{1.850000pt}{0.800000pt}}{0.000000pt}%
\pgfpathmoveto{\pgfqpoint{4.215255in}{0.456901in}}%
\pgfpathlineto{\pgfqpoint{4.215255in}{3.273333in}}%
\pgfusepath{stroke}%
\end{pgfscope}%
\begin{pgfscope}%
\pgfsetbuttcap%
\pgfsetroundjoin%
\definecolor{currentfill}{rgb}{1.000000,1.000000,1.000000}%
\pgfsetfillcolor{currentfill}%
\pgfsetlinewidth{0.803000pt}%
\definecolor{currentstroke}{rgb}{0.000000,0.000000,0.000000}%
\pgfsetstrokecolor{currentstroke}%
\pgfsetdash{}{0pt}%
\pgfsys@defobject{currentmarker}{\pgfqpoint{0.000000in}{0.000000in}}{\pgfqpoint{0.000000in}{0.048611in}}{%
\pgfpathmoveto{\pgfqpoint{0.000000in}{0.000000in}}%
\pgfpathlineto{\pgfqpoint{0.000000in}{0.048611in}}%
\pgfusepath{stroke,fill}%
}%
\begin{pgfscope}%
\pgfsys@transformshift{4.215255in}{0.456901in}%
\pgfsys@useobject{currentmarker}{}%
\end{pgfscope}%
\end{pgfscope}%
\begin{pgfscope}%
\definecolor{textcolor}{rgb}{0.000000,0.000000,0.000000}%
\pgfsetstrokecolor{textcolor}%
\pgfsetfillcolor{textcolor}%
\pgftext[x=4.215255in,y=0.408290in,,top]{\color{textcolor}{\rmfamily\fontsize{12.000000}{14.400000}\selectfont\catcode`\^=\active\def^{\ifmmode\sp\else\^{}\fi}\catcode`\%=\active\def%{\%}$\mathdefault{75}$}}%
\end{pgfscope}%
\begin{pgfscope}%
\pgfpathrectangle{\pgfqpoint{0.532060in}{0.456901in}}{\pgfqpoint{5.803720in}{2.816433in}}%
\pgfusepath{clip}%
\pgfsetbuttcap%
\pgfsetroundjoin%
\pgfsetlinewidth{0.501875pt}%
\definecolor{currentstroke}{rgb}{0.501961,0.501961,0.501961}%
\pgfsetstrokecolor{currentstroke}%
\pgfsetstrokeopacity{0.400000}%
\pgfsetdash{{1.850000pt}{0.800000pt}}{0.000000pt}%
\pgfpathmoveto{\pgfqpoint{4.915375in}{0.456901in}}%
\pgfpathlineto{\pgfqpoint{4.915375in}{3.273333in}}%
\pgfusepath{stroke}%
\end{pgfscope}%
\begin{pgfscope}%
\pgfsetbuttcap%
\pgfsetroundjoin%
\definecolor{currentfill}{rgb}{1.000000,1.000000,1.000000}%
\pgfsetfillcolor{currentfill}%
\pgfsetlinewidth{0.803000pt}%
\definecolor{currentstroke}{rgb}{0.000000,0.000000,0.000000}%
\pgfsetstrokecolor{currentstroke}%
\pgfsetdash{}{0pt}%
\pgfsys@defobject{currentmarker}{\pgfqpoint{0.000000in}{0.000000in}}{\pgfqpoint{0.000000in}{0.048611in}}{%
\pgfpathmoveto{\pgfqpoint{0.000000in}{0.000000in}}%
\pgfpathlineto{\pgfqpoint{0.000000in}{0.048611in}}%
\pgfusepath{stroke,fill}%
}%
\begin{pgfscope}%
\pgfsys@transformshift{4.915375in}{0.456901in}%
\pgfsys@useobject{currentmarker}{}%
\end{pgfscope}%
\end{pgfscope}%
\begin{pgfscope}%
\definecolor{textcolor}{rgb}{0.000000,0.000000,0.000000}%
\pgfsetstrokecolor{textcolor}%
\pgfsetfillcolor{textcolor}%
\pgftext[x=4.915375in,y=0.408290in,,top]{\color{textcolor}{\rmfamily\fontsize{12.000000}{14.400000}\selectfont\catcode`\^=\active\def^{\ifmmode\sp\else\^{}\fi}\catcode`\%=\active\def%{\%}$\mathdefault{100}$}}%
\end{pgfscope}%
\begin{pgfscope}%
\pgfpathrectangle{\pgfqpoint{0.532060in}{0.456901in}}{\pgfqpoint{5.803720in}{2.816433in}}%
\pgfusepath{clip}%
\pgfsetbuttcap%
\pgfsetroundjoin%
\pgfsetlinewidth{0.501875pt}%
\definecolor{currentstroke}{rgb}{0.501961,0.501961,0.501961}%
\pgfsetstrokecolor{currentstroke}%
\pgfsetstrokeopacity{0.400000}%
\pgfsetdash{{1.850000pt}{0.800000pt}}{0.000000pt}%
\pgfpathmoveto{\pgfqpoint{5.615496in}{0.456901in}}%
\pgfpathlineto{\pgfqpoint{5.615496in}{3.273333in}}%
\pgfusepath{stroke}%
\end{pgfscope}%
\begin{pgfscope}%
\pgfsetbuttcap%
\pgfsetroundjoin%
\definecolor{currentfill}{rgb}{1.000000,1.000000,1.000000}%
\pgfsetfillcolor{currentfill}%
\pgfsetlinewidth{0.803000pt}%
\definecolor{currentstroke}{rgb}{0.000000,0.000000,0.000000}%
\pgfsetstrokecolor{currentstroke}%
\pgfsetdash{}{0pt}%
\pgfsys@defobject{currentmarker}{\pgfqpoint{0.000000in}{0.000000in}}{\pgfqpoint{0.000000in}{0.048611in}}{%
\pgfpathmoveto{\pgfqpoint{0.000000in}{0.000000in}}%
\pgfpathlineto{\pgfqpoint{0.000000in}{0.048611in}}%
\pgfusepath{stroke,fill}%
}%
\begin{pgfscope}%
\pgfsys@transformshift{5.615496in}{0.456901in}%
\pgfsys@useobject{currentmarker}{}%
\end{pgfscope}%
\end{pgfscope}%
\begin{pgfscope}%
\definecolor{textcolor}{rgb}{0.000000,0.000000,0.000000}%
\pgfsetstrokecolor{textcolor}%
\pgfsetfillcolor{textcolor}%
\pgftext[x=5.615496in,y=0.408290in,,top]{\color{textcolor}{\rmfamily\fontsize{12.000000}{14.400000}\selectfont\catcode`\^=\active\def^{\ifmmode\sp\else\^{}\fi}\catcode`\%=\active\def%{\%}$\mathdefault{125}$}}%
\end{pgfscope}%
\begin{pgfscope}%
\pgfpathrectangle{\pgfqpoint{0.532060in}{0.456901in}}{\pgfqpoint{5.803720in}{2.816433in}}%
\pgfusepath{clip}%
\pgfsetbuttcap%
\pgfsetroundjoin%
\pgfsetlinewidth{0.501875pt}%
\definecolor{currentstroke}{rgb}{0.501961,0.501961,0.501961}%
\pgfsetstrokecolor{currentstroke}%
\pgfsetstrokeopacity{0.400000}%
\pgfsetdash{{1.850000pt}{0.800000pt}}{0.000000pt}%
\pgfpathmoveto{\pgfqpoint{6.315617in}{0.456901in}}%
\pgfpathlineto{\pgfqpoint{6.315617in}{3.273333in}}%
\pgfusepath{stroke}%
\end{pgfscope}%
\begin{pgfscope}%
\pgfsetbuttcap%
\pgfsetroundjoin%
\definecolor{currentfill}{rgb}{1.000000,1.000000,1.000000}%
\pgfsetfillcolor{currentfill}%
\pgfsetlinewidth{0.803000pt}%
\definecolor{currentstroke}{rgb}{0.000000,0.000000,0.000000}%
\pgfsetstrokecolor{currentstroke}%
\pgfsetdash{}{0pt}%
\pgfsys@defobject{currentmarker}{\pgfqpoint{0.000000in}{0.000000in}}{\pgfqpoint{0.000000in}{0.048611in}}{%
\pgfpathmoveto{\pgfqpoint{0.000000in}{0.000000in}}%
\pgfpathlineto{\pgfqpoint{0.000000in}{0.048611in}}%
\pgfusepath{stroke,fill}%
}%
\begin{pgfscope}%
\pgfsys@transformshift{6.315617in}{0.456901in}%
\pgfsys@useobject{currentmarker}{}%
\end{pgfscope}%
\end{pgfscope}%
\begin{pgfscope}%
\definecolor{textcolor}{rgb}{0.000000,0.000000,0.000000}%
\pgfsetstrokecolor{textcolor}%
\pgfsetfillcolor{textcolor}%
\pgftext[x=6.315617in,y=0.408290in,,top]{\color{textcolor}{\rmfamily\fontsize{12.000000}{14.400000}\selectfont\catcode`\^=\active\def^{\ifmmode\sp\else\^{}\fi}\catcode`\%=\active\def%{\%}$\mathdefault{150}$}}%
\end{pgfscope}%
\begin{pgfscope}%
\pgfpathrectangle{\pgfqpoint{0.532060in}{0.456901in}}{\pgfqpoint{5.803720in}{2.816433in}}%
\pgfusepath{clip}%
\pgfsetbuttcap%
\pgfsetroundjoin%
\pgfsetlinewidth{0.501875pt}%
\definecolor{currentstroke}{rgb}{0.501961,0.501961,0.501961}%
\pgfsetstrokecolor{currentstroke}%
\pgfsetstrokeopacity{0.400000}%
\pgfsetdash{{1.850000pt}{0.800000pt}}{0.000000pt}%
\pgfpathmoveto{\pgfqpoint{0.574627in}{0.456901in}}%
\pgfpathlineto{\pgfqpoint{0.574627in}{3.273333in}}%
\pgfusepath{stroke}%
\end{pgfscope}%
\begin{pgfscope}%
\pgfsetbuttcap%
\pgfsetroundjoin%
\definecolor{currentfill}{rgb}{1.000000,1.000000,1.000000}%
\pgfsetfillcolor{currentfill}%
\pgfsetlinewidth{0.602250pt}%
\definecolor{currentstroke}{rgb}{0.000000,0.000000,0.000000}%
\pgfsetstrokecolor{currentstroke}%
\pgfsetdash{}{0pt}%
\pgfsys@defobject{currentmarker}{\pgfqpoint{0.000000in}{0.000000in}}{\pgfqpoint{0.000000in}{0.027778in}}{%
\pgfpathmoveto{\pgfqpoint{0.000000in}{0.000000in}}%
\pgfpathlineto{\pgfqpoint{0.000000in}{0.027778in}}%
\pgfusepath{stroke,fill}%
}%
\begin{pgfscope}%
\pgfsys@transformshift{0.574627in}{0.456901in}%
\pgfsys@useobject{currentmarker}{}%
\end{pgfscope}%
\end{pgfscope}%
\begin{pgfscope}%
\pgfpathrectangle{\pgfqpoint{0.532060in}{0.456901in}}{\pgfqpoint{5.803720in}{2.816433in}}%
\pgfusepath{clip}%
\pgfsetbuttcap%
\pgfsetroundjoin%
\pgfsetlinewidth{0.501875pt}%
\definecolor{currentstroke}{rgb}{0.501961,0.501961,0.501961}%
\pgfsetstrokecolor{currentstroke}%
\pgfsetstrokeopacity{0.400000}%
\pgfsetdash{{1.850000pt}{0.800000pt}}{0.000000pt}%
\pgfpathmoveto{\pgfqpoint{0.854675in}{0.456901in}}%
\pgfpathlineto{\pgfqpoint{0.854675in}{3.273333in}}%
\pgfusepath{stroke}%
\end{pgfscope}%
\begin{pgfscope}%
\pgfsetbuttcap%
\pgfsetroundjoin%
\definecolor{currentfill}{rgb}{1.000000,1.000000,1.000000}%
\pgfsetfillcolor{currentfill}%
\pgfsetlinewidth{0.602250pt}%
\definecolor{currentstroke}{rgb}{0.000000,0.000000,0.000000}%
\pgfsetstrokecolor{currentstroke}%
\pgfsetdash{}{0pt}%
\pgfsys@defobject{currentmarker}{\pgfqpoint{0.000000in}{0.000000in}}{\pgfqpoint{0.000000in}{0.027778in}}{%
\pgfpathmoveto{\pgfqpoint{0.000000in}{0.000000in}}%
\pgfpathlineto{\pgfqpoint{0.000000in}{0.027778in}}%
\pgfusepath{stroke,fill}%
}%
\begin{pgfscope}%
\pgfsys@transformshift{0.854675in}{0.456901in}%
\pgfsys@useobject{currentmarker}{}%
\end{pgfscope}%
\end{pgfscope}%
\begin{pgfscope}%
\pgfpathrectangle{\pgfqpoint{0.532060in}{0.456901in}}{\pgfqpoint{5.803720in}{2.816433in}}%
\pgfusepath{clip}%
\pgfsetbuttcap%
\pgfsetroundjoin%
\pgfsetlinewidth{0.501875pt}%
\definecolor{currentstroke}{rgb}{0.501961,0.501961,0.501961}%
\pgfsetstrokecolor{currentstroke}%
\pgfsetstrokeopacity{0.400000}%
\pgfsetdash{{1.850000pt}{0.800000pt}}{0.000000pt}%
\pgfpathmoveto{\pgfqpoint{0.994699in}{0.456901in}}%
\pgfpathlineto{\pgfqpoint{0.994699in}{3.273333in}}%
\pgfusepath{stroke}%
\end{pgfscope}%
\begin{pgfscope}%
\pgfsetbuttcap%
\pgfsetroundjoin%
\definecolor{currentfill}{rgb}{1.000000,1.000000,1.000000}%
\pgfsetfillcolor{currentfill}%
\pgfsetlinewidth{0.602250pt}%
\definecolor{currentstroke}{rgb}{0.000000,0.000000,0.000000}%
\pgfsetstrokecolor{currentstroke}%
\pgfsetdash{}{0pt}%
\pgfsys@defobject{currentmarker}{\pgfqpoint{0.000000in}{0.000000in}}{\pgfqpoint{0.000000in}{0.027778in}}{%
\pgfpathmoveto{\pgfqpoint{0.000000in}{0.000000in}}%
\pgfpathlineto{\pgfqpoint{0.000000in}{0.027778in}}%
\pgfusepath{stroke,fill}%
}%
\begin{pgfscope}%
\pgfsys@transformshift{0.994699in}{0.456901in}%
\pgfsys@useobject{currentmarker}{}%
\end{pgfscope}%
\end{pgfscope}%
\begin{pgfscope}%
\pgfpathrectangle{\pgfqpoint{0.532060in}{0.456901in}}{\pgfqpoint{5.803720in}{2.816433in}}%
\pgfusepath{clip}%
\pgfsetbuttcap%
\pgfsetroundjoin%
\pgfsetlinewidth{0.501875pt}%
\definecolor{currentstroke}{rgb}{0.501961,0.501961,0.501961}%
\pgfsetstrokecolor{currentstroke}%
\pgfsetstrokeopacity{0.400000}%
\pgfsetdash{{1.850000pt}{0.800000pt}}{0.000000pt}%
\pgfpathmoveto{\pgfqpoint{1.134724in}{0.456901in}}%
\pgfpathlineto{\pgfqpoint{1.134724in}{3.273333in}}%
\pgfusepath{stroke}%
\end{pgfscope}%
\begin{pgfscope}%
\pgfsetbuttcap%
\pgfsetroundjoin%
\definecolor{currentfill}{rgb}{1.000000,1.000000,1.000000}%
\pgfsetfillcolor{currentfill}%
\pgfsetlinewidth{0.602250pt}%
\definecolor{currentstroke}{rgb}{0.000000,0.000000,0.000000}%
\pgfsetstrokecolor{currentstroke}%
\pgfsetdash{}{0pt}%
\pgfsys@defobject{currentmarker}{\pgfqpoint{0.000000in}{0.000000in}}{\pgfqpoint{0.000000in}{0.027778in}}{%
\pgfpathmoveto{\pgfqpoint{0.000000in}{0.000000in}}%
\pgfpathlineto{\pgfqpoint{0.000000in}{0.027778in}}%
\pgfusepath{stroke,fill}%
}%
\begin{pgfscope}%
\pgfsys@transformshift{1.134724in}{0.456901in}%
\pgfsys@useobject{currentmarker}{}%
\end{pgfscope}%
\end{pgfscope}%
\begin{pgfscope}%
\pgfpathrectangle{\pgfqpoint{0.532060in}{0.456901in}}{\pgfqpoint{5.803720in}{2.816433in}}%
\pgfusepath{clip}%
\pgfsetbuttcap%
\pgfsetroundjoin%
\pgfsetlinewidth{0.501875pt}%
\definecolor{currentstroke}{rgb}{0.501961,0.501961,0.501961}%
\pgfsetstrokecolor{currentstroke}%
\pgfsetstrokeopacity{0.400000}%
\pgfsetdash{{1.850000pt}{0.800000pt}}{0.000000pt}%
\pgfpathmoveto{\pgfqpoint{1.274748in}{0.456901in}}%
\pgfpathlineto{\pgfqpoint{1.274748in}{3.273333in}}%
\pgfusepath{stroke}%
\end{pgfscope}%
\begin{pgfscope}%
\pgfsetbuttcap%
\pgfsetroundjoin%
\definecolor{currentfill}{rgb}{1.000000,1.000000,1.000000}%
\pgfsetfillcolor{currentfill}%
\pgfsetlinewidth{0.602250pt}%
\definecolor{currentstroke}{rgb}{0.000000,0.000000,0.000000}%
\pgfsetstrokecolor{currentstroke}%
\pgfsetdash{}{0pt}%
\pgfsys@defobject{currentmarker}{\pgfqpoint{0.000000in}{0.000000in}}{\pgfqpoint{0.000000in}{0.027778in}}{%
\pgfpathmoveto{\pgfqpoint{0.000000in}{0.000000in}}%
\pgfpathlineto{\pgfqpoint{0.000000in}{0.027778in}}%
\pgfusepath{stroke,fill}%
}%
\begin{pgfscope}%
\pgfsys@transformshift{1.274748in}{0.456901in}%
\pgfsys@useobject{currentmarker}{}%
\end{pgfscope}%
\end{pgfscope}%
\begin{pgfscope}%
\pgfpathrectangle{\pgfqpoint{0.532060in}{0.456901in}}{\pgfqpoint{5.803720in}{2.816433in}}%
\pgfusepath{clip}%
\pgfsetbuttcap%
\pgfsetroundjoin%
\pgfsetlinewidth{0.501875pt}%
\definecolor{currentstroke}{rgb}{0.501961,0.501961,0.501961}%
\pgfsetstrokecolor{currentstroke}%
\pgfsetstrokeopacity{0.400000}%
\pgfsetdash{{1.850000pt}{0.800000pt}}{0.000000pt}%
\pgfpathmoveto{\pgfqpoint{1.554796in}{0.456901in}}%
\pgfpathlineto{\pgfqpoint{1.554796in}{3.273333in}}%
\pgfusepath{stroke}%
\end{pgfscope}%
\begin{pgfscope}%
\pgfsetbuttcap%
\pgfsetroundjoin%
\definecolor{currentfill}{rgb}{1.000000,1.000000,1.000000}%
\pgfsetfillcolor{currentfill}%
\pgfsetlinewidth{0.602250pt}%
\definecolor{currentstroke}{rgb}{0.000000,0.000000,0.000000}%
\pgfsetstrokecolor{currentstroke}%
\pgfsetdash{}{0pt}%
\pgfsys@defobject{currentmarker}{\pgfqpoint{0.000000in}{0.000000in}}{\pgfqpoint{0.000000in}{0.027778in}}{%
\pgfpathmoveto{\pgfqpoint{0.000000in}{0.000000in}}%
\pgfpathlineto{\pgfqpoint{0.000000in}{0.027778in}}%
\pgfusepath{stroke,fill}%
}%
\begin{pgfscope}%
\pgfsys@transformshift{1.554796in}{0.456901in}%
\pgfsys@useobject{currentmarker}{}%
\end{pgfscope}%
\end{pgfscope}%
\begin{pgfscope}%
\pgfpathrectangle{\pgfqpoint{0.532060in}{0.456901in}}{\pgfqpoint{5.803720in}{2.816433in}}%
\pgfusepath{clip}%
\pgfsetbuttcap%
\pgfsetroundjoin%
\pgfsetlinewidth{0.501875pt}%
\definecolor{currentstroke}{rgb}{0.501961,0.501961,0.501961}%
\pgfsetstrokecolor{currentstroke}%
\pgfsetstrokeopacity{0.400000}%
\pgfsetdash{{1.850000pt}{0.800000pt}}{0.000000pt}%
\pgfpathmoveto{\pgfqpoint{1.694820in}{0.456901in}}%
\pgfpathlineto{\pgfqpoint{1.694820in}{3.273333in}}%
\pgfusepath{stroke}%
\end{pgfscope}%
\begin{pgfscope}%
\pgfsetbuttcap%
\pgfsetroundjoin%
\definecolor{currentfill}{rgb}{1.000000,1.000000,1.000000}%
\pgfsetfillcolor{currentfill}%
\pgfsetlinewidth{0.602250pt}%
\definecolor{currentstroke}{rgb}{0.000000,0.000000,0.000000}%
\pgfsetstrokecolor{currentstroke}%
\pgfsetdash{}{0pt}%
\pgfsys@defobject{currentmarker}{\pgfqpoint{0.000000in}{0.000000in}}{\pgfqpoint{0.000000in}{0.027778in}}{%
\pgfpathmoveto{\pgfqpoint{0.000000in}{0.000000in}}%
\pgfpathlineto{\pgfqpoint{0.000000in}{0.027778in}}%
\pgfusepath{stroke,fill}%
}%
\begin{pgfscope}%
\pgfsys@transformshift{1.694820in}{0.456901in}%
\pgfsys@useobject{currentmarker}{}%
\end{pgfscope}%
\end{pgfscope}%
\begin{pgfscope}%
\pgfpathrectangle{\pgfqpoint{0.532060in}{0.456901in}}{\pgfqpoint{5.803720in}{2.816433in}}%
\pgfusepath{clip}%
\pgfsetbuttcap%
\pgfsetroundjoin%
\pgfsetlinewidth{0.501875pt}%
\definecolor{currentstroke}{rgb}{0.501961,0.501961,0.501961}%
\pgfsetstrokecolor{currentstroke}%
\pgfsetstrokeopacity{0.400000}%
\pgfsetdash{{1.850000pt}{0.800000pt}}{0.000000pt}%
\pgfpathmoveto{\pgfqpoint{1.834844in}{0.456901in}}%
\pgfpathlineto{\pgfqpoint{1.834844in}{3.273333in}}%
\pgfusepath{stroke}%
\end{pgfscope}%
\begin{pgfscope}%
\pgfsetbuttcap%
\pgfsetroundjoin%
\definecolor{currentfill}{rgb}{1.000000,1.000000,1.000000}%
\pgfsetfillcolor{currentfill}%
\pgfsetlinewidth{0.602250pt}%
\definecolor{currentstroke}{rgb}{0.000000,0.000000,0.000000}%
\pgfsetstrokecolor{currentstroke}%
\pgfsetdash{}{0pt}%
\pgfsys@defobject{currentmarker}{\pgfqpoint{0.000000in}{0.000000in}}{\pgfqpoint{0.000000in}{0.027778in}}{%
\pgfpathmoveto{\pgfqpoint{0.000000in}{0.000000in}}%
\pgfpathlineto{\pgfqpoint{0.000000in}{0.027778in}}%
\pgfusepath{stroke,fill}%
}%
\begin{pgfscope}%
\pgfsys@transformshift{1.834844in}{0.456901in}%
\pgfsys@useobject{currentmarker}{}%
\end{pgfscope}%
\end{pgfscope}%
\begin{pgfscope}%
\pgfpathrectangle{\pgfqpoint{0.532060in}{0.456901in}}{\pgfqpoint{5.803720in}{2.816433in}}%
\pgfusepath{clip}%
\pgfsetbuttcap%
\pgfsetroundjoin%
\pgfsetlinewidth{0.501875pt}%
\definecolor{currentstroke}{rgb}{0.501961,0.501961,0.501961}%
\pgfsetstrokecolor{currentstroke}%
\pgfsetstrokeopacity{0.400000}%
\pgfsetdash{{1.850000pt}{0.800000pt}}{0.000000pt}%
\pgfpathmoveto{\pgfqpoint{1.974868in}{0.456901in}}%
\pgfpathlineto{\pgfqpoint{1.974868in}{3.273333in}}%
\pgfusepath{stroke}%
\end{pgfscope}%
\begin{pgfscope}%
\pgfsetbuttcap%
\pgfsetroundjoin%
\definecolor{currentfill}{rgb}{1.000000,1.000000,1.000000}%
\pgfsetfillcolor{currentfill}%
\pgfsetlinewidth{0.602250pt}%
\definecolor{currentstroke}{rgb}{0.000000,0.000000,0.000000}%
\pgfsetstrokecolor{currentstroke}%
\pgfsetdash{}{0pt}%
\pgfsys@defobject{currentmarker}{\pgfqpoint{0.000000in}{0.000000in}}{\pgfqpoint{0.000000in}{0.027778in}}{%
\pgfpathmoveto{\pgfqpoint{0.000000in}{0.000000in}}%
\pgfpathlineto{\pgfqpoint{0.000000in}{0.027778in}}%
\pgfusepath{stroke,fill}%
}%
\begin{pgfscope}%
\pgfsys@transformshift{1.974868in}{0.456901in}%
\pgfsys@useobject{currentmarker}{}%
\end{pgfscope}%
\end{pgfscope}%
\begin{pgfscope}%
\pgfpathrectangle{\pgfqpoint{0.532060in}{0.456901in}}{\pgfqpoint{5.803720in}{2.816433in}}%
\pgfusepath{clip}%
\pgfsetbuttcap%
\pgfsetroundjoin%
\pgfsetlinewidth{0.501875pt}%
\definecolor{currentstroke}{rgb}{0.501961,0.501961,0.501961}%
\pgfsetstrokecolor{currentstroke}%
\pgfsetstrokeopacity{0.400000}%
\pgfsetdash{{1.850000pt}{0.800000pt}}{0.000000pt}%
\pgfpathmoveto{\pgfqpoint{2.254917in}{0.456901in}}%
\pgfpathlineto{\pgfqpoint{2.254917in}{3.273333in}}%
\pgfusepath{stroke}%
\end{pgfscope}%
\begin{pgfscope}%
\pgfsetbuttcap%
\pgfsetroundjoin%
\definecolor{currentfill}{rgb}{1.000000,1.000000,1.000000}%
\pgfsetfillcolor{currentfill}%
\pgfsetlinewidth{0.602250pt}%
\definecolor{currentstroke}{rgb}{0.000000,0.000000,0.000000}%
\pgfsetstrokecolor{currentstroke}%
\pgfsetdash{}{0pt}%
\pgfsys@defobject{currentmarker}{\pgfqpoint{0.000000in}{0.000000in}}{\pgfqpoint{0.000000in}{0.027778in}}{%
\pgfpathmoveto{\pgfqpoint{0.000000in}{0.000000in}}%
\pgfpathlineto{\pgfqpoint{0.000000in}{0.027778in}}%
\pgfusepath{stroke,fill}%
}%
\begin{pgfscope}%
\pgfsys@transformshift{2.254917in}{0.456901in}%
\pgfsys@useobject{currentmarker}{}%
\end{pgfscope}%
\end{pgfscope}%
\begin{pgfscope}%
\pgfpathrectangle{\pgfqpoint{0.532060in}{0.456901in}}{\pgfqpoint{5.803720in}{2.816433in}}%
\pgfusepath{clip}%
\pgfsetbuttcap%
\pgfsetroundjoin%
\pgfsetlinewidth{0.501875pt}%
\definecolor{currentstroke}{rgb}{0.501961,0.501961,0.501961}%
\pgfsetstrokecolor{currentstroke}%
\pgfsetstrokeopacity{0.400000}%
\pgfsetdash{{1.850000pt}{0.800000pt}}{0.000000pt}%
\pgfpathmoveto{\pgfqpoint{2.394941in}{0.456901in}}%
\pgfpathlineto{\pgfqpoint{2.394941in}{3.273333in}}%
\pgfusepath{stroke}%
\end{pgfscope}%
\begin{pgfscope}%
\pgfsetbuttcap%
\pgfsetroundjoin%
\definecolor{currentfill}{rgb}{1.000000,1.000000,1.000000}%
\pgfsetfillcolor{currentfill}%
\pgfsetlinewidth{0.602250pt}%
\definecolor{currentstroke}{rgb}{0.000000,0.000000,0.000000}%
\pgfsetstrokecolor{currentstroke}%
\pgfsetdash{}{0pt}%
\pgfsys@defobject{currentmarker}{\pgfqpoint{0.000000in}{0.000000in}}{\pgfqpoint{0.000000in}{0.027778in}}{%
\pgfpathmoveto{\pgfqpoint{0.000000in}{0.000000in}}%
\pgfpathlineto{\pgfqpoint{0.000000in}{0.027778in}}%
\pgfusepath{stroke,fill}%
}%
\begin{pgfscope}%
\pgfsys@transformshift{2.394941in}{0.456901in}%
\pgfsys@useobject{currentmarker}{}%
\end{pgfscope}%
\end{pgfscope}%
\begin{pgfscope}%
\pgfpathrectangle{\pgfqpoint{0.532060in}{0.456901in}}{\pgfqpoint{5.803720in}{2.816433in}}%
\pgfusepath{clip}%
\pgfsetbuttcap%
\pgfsetroundjoin%
\pgfsetlinewidth{0.501875pt}%
\definecolor{currentstroke}{rgb}{0.501961,0.501961,0.501961}%
\pgfsetstrokecolor{currentstroke}%
\pgfsetstrokeopacity{0.400000}%
\pgfsetdash{{1.850000pt}{0.800000pt}}{0.000000pt}%
\pgfpathmoveto{\pgfqpoint{2.534965in}{0.456901in}}%
\pgfpathlineto{\pgfqpoint{2.534965in}{3.273333in}}%
\pgfusepath{stroke}%
\end{pgfscope}%
\begin{pgfscope}%
\pgfsetbuttcap%
\pgfsetroundjoin%
\definecolor{currentfill}{rgb}{1.000000,1.000000,1.000000}%
\pgfsetfillcolor{currentfill}%
\pgfsetlinewidth{0.602250pt}%
\definecolor{currentstroke}{rgb}{0.000000,0.000000,0.000000}%
\pgfsetstrokecolor{currentstroke}%
\pgfsetdash{}{0pt}%
\pgfsys@defobject{currentmarker}{\pgfqpoint{0.000000in}{0.000000in}}{\pgfqpoint{0.000000in}{0.027778in}}{%
\pgfpathmoveto{\pgfqpoint{0.000000in}{0.000000in}}%
\pgfpathlineto{\pgfqpoint{0.000000in}{0.027778in}}%
\pgfusepath{stroke,fill}%
}%
\begin{pgfscope}%
\pgfsys@transformshift{2.534965in}{0.456901in}%
\pgfsys@useobject{currentmarker}{}%
\end{pgfscope}%
\end{pgfscope}%
\begin{pgfscope}%
\pgfpathrectangle{\pgfqpoint{0.532060in}{0.456901in}}{\pgfqpoint{5.803720in}{2.816433in}}%
\pgfusepath{clip}%
\pgfsetbuttcap%
\pgfsetroundjoin%
\pgfsetlinewidth{0.501875pt}%
\definecolor{currentstroke}{rgb}{0.501961,0.501961,0.501961}%
\pgfsetstrokecolor{currentstroke}%
\pgfsetstrokeopacity{0.400000}%
\pgfsetdash{{1.850000pt}{0.800000pt}}{0.000000pt}%
\pgfpathmoveto{\pgfqpoint{2.674989in}{0.456901in}}%
\pgfpathlineto{\pgfqpoint{2.674989in}{3.273333in}}%
\pgfusepath{stroke}%
\end{pgfscope}%
\begin{pgfscope}%
\pgfsetbuttcap%
\pgfsetroundjoin%
\definecolor{currentfill}{rgb}{1.000000,1.000000,1.000000}%
\pgfsetfillcolor{currentfill}%
\pgfsetlinewidth{0.602250pt}%
\definecolor{currentstroke}{rgb}{0.000000,0.000000,0.000000}%
\pgfsetstrokecolor{currentstroke}%
\pgfsetdash{}{0pt}%
\pgfsys@defobject{currentmarker}{\pgfqpoint{0.000000in}{0.000000in}}{\pgfqpoint{0.000000in}{0.027778in}}{%
\pgfpathmoveto{\pgfqpoint{0.000000in}{0.000000in}}%
\pgfpathlineto{\pgfqpoint{0.000000in}{0.027778in}}%
\pgfusepath{stroke,fill}%
}%
\begin{pgfscope}%
\pgfsys@transformshift{2.674989in}{0.456901in}%
\pgfsys@useobject{currentmarker}{}%
\end{pgfscope}%
\end{pgfscope}%
\begin{pgfscope}%
\pgfpathrectangle{\pgfqpoint{0.532060in}{0.456901in}}{\pgfqpoint{5.803720in}{2.816433in}}%
\pgfusepath{clip}%
\pgfsetbuttcap%
\pgfsetroundjoin%
\pgfsetlinewidth{0.501875pt}%
\definecolor{currentstroke}{rgb}{0.501961,0.501961,0.501961}%
\pgfsetstrokecolor{currentstroke}%
\pgfsetstrokeopacity{0.400000}%
\pgfsetdash{{1.850000pt}{0.800000pt}}{0.000000pt}%
\pgfpathmoveto{\pgfqpoint{2.955037in}{0.456901in}}%
\pgfpathlineto{\pgfqpoint{2.955037in}{3.273333in}}%
\pgfusepath{stroke}%
\end{pgfscope}%
\begin{pgfscope}%
\pgfsetbuttcap%
\pgfsetroundjoin%
\definecolor{currentfill}{rgb}{1.000000,1.000000,1.000000}%
\pgfsetfillcolor{currentfill}%
\pgfsetlinewidth{0.602250pt}%
\definecolor{currentstroke}{rgb}{0.000000,0.000000,0.000000}%
\pgfsetstrokecolor{currentstroke}%
\pgfsetdash{}{0pt}%
\pgfsys@defobject{currentmarker}{\pgfqpoint{0.000000in}{0.000000in}}{\pgfqpoint{0.000000in}{0.027778in}}{%
\pgfpathmoveto{\pgfqpoint{0.000000in}{0.000000in}}%
\pgfpathlineto{\pgfqpoint{0.000000in}{0.027778in}}%
\pgfusepath{stroke,fill}%
}%
\begin{pgfscope}%
\pgfsys@transformshift{2.955037in}{0.456901in}%
\pgfsys@useobject{currentmarker}{}%
\end{pgfscope}%
\end{pgfscope}%
\begin{pgfscope}%
\pgfpathrectangle{\pgfqpoint{0.532060in}{0.456901in}}{\pgfqpoint{5.803720in}{2.816433in}}%
\pgfusepath{clip}%
\pgfsetbuttcap%
\pgfsetroundjoin%
\pgfsetlinewidth{0.501875pt}%
\definecolor{currentstroke}{rgb}{0.501961,0.501961,0.501961}%
\pgfsetstrokecolor{currentstroke}%
\pgfsetstrokeopacity{0.400000}%
\pgfsetdash{{1.850000pt}{0.800000pt}}{0.000000pt}%
\pgfpathmoveto{\pgfqpoint{3.095061in}{0.456901in}}%
\pgfpathlineto{\pgfqpoint{3.095061in}{3.273333in}}%
\pgfusepath{stroke}%
\end{pgfscope}%
\begin{pgfscope}%
\pgfsetbuttcap%
\pgfsetroundjoin%
\definecolor{currentfill}{rgb}{1.000000,1.000000,1.000000}%
\pgfsetfillcolor{currentfill}%
\pgfsetlinewidth{0.602250pt}%
\definecolor{currentstroke}{rgb}{0.000000,0.000000,0.000000}%
\pgfsetstrokecolor{currentstroke}%
\pgfsetdash{}{0pt}%
\pgfsys@defobject{currentmarker}{\pgfqpoint{0.000000in}{0.000000in}}{\pgfqpoint{0.000000in}{0.027778in}}{%
\pgfpathmoveto{\pgfqpoint{0.000000in}{0.000000in}}%
\pgfpathlineto{\pgfqpoint{0.000000in}{0.027778in}}%
\pgfusepath{stroke,fill}%
}%
\begin{pgfscope}%
\pgfsys@transformshift{3.095061in}{0.456901in}%
\pgfsys@useobject{currentmarker}{}%
\end{pgfscope}%
\end{pgfscope}%
\begin{pgfscope}%
\pgfpathrectangle{\pgfqpoint{0.532060in}{0.456901in}}{\pgfqpoint{5.803720in}{2.816433in}}%
\pgfusepath{clip}%
\pgfsetbuttcap%
\pgfsetroundjoin%
\pgfsetlinewidth{0.501875pt}%
\definecolor{currentstroke}{rgb}{0.501961,0.501961,0.501961}%
\pgfsetstrokecolor{currentstroke}%
\pgfsetstrokeopacity{0.400000}%
\pgfsetdash{{1.850000pt}{0.800000pt}}{0.000000pt}%
\pgfpathmoveto{\pgfqpoint{3.235086in}{0.456901in}}%
\pgfpathlineto{\pgfqpoint{3.235086in}{3.273333in}}%
\pgfusepath{stroke}%
\end{pgfscope}%
\begin{pgfscope}%
\pgfsetbuttcap%
\pgfsetroundjoin%
\definecolor{currentfill}{rgb}{1.000000,1.000000,1.000000}%
\pgfsetfillcolor{currentfill}%
\pgfsetlinewidth{0.602250pt}%
\definecolor{currentstroke}{rgb}{0.000000,0.000000,0.000000}%
\pgfsetstrokecolor{currentstroke}%
\pgfsetdash{}{0pt}%
\pgfsys@defobject{currentmarker}{\pgfqpoint{0.000000in}{0.000000in}}{\pgfqpoint{0.000000in}{0.027778in}}{%
\pgfpathmoveto{\pgfqpoint{0.000000in}{0.000000in}}%
\pgfpathlineto{\pgfqpoint{0.000000in}{0.027778in}}%
\pgfusepath{stroke,fill}%
}%
\begin{pgfscope}%
\pgfsys@transformshift{3.235086in}{0.456901in}%
\pgfsys@useobject{currentmarker}{}%
\end{pgfscope}%
\end{pgfscope}%
\begin{pgfscope}%
\pgfpathrectangle{\pgfqpoint{0.532060in}{0.456901in}}{\pgfqpoint{5.803720in}{2.816433in}}%
\pgfusepath{clip}%
\pgfsetbuttcap%
\pgfsetroundjoin%
\pgfsetlinewidth{0.501875pt}%
\definecolor{currentstroke}{rgb}{0.501961,0.501961,0.501961}%
\pgfsetstrokecolor{currentstroke}%
\pgfsetstrokeopacity{0.400000}%
\pgfsetdash{{1.850000pt}{0.800000pt}}{0.000000pt}%
\pgfpathmoveto{\pgfqpoint{3.375110in}{0.456901in}}%
\pgfpathlineto{\pgfqpoint{3.375110in}{3.273333in}}%
\pgfusepath{stroke}%
\end{pgfscope}%
\begin{pgfscope}%
\pgfsetbuttcap%
\pgfsetroundjoin%
\definecolor{currentfill}{rgb}{1.000000,1.000000,1.000000}%
\pgfsetfillcolor{currentfill}%
\pgfsetlinewidth{0.602250pt}%
\definecolor{currentstroke}{rgb}{0.000000,0.000000,0.000000}%
\pgfsetstrokecolor{currentstroke}%
\pgfsetdash{}{0pt}%
\pgfsys@defobject{currentmarker}{\pgfqpoint{0.000000in}{0.000000in}}{\pgfqpoint{0.000000in}{0.027778in}}{%
\pgfpathmoveto{\pgfqpoint{0.000000in}{0.000000in}}%
\pgfpathlineto{\pgfqpoint{0.000000in}{0.027778in}}%
\pgfusepath{stroke,fill}%
}%
\begin{pgfscope}%
\pgfsys@transformshift{3.375110in}{0.456901in}%
\pgfsys@useobject{currentmarker}{}%
\end{pgfscope}%
\end{pgfscope}%
\begin{pgfscope}%
\pgfpathrectangle{\pgfqpoint{0.532060in}{0.456901in}}{\pgfqpoint{5.803720in}{2.816433in}}%
\pgfusepath{clip}%
\pgfsetbuttcap%
\pgfsetroundjoin%
\pgfsetlinewidth{0.501875pt}%
\definecolor{currentstroke}{rgb}{0.501961,0.501961,0.501961}%
\pgfsetstrokecolor{currentstroke}%
\pgfsetstrokeopacity{0.400000}%
\pgfsetdash{{1.850000pt}{0.800000pt}}{0.000000pt}%
\pgfpathmoveto{\pgfqpoint{3.655158in}{0.456901in}}%
\pgfpathlineto{\pgfqpoint{3.655158in}{3.273333in}}%
\pgfusepath{stroke}%
\end{pgfscope}%
\begin{pgfscope}%
\pgfsetbuttcap%
\pgfsetroundjoin%
\definecolor{currentfill}{rgb}{1.000000,1.000000,1.000000}%
\pgfsetfillcolor{currentfill}%
\pgfsetlinewidth{0.602250pt}%
\definecolor{currentstroke}{rgb}{0.000000,0.000000,0.000000}%
\pgfsetstrokecolor{currentstroke}%
\pgfsetdash{}{0pt}%
\pgfsys@defobject{currentmarker}{\pgfqpoint{0.000000in}{0.000000in}}{\pgfqpoint{0.000000in}{0.027778in}}{%
\pgfpathmoveto{\pgfqpoint{0.000000in}{0.000000in}}%
\pgfpathlineto{\pgfqpoint{0.000000in}{0.027778in}}%
\pgfusepath{stroke,fill}%
}%
\begin{pgfscope}%
\pgfsys@transformshift{3.655158in}{0.456901in}%
\pgfsys@useobject{currentmarker}{}%
\end{pgfscope}%
\end{pgfscope}%
\begin{pgfscope}%
\pgfpathrectangle{\pgfqpoint{0.532060in}{0.456901in}}{\pgfqpoint{5.803720in}{2.816433in}}%
\pgfusepath{clip}%
\pgfsetbuttcap%
\pgfsetroundjoin%
\pgfsetlinewidth{0.501875pt}%
\definecolor{currentstroke}{rgb}{0.501961,0.501961,0.501961}%
\pgfsetstrokecolor{currentstroke}%
\pgfsetstrokeopacity{0.400000}%
\pgfsetdash{{1.850000pt}{0.800000pt}}{0.000000pt}%
\pgfpathmoveto{\pgfqpoint{3.795182in}{0.456901in}}%
\pgfpathlineto{\pgfqpoint{3.795182in}{3.273333in}}%
\pgfusepath{stroke}%
\end{pgfscope}%
\begin{pgfscope}%
\pgfsetbuttcap%
\pgfsetroundjoin%
\definecolor{currentfill}{rgb}{1.000000,1.000000,1.000000}%
\pgfsetfillcolor{currentfill}%
\pgfsetlinewidth{0.602250pt}%
\definecolor{currentstroke}{rgb}{0.000000,0.000000,0.000000}%
\pgfsetstrokecolor{currentstroke}%
\pgfsetdash{}{0pt}%
\pgfsys@defobject{currentmarker}{\pgfqpoint{0.000000in}{0.000000in}}{\pgfqpoint{0.000000in}{0.027778in}}{%
\pgfpathmoveto{\pgfqpoint{0.000000in}{0.000000in}}%
\pgfpathlineto{\pgfqpoint{0.000000in}{0.027778in}}%
\pgfusepath{stroke,fill}%
}%
\begin{pgfscope}%
\pgfsys@transformshift{3.795182in}{0.456901in}%
\pgfsys@useobject{currentmarker}{}%
\end{pgfscope}%
\end{pgfscope}%
\begin{pgfscope}%
\pgfpathrectangle{\pgfqpoint{0.532060in}{0.456901in}}{\pgfqpoint{5.803720in}{2.816433in}}%
\pgfusepath{clip}%
\pgfsetbuttcap%
\pgfsetroundjoin%
\pgfsetlinewidth{0.501875pt}%
\definecolor{currentstroke}{rgb}{0.501961,0.501961,0.501961}%
\pgfsetstrokecolor{currentstroke}%
\pgfsetstrokeopacity{0.400000}%
\pgfsetdash{{1.850000pt}{0.800000pt}}{0.000000pt}%
\pgfpathmoveto{\pgfqpoint{3.935206in}{0.456901in}}%
\pgfpathlineto{\pgfqpoint{3.935206in}{3.273333in}}%
\pgfusepath{stroke}%
\end{pgfscope}%
\begin{pgfscope}%
\pgfsetbuttcap%
\pgfsetroundjoin%
\definecolor{currentfill}{rgb}{1.000000,1.000000,1.000000}%
\pgfsetfillcolor{currentfill}%
\pgfsetlinewidth{0.602250pt}%
\definecolor{currentstroke}{rgb}{0.000000,0.000000,0.000000}%
\pgfsetstrokecolor{currentstroke}%
\pgfsetdash{}{0pt}%
\pgfsys@defobject{currentmarker}{\pgfqpoint{0.000000in}{0.000000in}}{\pgfqpoint{0.000000in}{0.027778in}}{%
\pgfpathmoveto{\pgfqpoint{0.000000in}{0.000000in}}%
\pgfpathlineto{\pgfqpoint{0.000000in}{0.027778in}}%
\pgfusepath{stroke,fill}%
}%
\begin{pgfscope}%
\pgfsys@transformshift{3.935206in}{0.456901in}%
\pgfsys@useobject{currentmarker}{}%
\end{pgfscope}%
\end{pgfscope}%
\begin{pgfscope}%
\pgfpathrectangle{\pgfqpoint{0.532060in}{0.456901in}}{\pgfqpoint{5.803720in}{2.816433in}}%
\pgfusepath{clip}%
\pgfsetbuttcap%
\pgfsetroundjoin%
\pgfsetlinewidth{0.501875pt}%
\definecolor{currentstroke}{rgb}{0.501961,0.501961,0.501961}%
\pgfsetstrokecolor{currentstroke}%
\pgfsetstrokeopacity{0.400000}%
\pgfsetdash{{1.850000pt}{0.800000pt}}{0.000000pt}%
\pgfpathmoveto{\pgfqpoint{4.075230in}{0.456901in}}%
\pgfpathlineto{\pgfqpoint{4.075230in}{3.273333in}}%
\pgfusepath{stroke}%
\end{pgfscope}%
\begin{pgfscope}%
\pgfsetbuttcap%
\pgfsetroundjoin%
\definecolor{currentfill}{rgb}{1.000000,1.000000,1.000000}%
\pgfsetfillcolor{currentfill}%
\pgfsetlinewidth{0.602250pt}%
\definecolor{currentstroke}{rgb}{0.000000,0.000000,0.000000}%
\pgfsetstrokecolor{currentstroke}%
\pgfsetdash{}{0pt}%
\pgfsys@defobject{currentmarker}{\pgfqpoint{0.000000in}{0.000000in}}{\pgfqpoint{0.000000in}{0.027778in}}{%
\pgfpathmoveto{\pgfqpoint{0.000000in}{0.000000in}}%
\pgfpathlineto{\pgfqpoint{0.000000in}{0.027778in}}%
\pgfusepath{stroke,fill}%
}%
\begin{pgfscope}%
\pgfsys@transformshift{4.075230in}{0.456901in}%
\pgfsys@useobject{currentmarker}{}%
\end{pgfscope}%
\end{pgfscope}%
\begin{pgfscope}%
\pgfpathrectangle{\pgfqpoint{0.532060in}{0.456901in}}{\pgfqpoint{5.803720in}{2.816433in}}%
\pgfusepath{clip}%
\pgfsetbuttcap%
\pgfsetroundjoin%
\pgfsetlinewidth{0.501875pt}%
\definecolor{currentstroke}{rgb}{0.501961,0.501961,0.501961}%
\pgfsetstrokecolor{currentstroke}%
\pgfsetstrokeopacity{0.400000}%
\pgfsetdash{{1.850000pt}{0.800000pt}}{0.000000pt}%
\pgfpathmoveto{\pgfqpoint{4.355279in}{0.456901in}}%
\pgfpathlineto{\pgfqpoint{4.355279in}{3.273333in}}%
\pgfusepath{stroke}%
\end{pgfscope}%
\begin{pgfscope}%
\pgfsetbuttcap%
\pgfsetroundjoin%
\definecolor{currentfill}{rgb}{1.000000,1.000000,1.000000}%
\pgfsetfillcolor{currentfill}%
\pgfsetlinewidth{0.602250pt}%
\definecolor{currentstroke}{rgb}{0.000000,0.000000,0.000000}%
\pgfsetstrokecolor{currentstroke}%
\pgfsetdash{}{0pt}%
\pgfsys@defobject{currentmarker}{\pgfqpoint{0.000000in}{0.000000in}}{\pgfqpoint{0.000000in}{0.027778in}}{%
\pgfpathmoveto{\pgfqpoint{0.000000in}{0.000000in}}%
\pgfpathlineto{\pgfqpoint{0.000000in}{0.027778in}}%
\pgfusepath{stroke,fill}%
}%
\begin{pgfscope}%
\pgfsys@transformshift{4.355279in}{0.456901in}%
\pgfsys@useobject{currentmarker}{}%
\end{pgfscope}%
\end{pgfscope}%
\begin{pgfscope}%
\pgfpathrectangle{\pgfqpoint{0.532060in}{0.456901in}}{\pgfqpoint{5.803720in}{2.816433in}}%
\pgfusepath{clip}%
\pgfsetbuttcap%
\pgfsetroundjoin%
\pgfsetlinewidth{0.501875pt}%
\definecolor{currentstroke}{rgb}{0.501961,0.501961,0.501961}%
\pgfsetstrokecolor{currentstroke}%
\pgfsetstrokeopacity{0.400000}%
\pgfsetdash{{1.850000pt}{0.800000pt}}{0.000000pt}%
\pgfpathmoveto{\pgfqpoint{4.495303in}{0.456901in}}%
\pgfpathlineto{\pgfqpoint{4.495303in}{3.273333in}}%
\pgfusepath{stroke}%
\end{pgfscope}%
\begin{pgfscope}%
\pgfsetbuttcap%
\pgfsetroundjoin%
\definecolor{currentfill}{rgb}{1.000000,1.000000,1.000000}%
\pgfsetfillcolor{currentfill}%
\pgfsetlinewidth{0.602250pt}%
\definecolor{currentstroke}{rgb}{0.000000,0.000000,0.000000}%
\pgfsetstrokecolor{currentstroke}%
\pgfsetdash{}{0pt}%
\pgfsys@defobject{currentmarker}{\pgfqpoint{0.000000in}{0.000000in}}{\pgfqpoint{0.000000in}{0.027778in}}{%
\pgfpathmoveto{\pgfqpoint{0.000000in}{0.000000in}}%
\pgfpathlineto{\pgfqpoint{0.000000in}{0.027778in}}%
\pgfusepath{stroke,fill}%
}%
\begin{pgfscope}%
\pgfsys@transformshift{4.495303in}{0.456901in}%
\pgfsys@useobject{currentmarker}{}%
\end{pgfscope}%
\end{pgfscope}%
\begin{pgfscope}%
\pgfpathrectangle{\pgfqpoint{0.532060in}{0.456901in}}{\pgfqpoint{5.803720in}{2.816433in}}%
\pgfusepath{clip}%
\pgfsetbuttcap%
\pgfsetroundjoin%
\pgfsetlinewidth{0.501875pt}%
\definecolor{currentstroke}{rgb}{0.501961,0.501961,0.501961}%
\pgfsetstrokecolor{currentstroke}%
\pgfsetstrokeopacity{0.400000}%
\pgfsetdash{{1.850000pt}{0.800000pt}}{0.000000pt}%
\pgfpathmoveto{\pgfqpoint{4.635327in}{0.456901in}}%
\pgfpathlineto{\pgfqpoint{4.635327in}{3.273333in}}%
\pgfusepath{stroke}%
\end{pgfscope}%
\begin{pgfscope}%
\pgfsetbuttcap%
\pgfsetroundjoin%
\definecolor{currentfill}{rgb}{1.000000,1.000000,1.000000}%
\pgfsetfillcolor{currentfill}%
\pgfsetlinewidth{0.602250pt}%
\definecolor{currentstroke}{rgb}{0.000000,0.000000,0.000000}%
\pgfsetstrokecolor{currentstroke}%
\pgfsetdash{}{0pt}%
\pgfsys@defobject{currentmarker}{\pgfqpoint{0.000000in}{0.000000in}}{\pgfqpoint{0.000000in}{0.027778in}}{%
\pgfpathmoveto{\pgfqpoint{0.000000in}{0.000000in}}%
\pgfpathlineto{\pgfqpoint{0.000000in}{0.027778in}}%
\pgfusepath{stroke,fill}%
}%
\begin{pgfscope}%
\pgfsys@transformshift{4.635327in}{0.456901in}%
\pgfsys@useobject{currentmarker}{}%
\end{pgfscope}%
\end{pgfscope}%
\begin{pgfscope}%
\pgfpathrectangle{\pgfqpoint{0.532060in}{0.456901in}}{\pgfqpoint{5.803720in}{2.816433in}}%
\pgfusepath{clip}%
\pgfsetbuttcap%
\pgfsetroundjoin%
\pgfsetlinewidth{0.501875pt}%
\definecolor{currentstroke}{rgb}{0.501961,0.501961,0.501961}%
\pgfsetstrokecolor{currentstroke}%
\pgfsetstrokeopacity{0.400000}%
\pgfsetdash{{1.850000pt}{0.800000pt}}{0.000000pt}%
\pgfpathmoveto{\pgfqpoint{4.775351in}{0.456901in}}%
\pgfpathlineto{\pgfqpoint{4.775351in}{3.273333in}}%
\pgfusepath{stroke}%
\end{pgfscope}%
\begin{pgfscope}%
\pgfsetbuttcap%
\pgfsetroundjoin%
\definecolor{currentfill}{rgb}{1.000000,1.000000,1.000000}%
\pgfsetfillcolor{currentfill}%
\pgfsetlinewidth{0.602250pt}%
\definecolor{currentstroke}{rgb}{0.000000,0.000000,0.000000}%
\pgfsetstrokecolor{currentstroke}%
\pgfsetdash{}{0pt}%
\pgfsys@defobject{currentmarker}{\pgfqpoint{0.000000in}{0.000000in}}{\pgfqpoint{0.000000in}{0.027778in}}{%
\pgfpathmoveto{\pgfqpoint{0.000000in}{0.000000in}}%
\pgfpathlineto{\pgfqpoint{0.000000in}{0.027778in}}%
\pgfusepath{stroke,fill}%
}%
\begin{pgfscope}%
\pgfsys@transformshift{4.775351in}{0.456901in}%
\pgfsys@useobject{currentmarker}{}%
\end{pgfscope}%
\end{pgfscope}%
\begin{pgfscope}%
\pgfpathrectangle{\pgfqpoint{0.532060in}{0.456901in}}{\pgfqpoint{5.803720in}{2.816433in}}%
\pgfusepath{clip}%
\pgfsetbuttcap%
\pgfsetroundjoin%
\pgfsetlinewidth{0.501875pt}%
\definecolor{currentstroke}{rgb}{0.501961,0.501961,0.501961}%
\pgfsetstrokecolor{currentstroke}%
\pgfsetstrokeopacity{0.400000}%
\pgfsetdash{{1.850000pt}{0.800000pt}}{0.000000pt}%
\pgfpathmoveto{\pgfqpoint{5.055399in}{0.456901in}}%
\pgfpathlineto{\pgfqpoint{5.055399in}{3.273333in}}%
\pgfusepath{stroke}%
\end{pgfscope}%
\begin{pgfscope}%
\pgfsetbuttcap%
\pgfsetroundjoin%
\definecolor{currentfill}{rgb}{1.000000,1.000000,1.000000}%
\pgfsetfillcolor{currentfill}%
\pgfsetlinewidth{0.602250pt}%
\definecolor{currentstroke}{rgb}{0.000000,0.000000,0.000000}%
\pgfsetstrokecolor{currentstroke}%
\pgfsetdash{}{0pt}%
\pgfsys@defobject{currentmarker}{\pgfqpoint{0.000000in}{0.000000in}}{\pgfqpoint{0.000000in}{0.027778in}}{%
\pgfpathmoveto{\pgfqpoint{0.000000in}{0.000000in}}%
\pgfpathlineto{\pgfqpoint{0.000000in}{0.027778in}}%
\pgfusepath{stroke,fill}%
}%
\begin{pgfscope}%
\pgfsys@transformshift{5.055399in}{0.456901in}%
\pgfsys@useobject{currentmarker}{}%
\end{pgfscope}%
\end{pgfscope}%
\begin{pgfscope}%
\pgfpathrectangle{\pgfqpoint{0.532060in}{0.456901in}}{\pgfqpoint{5.803720in}{2.816433in}}%
\pgfusepath{clip}%
\pgfsetbuttcap%
\pgfsetroundjoin%
\pgfsetlinewidth{0.501875pt}%
\definecolor{currentstroke}{rgb}{0.501961,0.501961,0.501961}%
\pgfsetstrokecolor{currentstroke}%
\pgfsetstrokeopacity{0.400000}%
\pgfsetdash{{1.850000pt}{0.800000pt}}{0.000000pt}%
\pgfpathmoveto{\pgfqpoint{5.195423in}{0.456901in}}%
\pgfpathlineto{\pgfqpoint{5.195423in}{3.273333in}}%
\pgfusepath{stroke}%
\end{pgfscope}%
\begin{pgfscope}%
\pgfsetbuttcap%
\pgfsetroundjoin%
\definecolor{currentfill}{rgb}{1.000000,1.000000,1.000000}%
\pgfsetfillcolor{currentfill}%
\pgfsetlinewidth{0.602250pt}%
\definecolor{currentstroke}{rgb}{0.000000,0.000000,0.000000}%
\pgfsetstrokecolor{currentstroke}%
\pgfsetdash{}{0pt}%
\pgfsys@defobject{currentmarker}{\pgfqpoint{0.000000in}{0.000000in}}{\pgfqpoint{0.000000in}{0.027778in}}{%
\pgfpathmoveto{\pgfqpoint{0.000000in}{0.000000in}}%
\pgfpathlineto{\pgfqpoint{0.000000in}{0.027778in}}%
\pgfusepath{stroke,fill}%
}%
\begin{pgfscope}%
\pgfsys@transformshift{5.195423in}{0.456901in}%
\pgfsys@useobject{currentmarker}{}%
\end{pgfscope}%
\end{pgfscope}%
\begin{pgfscope}%
\pgfpathrectangle{\pgfqpoint{0.532060in}{0.456901in}}{\pgfqpoint{5.803720in}{2.816433in}}%
\pgfusepath{clip}%
\pgfsetbuttcap%
\pgfsetroundjoin%
\pgfsetlinewidth{0.501875pt}%
\definecolor{currentstroke}{rgb}{0.501961,0.501961,0.501961}%
\pgfsetstrokecolor{currentstroke}%
\pgfsetstrokeopacity{0.400000}%
\pgfsetdash{{1.850000pt}{0.800000pt}}{0.000000pt}%
\pgfpathmoveto{\pgfqpoint{5.335448in}{0.456901in}}%
\pgfpathlineto{\pgfqpoint{5.335448in}{3.273333in}}%
\pgfusepath{stroke}%
\end{pgfscope}%
\begin{pgfscope}%
\pgfsetbuttcap%
\pgfsetroundjoin%
\definecolor{currentfill}{rgb}{1.000000,1.000000,1.000000}%
\pgfsetfillcolor{currentfill}%
\pgfsetlinewidth{0.602250pt}%
\definecolor{currentstroke}{rgb}{0.000000,0.000000,0.000000}%
\pgfsetstrokecolor{currentstroke}%
\pgfsetdash{}{0pt}%
\pgfsys@defobject{currentmarker}{\pgfqpoint{0.000000in}{0.000000in}}{\pgfqpoint{0.000000in}{0.027778in}}{%
\pgfpathmoveto{\pgfqpoint{0.000000in}{0.000000in}}%
\pgfpathlineto{\pgfqpoint{0.000000in}{0.027778in}}%
\pgfusepath{stroke,fill}%
}%
\begin{pgfscope}%
\pgfsys@transformshift{5.335448in}{0.456901in}%
\pgfsys@useobject{currentmarker}{}%
\end{pgfscope}%
\end{pgfscope}%
\begin{pgfscope}%
\pgfpathrectangle{\pgfqpoint{0.532060in}{0.456901in}}{\pgfqpoint{5.803720in}{2.816433in}}%
\pgfusepath{clip}%
\pgfsetbuttcap%
\pgfsetroundjoin%
\pgfsetlinewidth{0.501875pt}%
\definecolor{currentstroke}{rgb}{0.501961,0.501961,0.501961}%
\pgfsetstrokecolor{currentstroke}%
\pgfsetstrokeopacity{0.400000}%
\pgfsetdash{{1.850000pt}{0.800000pt}}{0.000000pt}%
\pgfpathmoveto{\pgfqpoint{5.475472in}{0.456901in}}%
\pgfpathlineto{\pgfqpoint{5.475472in}{3.273333in}}%
\pgfusepath{stroke}%
\end{pgfscope}%
\begin{pgfscope}%
\pgfsetbuttcap%
\pgfsetroundjoin%
\definecolor{currentfill}{rgb}{1.000000,1.000000,1.000000}%
\pgfsetfillcolor{currentfill}%
\pgfsetlinewidth{0.602250pt}%
\definecolor{currentstroke}{rgb}{0.000000,0.000000,0.000000}%
\pgfsetstrokecolor{currentstroke}%
\pgfsetdash{}{0pt}%
\pgfsys@defobject{currentmarker}{\pgfqpoint{0.000000in}{0.000000in}}{\pgfqpoint{0.000000in}{0.027778in}}{%
\pgfpathmoveto{\pgfqpoint{0.000000in}{0.000000in}}%
\pgfpathlineto{\pgfqpoint{0.000000in}{0.027778in}}%
\pgfusepath{stroke,fill}%
}%
\begin{pgfscope}%
\pgfsys@transformshift{5.475472in}{0.456901in}%
\pgfsys@useobject{currentmarker}{}%
\end{pgfscope}%
\end{pgfscope}%
\begin{pgfscope}%
\pgfpathrectangle{\pgfqpoint{0.532060in}{0.456901in}}{\pgfqpoint{5.803720in}{2.816433in}}%
\pgfusepath{clip}%
\pgfsetbuttcap%
\pgfsetroundjoin%
\pgfsetlinewidth{0.501875pt}%
\definecolor{currentstroke}{rgb}{0.501961,0.501961,0.501961}%
\pgfsetstrokecolor{currentstroke}%
\pgfsetstrokeopacity{0.400000}%
\pgfsetdash{{1.850000pt}{0.800000pt}}{0.000000pt}%
\pgfpathmoveto{\pgfqpoint{5.755520in}{0.456901in}}%
\pgfpathlineto{\pgfqpoint{5.755520in}{3.273333in}}%
\pgfusepath{stroke}%
\end{pgfscope}%
\begin{pgfscope}%
\pgfsetbuttcap%
\pgfsetroundjoin%
\definecolor{currentfill}{rgb}{1.000000,1.000000,1.000000}%
\pgfsetfillcolor{currentfill}%
\pgfsetlinewidth{0.602250pt}%
\definecolor{currentstroke}{rgb}{0.000000,0.000000,0.000000}%
\pgfsetstrokecolor{currentstroke}%
\pgfsetdash{}{0pt}%
\pgfsys@defobject{currentmarker}{\pgfqpoint{0.000000in}{0.000000in}}{\pgfqpoint{0.000000in}{0.027778in}}{%
\pgfpathmoveto{\pgfqpoint{0.000000in}{0.000000in}}%
\pgfpathlineto{\pgfqpoint{0.000000in}{0.027778in}}%
\pgfusepath{stroke,fill}%
}%
\begin{pgfscope}%
\pgfsys@transformshift{5.755520in}{0.456901in}%
\pgfsys@useobject{currentmarker}{}%
\end{pgfscope}%
\end{pgfscope}%
\begin{pgfscope}%
\pgfpathrectangle{\pgfqpoint{0.532060in}{0.456901in}}{\pgfqpoint{5.803720in}{2.816433in}}%
\pgfusepath{clip}%
\pgfsetbuttcap%
\pgfsetroundjoin%
\pgfsetlinewidth{0.501875pt}%
\definecolor{currentstroke}{rgb}{0.501961,0.501961,0.501961}%
\pgfsetstrokecolor{currentstroke}%
\pgfsetstrokeopacity{0.400000}%
\pgfsetdash{{1.850000pt}{0.800000pt}}{0.000000pt}%
\pgfpathmoveto{\pgfqpoint{5.895544in}{0.456901in}}%
\pgfpathlineto{\pgfqpoint{5.895544in}{3.273333in}}%
\pgfusepath{stroke}%
\end{pgfscope}%
\begin{pgfscope}%
\pgfsetbuttcap%
\pgfsetroundjoin%
\definecolor{currentfill}{rgb}{1.000000,1.000000,1.000000}%
\pgfsetfillcolor{currentfill}%
\pgfsetlinewidth{0.602250pt}%
\definecolor{currentstroke}{rgb}{0.000000,0.000000,0.000000}%
\pgfsetstrokecolor{currentstroke}%
\pgfsetdash{}{0pt}%
\pgfsys@defobject{currentmarker}{\pgfqpoint{0.000000in}{0.000000in}}{\pgfqpoint{0.000000in}{0.027778in}}{%
\pgfpathmoveto{\pgfqpoint{0.000000in}{0.000000in}}%
\pgfpathlineto{\pgfqpoint{0.000000in}{0.027778in}}%
\pgfusepath{stroke,fill}%
}%
\begin{pgfscope}%
\pgfsys@transformshift{5.895544in}{0.456901in}%
\pgfsys@useobject{currentmarker}{}%
\end{pgfscope}%
\end{pgfscope}%
\begin{pgfscope}%
\pgfpathrectangle{\pgfqpoint{0.532060in}{0.456901in}}{\pgfqpoint{5.803720in}{2.816433in}}%
\pgfusepath{clip}%
\pgfsetbuttcap%
\pgfsetroundjoin%
\pgfsetlinewidth{0.501875pt}%
\definecolor{currentstroke}{rgb}{0.501961,0.501961,0.501961}%
\pgfsetstrokecolor{currentstroke}%
\pgfsetstrokeopacity{0.400000}%
\pgfsetdash{{1.850000pt}{0.800000pt}}{0.000000pt}%
\pgfpathmoveto{\pgfqpoint{6.035568in}{0.456901in}}%
\pgfpathlineto{\pgfqpoint{6.035568in}{3.273333in}}%
\pgfusepath{stroke}%
\end{pgfscope}%
\begin{pgfscope}%
\pgfsetbuttcap%
\pgfsetroundjoin%
\definecolor{currentfill}{rgb}{1.000000,1.000000,1.000000}%
\pgfsetfillcolor{currentfill}%
\pgfsetlinewidth{0.602250pt}%
\definecolor{currentstroke}{rgb}{0.000000,0.000000,0.000000}%
\pgfsetstrokecolor{currentstroke}%
\pgfsetdash{}{0pt}%
\pgfsys@defobject{currentmarker}{\pgfqpoint{0.000000in}{0.000000in}}{\pgfqpoint{0.000000in}{0.027778in}}{%
\pgfpathmoveto{\pgfqpoint{0.000000in}{0.000000in}}%
\pgfpathlineto{\pgfqpoint{0.000000in}{0.027778in}}%
\pgfusepath{stroke,fill}%
}%
\begin{pgfscope}%
\pgfsys@transformshift{6.035568in}{0.456901in}%
\pgfsys@useobject{currentmarker}{}%
\end{pgfscope}%
\end{pgfscope}%
\begin{pgfscope}%
\pgfpathrectangle{\pgfqpoint{0.532060in}{0.456901in}}{\pgfqpoint{5.803720in}{2.816433in}}%
\pgfusepath{clip}%
\pgfsetbuttcap%
\pgfsetroundjoin%
\pgfsetlinewidth{0.501875pt}%
\definecolor{currentstroke}{rgb}{0.501961,0.501961,0.501961}%
\pgfsetstrokecolor{currentstroke}%
\pgfsetstrokeopacity{0.400000}%
\pgfsetdash{{1.850000pt}{0.800000pt}}{0.000000pt}%
\pgfpathmoveto{\pgfqpoint{6.175592in}{0.456901in}}%
\pgfpathlineto{\pgfqpoint{6.175592in}{3.273333in}}%
\pgfusepath{stroke}%
\end{pgfscope}%
\begin{pgfscope}%
\pgfsetbuttcap%
\pgfsetroundjoin%
\definecolor{currentfill}{rgb}{1.000000,1.000000,1.000000}%
\pgfsetfillcolor{currentfill}%
\pgfsetlinewidth{0.602250pt}%
\definecolor{currentstroke}{rgb}{0.000000,0.000000,0.000000}%
\pgfsetstrokecolor{currentstroke}%
\pgfsetdash{}{0pt}%
\pgfsys@defobject{currentmarker}{\pgfqpoint{0.000000in}{0.000000in}}{\pgfqpoint{0.000000in}{0.027778in}}{%
\pgfpathmoveto{\pgfqpoint{0.000000in}{0.000000in}}%
\pgfpathlineto{\pgfqpoint{0.000000in}{0.027778in}}%
\pgfusepath{stroke,fill}%
}%
\begin{pgfscope}%
\pgfsys@transformshift{6.175592in}{0.456901in}%
\pgfsys@useobject{currentmarker}{}%
\end{pgfscope}%
\end{pgfscope}%
\begin{pgfscope}%
\definecolor{textcolor}{rgb}{0.000000,0.000000,0.000000}%
\pgfsetstrokecolor{textcolor}%
\pgfsetfillcolor{textcolor}%
\pgftext[x=3.433920in,y=0.201367in,,top]{\color{textcolor}{\rmfamily\fontsize{12.000000}{14.400000}\selectfont\catcode`\^=\active\def^{\ifmmode\sp\else\^{}\fi}\catcode`\%=\active\def%{\%}$t$}}%
\end{pgfscope}%
\begin{pgfscope}%
\pgfpathrectangle{\pgfqpoint{0.532060in}{0.456901in}}{\pgfqpoint{5.803720in}{2.816433in}}%
\pgfusepath{clip}%
\pgfsetbuttcap%
\pgfsetroundjoin%
\pgfsetlinewidth{0.501875pt}%
\definecolor{currentstroke}{rgb}{0.501961,0.501961,0.501961}%
\pgfsetstrokecolor{currentstroke}%
\pgfsetstrokeopacity{0.400000}%
\pgfsetdash{{1.850000pt}{0.800000pt}}{0.000000pt}%
\pgfpathmoveto{\pgfqpoint{0.532060in}{0.553642in}}%
\pgfpathlineto{\pgfqpoint{6.335780in}{0.553642in}}%
\pgfusepath{stroke}%
\end{pgfscope}%
\begin{pgfscope}%
\pgfsetbuttcap%
\pgfsetroundjoin%
\definecolor{currentfill}{rgb}{1.000000,1.000000,1.000000}%
\pgfsetfillcolor{currentfill}%
\pgfsetlinewidth{0.803000pt}%
\definecolor{currentstroke}{rgb}{0.000000,0.000000,0.000000}%
\pgfsetstrokecolor{currentstroke}%
\pgfsetdash{}{0pt}%
\pgfsys@defobject{currentmarker}{\pgfqpoint{0.000000in}{0.000000in}}{\pgfqpoint{0.048611in}{0.000000in}}{%
\pgfpathmoveto{\pgfqpoint{0.000000in}{0.000000in}}%
\pgfpathlineto{\pgfqpoint{0.048611in}{0.000000in}}%
\pgfusepath{stroke,fill}%
}%
\begin{pgfscope}%
\pgfsys@transformshift{0.532060in}{0.553642in}%
\pgfsys@useobject{currentmarker}{}%
\end{pgfscope}%
\end{pgfscope}%
\begin{pgfscope}%
\definecolor{textcolor}{rgb}{0.000000,0.000000,0.000000}%
\pgfsetstrokecolor{textcolor}%
\pgfsetfillcolor{textcolor}%
\pgftext[x=0.272222in, y=0.495780in, left, base]{\color{textcolor}{\rmfamily\fontsize{12.000000}{14.400000}\selectfont\catcode`\^=\active\def^{\ifmmode\sp\else\^{}\fi}\catcode`\%=\active\def%{\%}$\mathdefault{\ensuremath{-}4}$}}%
\end{pgfscope}%
\begin{pgfscope}%
\pgfpathrectangle{\pgfqpoint{0.532060in}{0.456901in}}{\pgfqpoint{5.803720in}{2.816433in}}%
\pgfusepath{clip}%
\pgfsetbuttcap%
\pgfsetroundjoin%
\pgfsetlinewidth{0.501875pt}%
\definecolor{currentstroke}{rgb}{0.501961,0.501961,0.501961}%
\pgfsetstrokecolor{currentstroke}%
\pgfsetstrokeopacity{0.400000}%
\pgfsetdash{{1.850000pt}{0.800000pt}}{0.000000pt}%
\pgfpathmoveto{\pgfqpoint{0.532060in}{1.209389in}}%
\pgfpathlineto{\pgfqpoint{6.335780in}{1.209389in}}%
\pgfusepath{stroke}%
\end{pgfscope}%
\begin{pgfscope}%
\pgfsetbuttcap%
\pgfsetroundjoin%
\definecolor{currentfill}{rgb}{1.000000,1.000000,1.000000}%
\pgfsetfillcolor{currentfill}%
\pgfsetlinewidth{0.803000pt}%
\definecolor{currentstroke}{rgb}{0.000000,0.000000,0.000000}%
\pgfsetstrokecolor{currentstroke}%
\pgfsetdash{}{0pt}%
\pgfsys@defobject{currentmarker}{\pgfqpoint{0.000000in}{0.000000in}}{\pgfqpoint{0.048611in}{0.000000in}}{%
\pgfpathmoveto{\pgfqpoint{0.000000in}{0.000000in}}%
\pgfpathlineto{\pgfqpoint{0.048611in}{0.000000in}}%
\pgfusepath{stroke,fill}%
}%
\begin{pgfscope}%
\pgfsys@transformshift{0.532060in}{1.209389in}%
\pgfsys@useobject{currentmarker}{}%
\end{pgfscope}%
\end{pgfscope}%
\begin{pgfscope}%
\definecolor{textcolor}{rgb}{0.000000,0.000000,0.000000}%
\pgfsetstrokecolor{textcolor}%
\pgfsetfillcolor{textcolor}%
\pgftext[x=0.272222in, y=1.151527in, left, base]{\color{textcolor}{\rmfamily\fontsize{12.000000}{14.400000}\selectfont\catcode`\^=\active\def^{\ifmmode\sp\else\^{}\fi}\catcode`\%=\active\def%{\%}$\mathdefault{\ensuremath{-}2}$}}%
\end{pgfscope}%
\begin{pgfscope}%
\pgfpathrectangle{\pgfqpoint{0.532060in}{0.456901in}}{\pgfqpoint{5.803720in}{2.816433in}}%
\pgfusepath{clip}%
\pgfsetbuttcap%
\pgfsetroundjoin%
\pgfsetlinewidth{0.501875pt}%
\definecolor{currentstroke}{rgb}{0.501961,0.501961,0.501961}%
\pgfsetstrokecolor{currentstroke}%
\pgfsetstrokeopacity{0.400000}%
\pgfsetdash{{1.850000pt}{0.800000pt}}{0.000000pt}%
\pgfpathmoveto{\pgfqpoint{0.532060in}{1.865136in}}%
\pgfpathlineto{\pgfqpoint{6.335780in}{1.865136in}}%
\pgfusepath{stroke}%
\end{pgfscope}%
\begin{pgfscope}%
\pgfsetbuttcap%
\pgfsetroundjoin%
\definecolor{currentfill}{rgb}{1.000000,1.000000,1.000000}%
\pgfsetfillcolor{currentfill}%
\pgfsetlinewidth{0.803000pt}%
\definecolor{currentstroke}{rgb}{0.000000,0.000000,0.000000}%
\pgfsetstrokecolor{currentstroke}%
\pgfsetdash{}{0pt}%
\pgfsys@defobject{currentmarker}{\pgfqpoint{0.000000in}{0.000000in}}{\pgfqpoint{0.048611in}{0.000000in}}{%
\pgfpathmoveto{\pgfqpoint{0.000000in}{0.000000in}}%
\pgfpathlineto{\pgfqpoint{0.048611in}{0.000000in}}%
\pgfusepath{stroke,fill}%
}%
\begin{pgfscope}%
\pgfsys@transformshift{0.532060in}{1.865136in}%
\pgfsys@useobject{currentmarker}{}%
\end{pgfscope}%
\end{pgfscope}%
\begin{pgfscope}%
\definecolor{textcolor}{rgb}{0.000000,0.000000,0.000000}%
\pgfsetstrokecolor{textcolor}%
\pgfsetfillcolor{textcolor}%
\pgftext[x=0.401852in, y=1.807274in, left, base]{\color{textcolor}{\rmfamily\fontsize{12.000000}{14.400000}\selectfont\catcode`\^=\active\def^{\ifmmode\sp\else\^{}\fi}\catcode`\%=\active\def%{\%}$\mathdefault{0}$}}%
\end{pgfscope}%
\begin{pgfscope}%
\pgfpathrectangle{\pgfqpoint{0.532060in}{0.456901in}}{\pgfqpoint{5.803720in}{2.816433in}}%
\pgfusepath{clip}%
\pgfsetbuttcap%
\pgfsetroundjoin%
\pgfsetlinewidth{0.501875pt}%
\definecolor{currentstroke}{rgb}{0.501961,0.501961,0.501961}%
\pgfsetstrokecolor{currentstroke}%
\pgfsetstrokeopacity{0.400000}%
\pgfsetdash{{1.850000pt}{0.800000pt}}{0.000000pt}%
\pgfpathmoveto{\pgfqpoint{0.532060in}{2.520883in}}%
\pgfpathlineto{\pgfqpoint{6.335780in}{2.520883in}}%
\pgfusepath{stroke}%
\end{pgfscope}%
\begin{pgfscope}%
\pgfsetbuttcap%
\pgfsetroundjoin%
\definecolor{currentfill}{rgb}{1.000000,1.000000,1.000000}%
\pgfsetfillcolor{currentfill}%
\pgfsetlinewidth{0.803000pt}%
\definecolor{currentstroke}{rgb}{0.000000,0.000000,0.000000}%
\pgfsetstrokecolor{currentstroke}%
\pgfsetdash{}{0pt}%
\pgfsys@defobject{currentmarker}{\pgfqpoint{0.000000in}{0.000000in}}{\pgfqpoint{0.048611in}{0.000000in}}{%
\pgfpathmoveto{\pgfqpoint{0.000000in}{0.000000in}}%
\pgfpathlineto{\pgfqpoint{0.048611in}{0.000000in}}%
\pgfusepath{stroke,fill}%
}%
\begin{pgfscope}%
\pgfsys@transformshift{0.532060in}{2.520883in}%
\pgfsys@useobject{currentmarker}{}%
\end{pgfscope}%
\end{pgfscope}%
\begin{pgfscope}%
\definecolor{textcolor}{rgb}{0.000000,0.000000,0.000000}%
\pgfsetstrokecolor{textcolor}%
\pgfsetfillcolor{textcolor}%
\pgftext[x=0.401852in, y=2.463021in, left, base]{\color{textcolor}{\rmfamily\fontsize{12.000000}{14.400000}\selectfont\catcode`\^=\active\def^{\ifmmode\sp\else\^{}\fi}\catcode`\%=\active\def%{\%}$\mathdefault{2}$}}%
\end{pgfscope}%
\begin{pgfscope}%
\pgfpathrectangle{\pgfqpoint{0.532060in}{0.456901in}}{\pgfqpoint{5.803720in}{2.816433in}}%
\pgfusepath{clip}%
\pgfsetbuttcap%
\pgfsetroundjoin%
\pgfsetlinewidth{0.501875pt}%
\definecolor{currentstroke}{rgb}{0.501961,0.501961,0.501961}%
\pgfsetstrokecolor{currentstroke}%
\pgfsetstrokeopacity{0.400000}%
\pgfsetdash{{1.850000pt}{0.800000pt}}{0.000000pt}%
\pgfpathmoveto{\pgfqpoint{0.532060in}{3.176630in}}%
\pgfpathlineto{\pgfqpoint{6.335780in}{3.176630in}}%
\pgfusepath{stroke}%
\end{pgfscope}%
\begin{pgfscope}%
\pgfsetbuttcap%
\pgfsetroundjoin%
\definecolor{currentfill}{rgb}{1.000000,1.000000,1.000000}%
\pgfsetfillcolor{currentfill}%
\pgfsetlinewidth{0.803000pt}%
\definecolor{currentstroke}{rgb}{0.000000,0.000000,0.000000}%
\pgfsetstrokecolor{currentstroke}%
\pgfsetdash{}{0pt}%
\pgfsys@defobject{currentmarker}{\pgfqpoint{0.000000in}{0.000000in}}{\pgfqpoint{0.048611in}{0.000000in}}{%
\pgfpathmoveto{\pgfqpoint{0.000000in}{0.000000in}}%
\pgfpathlineto{\pgfqpoint{0.048611in}{0.000000in}}%
\pgfusepath{stroke,fill}%
}%
\begin{pgfscope}%
\pgfsys@transformshift{0.532060in}{3.176630in}%
\pgfsys@useobject{currentmarker}{}%
\end{pgfscope}%
\end{pgfscope}%
\begin{pgfscope}%
\definecolor{textcolor}{rgb}{0.000000,0.000000,0.000000}%
\pgfsetstrokecolor{textcolor}%
\pgfsetfillcolor{textcolor}%
\pgftext[x=0.401852in, y=3.118768in, left, base]{\color{textcolor}{\rmfamily\fontsize{12.000000}{14.400000}\selectfont\catcode`\^=\active\def^{\ifmmode\sp\else\^{}\fi}\catcode`\%=\active\def%{\%}$\mathdefault{4}$}}%
\end{pgfscope}%
\begin{pgfscope}%
\definecolor{textcolor}{rgb}{0.000000,0.000000,0.000000}%
\pgfsetstrokecolor{textcolor}%
\pgfsetfillcolor{textcolor}%
\pgftext[x=0.216667in,y=1.865117in,,bottom,rotate=90.000000]{\color{textcolor}{\rmfamily\fontsize{12.000000}{14.400000}\selectfont\catcode`\^=\active\def^{\ifmmode\sp\else\^{}\fi}\catcode`\%=\active\def%{\%}$i_1(t)$}}%
\end{pgfscope}%
\begin{pgfscope}%
\pgfpathrectangle{\pgfqpoint{0.532060in}{0.456901in}}{\pgfqpoint{5.803720in}{2.816433in}}%
\pgfusepath{clip}%
\pgfsetbuttcap%
\pgfsetroundjoin%
\pgfsetlinewidth{2.007500pt}%
\definecolor{currentstroke}{rgb}{0.000000,0.000000,0.000000}%
\pgfsetstrokecolor{currentstroke}%
\pgfsetdash{{7.400000pt}{3.200000pt}}{0.000000pt}%
\pgfpathmoveto{\pgfqpoint{0.795865in}{1.209389in}}%
\pgfpathlineto{\pgfqpoint{2.110931in}{1.209389in}}%
\pgfpathlineto{\pgfqpoint{2.116213in}{2.520883in}}%
\pgfpathlineto{\pgfqpoint{3.431279in}{2.520883in}}%
\pgfpathlineto{\pgfqpoint{3.436561in}{1.209389in}}%
\pgfpathlineto{\pgfqpoint{4.751627in}{1.209389in}}%
\pgfpathlineto{\pgfqpoint{4.756908in}{2.520883in}}%
\pgfpathlineto{\pgfqpoint{6.066693in}{2.520883in}}%
\pgfpathlineto{\pgfqpoint{6.071975in}{1.209389in}}%
\pgfpathlineto{\pgfqpoint{6.071975in}{1.209389in}}%
\pgfusepath{stroke}%
\end{pgfscope}%
\begin{pgfscope}%
\pgfpathrectangle{\pgfqpoint{0.532060in}{0.456901in}}{\pgfqpoint{5.803720in}{2.816433in}}%
\pgfusepath{clip}%
\pgfsetrectcap%
\pgfsetroundjoin%
\pgfsetlinewidth{2.007500pt}%
\definecolor{currentstroke}{rgb}{1.000000,0.000000,0.000000}%
\pgfsetstrokecolor{currentstroke}%
\pgfsetdash{}{0pt}%
\pgfpathmoveto{\pgfqpoint{0.795865in}{0.584920in}}%
\pgfpathlineto{\pgfqpoint{0.801147in}{0.585515in}}%
\pgfpathlineto{\pgfqpoint{0.806428in}{0.587296in}}%
\pgfpathlineto{\pgfqpoint{0.816991in}{0.594398in}}%
\pgfpathlineto{\pgfqpoint{0.827554in}{0.606150in}}%
\pgfpathlineto{\pgfqpoint{0.838116in}{0.622426in}}%
\pgfpathlineto{\pgfqpoint{0.848679in}{0.643053in}}%
\pgfpathlineto{\pgfqpoint{0.859242in}{0.667810in}}%
\pgfpathlineto{\pgfqpoint{0.875086in}{0.712109in}}%
\pgfpathlineto{\pgfqpoint{0.890930in}{0.764062in}}%
\pgfpathlineto{\pgfqpoint{0.912056in}{0.843137in}}%
\pgfpathlineto{\pgfqpoint{0.938463in}{0.953197in}}%
\pgfpathlineto{\pgfqpoint{1.017684in}{1.294528in}}%
\pgfpathlineto{\pgfqpoint{1.038809in}{1.373726in}}%
\pgfpathlineto{\pgfqpoint{1.059935in}{1.443462in}}%
\pgfpathlineto{\pgfqpoint{1.075779in}{1.488741in}}%
\pgfpathlineto{\pgfqpoint{1.091623in}{1.527653in}}%
\pgfpathlineto{\pgfqpoint{1.107467in}{1.560128in}}%
\pgfpathlineto{\pgfqpoint{1.123311in}{1.586313in}}%
\pgfpathlineto{\pgfqpoint{1.133874in}{1.600442in}}%
\pgfpathlineto{\pgfqpoint{1.144437in}{1.612079in}}%
\pgfpathlineto{\pgfqpoint{1.155000in}{1.621408in}}%
\pgfpathlineto{\pgfqpoint{1.165563in}{1.628640in}}%
\pgfpathlineto{\pgfqpoint{1.181407in}{1.636084in}}%
\pgfpathlineto{\pgfqpoint{1.197251in}{1.640234in}}%
\pgfpathlineto{\pgfqpoint{1.213095in}{1.642039in}}%
\pgfpathlineto{\pgfqpoint{1.244783in}{1.642520in}}%
\pgfpathlineto{\pgfqpoint{1.265909in}{1.643542in}}%
\pgfpathlineto{\pgfqpoint{1.281753in}{1.646105in}}%
\pgfpathlineto{\pgfqpoint{1.297597in}{1.650985in}}%
\pgfpathlineto{\pgfqpoint{1.313441in}{1.658730in}}%
\pgfpathlineto{\pgfqpoint{1.329286in}{1.669730in}}%
\pgfpathlineto{\pgfqpoint{1.345130in}{1.684201in}}%
\pgfpathlineto{\pgfqpoint{1.360974in}{1.702187in}}%
\pgfpathlineto{\pgfqpoint{1.376818in}{1.723559in}}%
\pgfpathlineto{\pgfqpoint{1.397944in}{1.756798in}}%
\pgfpathlineto{\pgfqpoint{1.419069in}{1.794443in}}%
\pgfpathlineto{\pgfqpoint{1.450758in}{1.855949in}}%
\pgfpathlineto{\pgfqpoint{1.498290in}{1.947985in}}%
\pgfpathlineto{\pgfqpoint{1.519416in}{1.984393in}}%
\pgfpathlineto{\pgfqpoint{1.540541in}{2.016014in}}%
\pgfpathlineto{\pgfqpoint{1.556385in}{2.036001in}}%
\pgfpathlineto{\pgfqpoint{1.572230in}{2.052528in}}%
\pgfpathlineto{\pgfqpoint{1.588074in}{2.065538in}}%
\pgfpathlineto{\pgfqpoint{1.603918in}{2.075145in}}%
\pgfpathlineto{\pgfqpoint{1.619762in}{2.081640in}}%
\pgfpathlineto{\pgfqpoint{1.635606in}{2.085479in}}%
\pgfpathlineto{\pgfqpoint{1.656732in}{2.087540in}}%
\pgfpathlineto{\pgfqpoint{1.714827in}{2.090281in}}%
\pgfpathlineto{\pgfqpoint{1.730671in}{2.094669in}}%
\pgfpathlineto{\pgfqpoint{1.741234in}{2.099412in}}%
\pgfpathlineto{\pgfqpoint{1.751797in}{2.105918in}}%
\pgfpathlineto{\pgfqpoint{1.762360in}{2.114434in}}%
\pgfpathlineto{\pgfqpoint{1.772922in}{2.125182in}}%
\pgfpathlineto{\pgfqpoint{1.783485in}{2.138357in}}%
\pgfpathlineto{\pgfqpoint{1.794048in}{2.154123in}}%
\pgfpathlineto{\pgfqpoint{1.809892in}{2.182897in}}%
\pgfpathlineto{\pgfqpoint{1.825736in}{2.218048in}}%
\pgfpathlineto{\pgfqpoint{1.841581in}{2.259634in}}%
\pgfpathlineto{\pgfqpoint{1.857425in}{2.307492in}}%
\pgfpathlineto{\pgfqpoint{1.878550in}{2.380358in}}%
\pgfpathlineto{\pgfqpoint{1.899676in}{2.462138in}}%
\pgfpathlineto{\pgfqpoint{1.926083in}{2.573548in}}%
\pgfpathlineto{\pgfqpoint{2.000022in}{2.892353in}}%
\pgfpathlineto{\pgfqpoint{2.021148in}{2.970802in}}%
\pgfpathlineto{\pgfqpoint{2.036992in}{3.022162in}}%
\pgfpathlineto{\pgfqpoint{2.052836in}{3.065773in}}%
\pgfpathlineto{\pgfqpoint{2.063399in}{3.090028in}}%
\pgfpathlineto{\pgfqpoint{2.073962in}{3.110121in}}%
\pgfpathlineto{\pgfqpoint{2.084525in}{3.125840in}}%
\pgfpathlineto{\pgfqpoint{2.095087in}{3.137017in}}%
\pgfpathlineto{\pgfqpoint{2.105650in}{3.143531in}}%
\pgfpathlineto{\pgfqpoint{2.110931in}{3.145016in}}%
\pgfpathlineto{\pgfqpoint{2.116213in}{3.145314in}}%
\pgfpathlineto{\pgfqpoint{2.121494in}{3.144422in}}%
\pgfpathlineto{\pgfqpoint{2.132057in}{3.139086in}}%
\pgfpathlineto{\pgfqpoint{2.142620in}{3.129066in}}%
\pgfpathlineto{\pgfqpoint{2.153183in}{3.114469in}}%
\pgfpathlineto{\pgfqpoint{2.163745in}{3.095450in}}%
\pgfpathlineto{\pgfqpoint{2.174308in}{3.072214in}}%
\pgfpathlineto{\pgfqpoint{2.190152in}{3.029998in}}%
\pgfpathlineto{\pgfqpoint{2.205997in}{2.979851in}}%
\pgfpathlineto{\pgfqpoint{2.221841in}{2.922941in}}%
\pgfpathlineto{\pgfqpoint{2.242966in}{2.838809in}}%
\pgfpathlineto{\pgfqpoint{2.274655in}{2.701781in}}%
\pgfpathlineto{\pgfqpoint{2.322187in}{2.494680in}}%
\pgfpathlineto{\pgfqpoint{2.348594in}{2.390132in}}%
\pgfpathlineto{\pgfqpoint{2.369720in}{2.316054in}}%
\pgfpathlineto{\pgfqpoint{2.385564in}{2.267183in}}%
\pgfpathlineto{\pgfqpoint{2.401408in}{2.224534in}}%
\pgfpathlineto{\pgfqpoint{2.417252in}{2.188309in}}%
\pgfpathlineto{\pgfqpoint{2.433096in}{2.158484in}}%
\pgfpathlineto{\pgfqpoint{2.448940in}{2.134826in}}%
\pgfpathlineto{\pgfqpoint{2.459503in}{2.122275in}}%
\pgfpathlineto{\pgfqpoint{2.470066in}{2.112104in}}%
\pgfpathlineto{\pgfqpoint{2.480629in}{2.104112in}}%
\pgfpathlineto{\pgfqpoint{2.491192in}{2.098071in}}%
\pgfpathlineto{\pgfqpoint{2.507036in}{2.092112in}}%
\pgfpathlineto{\pgfqpoint{2.522880in}{2.089056in}}%
\pgfpathlineto{\pgfqpoint{2.544006in}{2.087828in}}%
\pgfpathlineto{\pgfqpoint{2.586257in}{2.086588in}}%
\pgfpathlineto{\pgfqpoint{2.602101in}{2.083857in}}%
\pgfpathlineto{\pgfqpoint{2.617945in}{2.078757in}}%
\pgfpathlineto{\pgfqpoint{2.633789in}{2.070753in}}%
\pgfpathlineto{\pgfqpoint{2.649633in}{2.059470in}}%
\pgfpathlineto{\pgfqpoint{2.665478in}{2.044705in}}%
\pgfpathlineto{\pgfqpoint{2.681322in}{2.026430in}}%
\pgfpathlineto{\pgfqpoint{2.697166in}{2.004787in}}%
\pgfpathlineto{\pgfqpoint{2.718291in}{1.971236in}}%
\pgfpathlineto{\pgfqpoint{2.739417in}{1.933360in}}%
\pgfpathlineto{\pgfqpoint{2.776387in}{1.861197in}}%
\pgfpathlineto{\pgfqpoint{2.813356in}{1.789543in}}%
\pgfpathlineto{\pgfqpoint{2.834482in}{1.752376in}}%
\pgfpathlineto{\pgfqpoint{2.855608in}{1.719772in}}%
\pgfpathlineto{\pgfqpoint{2.871452in}{1.698949in}}%
\pgfpathlineto{\pgfqpoint{2.887296in}{1.681544in}}%
\pgfpathlineto{\pgfqpoint{2.903140in}{1.667659in}}%
\pgfpathlineto{\pgfqpoint{2.918984in}{1.657222in}}%
\pgfpathlineto{\pgfqpoint{2.934828in}{1.649985in}}%
\pgfpathlineto{\pgfqpoint{2.950673in}{1.645534in}}%
\pgfpathlineto{\pgfqpoint{2.966517in}{1.643294in}}%
\pgfpathlineto{\pgfqpoint{2.992924in}{1.642490in}}%
\pgfpathlineto{\pgfqpoint{3.024612in}{1.641369in}}%
\pgfpathlineto{\pgfqpoint{3.040456in}{1.638512in}}%
\pgfpathlineto{\pgfqpoint{3.056300in}{1.632831in}}%
\pgfpathlineto{\pgfqpoint{3.066863in}{1.627016in}}%
\pgfpathlineto{\pgfqpoint{3.077426in}{1.619281in}}%
\pgfpathlineto{\pgfqpoint{3.087989in}{1.609394in}}%
\pgfpathlineto{\pgfqpoint{3.098552in}{1.597150in}}%
\pgfpathlineto{\pgfqpoint{3.109114in}{1.582373in}}%
\pgfpathlineto{\pgfqpoint{3.124959in}{1.555159in}}%
\pgfpathlineto{\pgfqpoint{3.140803in}{1.521616in}}%
\pgfpathlineto{\pgfqpoint{3.156647in}{1.481632in}}%
\pgfpathlineto{\pgfqpoint{3.172491in}{1.435314in}}%
\pgfpathlineto{\pgfqpoint{3.188335in}{1.382994in}}%
\pgfpathlineto{\pgfqpoint{3.209461in}{1.304884in}}%
\pgfpathlineto{\pgfqpoint{3.235868in}{1.196599in}}%
\pgfpathlineto{\pgfqpoint{3.283400in}{0.987735in}}%
\pgfpathlineto{\pgfqpoint{3.315089in}{0.853671in}}%
\pgfpathlineto{\pgfqpoint{3.336214in}{0.773383in}}%
\pgfpathlineto{\pgfqpoint{3.352058in}{0.720267in}}%
\pgfpathlineto{\pgfqpoint{3.367902in}{0.674615in}}%
\pgfpathlineto{\pgfqpoint{3.383747in}{0.637499in}}%
\pgfpathlineto{\pgfqpoint{3.394309in}{0.617942in}}%
\pgfpathlineto{\pgfqpoint{3.404872in}{0.602782in}}%
\pgfpathlineto{\pgfqpoint{3.415435in}{0.592183in}}%
\pgfpathlineto{\pgfqpoint{3.425998in}{0.586257in}}%
\pgfpathlineto{\pgfqpoint{3.431279in}{0.585069in}}%
\pgfpathlineto{\pgfqpoint{3.436561in}{0.585069in}}%
\pgfpathlineto{\pgfqpoint{3.441842in}{0.586257in}}%
\pgfpathlineto{\pgfqpoint{3.452405in}{0.592183in}}%
\pgfpathlineto{\pgfqpoint{3.462968in}{0.602782in}}%
\pgfpathlineto{\pgfqpoint{3.473530in}{0.617942in}}%
\pgfpathlineto{\pgfqpoint{3.484093in}{0.637499in}}%
\pgfpathlineto{\pgfqpoint{3.494656in}{0.661247in}}%
\pgfpathlineto{\pgfqpoint{3.510500in}{0.704164in}}%
\pgfpathlineto{\pgfqpoint{3.526344in}{0.754921in}}%
\pgfpathlineto{\pgfqpoint{3.547470in}{0.832733in}}%
\pgfpathlineto{\pgfqpoint{3.573877in}{0.941799in}}%
\pgfpathlineto{\pgfqpoint{3.658379in}{1.304884in}}%
\pgfpathlineto{\pgfqpoint{3.679505in}{1.382994in}}%
\pgfpathlineto{\pgfqpoint{3.700630in}{1.451441in}}%
\pgfpathlineto{\pgfqpoint{3.716474in}{1.495673in}}%
\pgfpathlineto{\pgfqpoint{3.732318in}{1.533511in}}%
\pgfpathlineto{\pgfqpoint{3.748163in}{1.564922in}}%
\pgfpathlineto{\pgfqpoint{3.764007in}{1.590088in}}%
\pgfpathlineto{\pgfqpoint{3.774570in}{1.603579in}}%
\pgfpathlineto{\pgfqpoint{3.785132in}{1.614620in}}%
\pgfpathlineto{\pgfqpoint{3.795695in}{1.623403in}}%
\pgfpathlineto{\pgfqpoint{3.806258in}{1.630148in}}%
\pgfpathlineto{\pgfqpoint{3.822102in}{1.636979in}}%
\pgfpathlineto{\pgfqpoint{3.837946in}{1.640671in}}%
\pgfpathlineto{\pgfqpoint{3.859072in}{1.642368in}}%
\pgfpathlineto{\pgfqpoint{3.906604in}{1.643836in}}%
\pgfpathlineto{\pgfqpoint{3.922449in}{1.646740in}}%
\pgfpathlineto{\pgfqpoint{3.938293in}{1.652064in}}%
\pgfpathlineto{\pgfqpoint{3.954137in}{1.660329in}}%
\pgfpathlineto{\pgfqpoint{3.969981in}{1.671897in}}%
\pgfpathlineto{\pgfqpoint{3.985825in}{1.686956in}}%
\pgfpathlineto{\pgfqpoint{4.001669in}{1.705520in}}%
\pgfpathlineto{\pgfqpoint{4.017514in}{1.727432in}}%
\pgfpathlineto{\pgfqpoint{4.038639in}{1.761289in}}%
\pgfpathlineto{\pgfqpoint{4.065046in}{1.809411in}}%
\pgfpathlineto{\pgfqpoint{4.154830in}{1.980080in}}%
\pgfpathlineto{\pgfqpoint{4.175955in}{2.012360in}}%
\pgfpathlineto{\pgfqpoint{4.191799in}{2.032906in}}%
\pgfpathlineto{\pgfqpoint{4.207644in}{2.050019in}}%
\pgfpathlineto{\pgfqpoint{4.223488in}{2.063611in}}%
\pgfpathlineto{\pgfqpoint{4.239332in}{2.073770in}}%
\pgfpathlineto{\pgfqpoint{4.255176in}{2.080757in}}%
\pgfpathlineto{\pgfqpoint{4.271020in}{2.085000in}}%
\pgfpathlineto{\pgfqpoint{4.286864in}{2.087085in}}%
\pgfpathlineto{\pgfqpoint{4.313271in}{2.087782in}}%
\pgfpathlineto{\pgfqpoint{4.339678in}{2.088517in}}%
\pgfpathlineto{\pgfqpoint{4.355523in}{2.090817in}}%
\pgfpathlineto{\pgfqpoint{4.371367in}{2.095704in}}%
\pgfpathlineto{\pgfqpoint{4.381929in}{2.100863in}}%
\pgfpathlineto{\pgfqpoint{4.392492in}{2.107850in}}%
\pgfpathlineto{\pgfqpoint{4.403055in}{2.116904in}}%
\pgfpathlineto{\pgfqpoint{4.413618in}{2.128242in}}%
\pgfpathlineto{\pgfqpoint{4.424181in}{2.142050in}}%
\pgfpathlineto{\pgfqpoint{4.434743in}{2.158484in}}%
\pgfpathlineto{\pgfqpoint{4.450588in}{2.188309in}}%
\pgfpathlineto{\pgfqpoint{4.466432in}{2.224534in}}%
\pgfpathlineto{\pgfqpoint{4.482276in}{2.267183in}}%
\pgfpathlineto{\pgfqpoint{4.498120in}{2.316054in}}%
\pgfpathlineto{\pgfqpoint{4.519246in}{2.390132in}}%
\pgfpathlineto{\pgfqpoint{4.540371in}{2.472885in}}%
\pgfpathlineto{\pgfqpoint{4.572060in}{2.608244in}}%
\pgfpathlineto{\pgfqpoint{4.630155in}{2.860572in}}%
\pgfpathlineto{\pgfqpoint{4.651280in}{2.942584in}}%
\pgfpathlineto{\pgfqpoint{4.667125in}{2.997379in}}%
\pgfpathlineto{\pgfqpoint{4.682969in}{3.045005in}}%
\pgfpathlineto{\pgfqpoint{4.698813in}{3.084345in}}%
\pgfpathlineto{\pgfqpoint{4.709376in}{3.105501in}}%
\pgfpathlineto{\pgfqpoint{4.719939in}{3.122331in}}%
\pgfpathlineto{\pgfqpoint{4.730501in}{3.134656in}}%
\pgfpathlineto{\pgfqpoint{4.741064in}{3.142345in}}%
\pgfpathlineto{\pgfqpoint{4.751627in}{3.145314in}}%
\pgfpathlineto{\pgfqpoint{4.756908in}{3.145016in}}%
\pgfpathlineto{\pgfqpoint{4.762190in}{3.143531in}}%
\pgfpathlineto{\pgfqpoint{4.772752in}{3.137017in}}%
\pgfpathlineto{\pgfqpoint{4.783315in}{3.125840in}}%
\pgfpathlineto{\pgfqpoint{4.793878in}{3.110121in}}%
\pgfpathlineto{\pgfqpoint{4.804441in}{3.090028in}}%
\pgfpathlineto{\pgfqpoint{4.815004in}{3.065773in}}%
\pgfpathlineto{\pgfqpoint{4.830848in}{3.022162in}}%
\pgfpathlineto{\pgfqpoint{4.846692in}{2.970802in}}%
\pgfpathlineto{\pgfqpoint{4.867817in}{2.892353in}}%
\pgfpathlineto{\pgfqpoint{4.894224in}{2.782781in}}%
\pgfpathlineto{\pgfqpoint{4.978727in}{2.420255in}}%
\pgfpathlineto{\pgfqpoint{4.999852in}{2.342700in}}%
\pgfpathlineto{\pgfqpoint{5.020978in}{2.274904in}}%
\pgfpathlineto{\pgfqpoint{5.036822in}{2.231199in}}%
\pgfpathlineto{\pgfqpoint{5.052666in}{2.193898in}}%
\pgfpathlineto{\pgfqpoint{5.068510in}{2.163018in}}%
\pgfpathlineto{\pgfqpoint{5.084354in}{2.138357in}}%
\pgfpathlineto{\pgfqpoint{5.094917in}{2.125182in}}%
\pgfpathlineto{\pgfqpoint{5.105480in}{2.114434in}}%
\pgfpathlineto{\pgfqpoint{5.116043in}{2.105918in}}%
\pgfpathlineto{\pgfqpoint{5.126606in}{2.099412in}}%
\pgfpathlineto{\pgfqpoint{5.142450in}{2.092878in}}%
\pgfpathlineto{\pgfqpoint{5.158294in}{2.089405in}}%
\pgfpathlineto{\pgfqpoint{5.179419in}{2.087873in}}%
\pgfpathlineto{\pgfqpoint{5.226952in}{2.086270in}}%
\pgfpathlineto{\pgfqpoint{5.242796in}{2.083188in}}%
\pgfpathlineto{\pgfqpoint{5.258640in}{2.077637in}}%
\pgfpathlineto{\pgfqpoint{5.274485in}{2.069108in}}%
\pgfpathlineto{\pgfqpoint{5.290329in}{2.057254in}}%
\pgfpathlineto{\pgfqpoint{5.306173in}{2.041901in}}%
\pgfpathlineto{\pgfqpoint{5.322017in}{2.023050in}}%
\pgfpathlineto{\pgfqpoint{5.337861in}{2.000872in}}%
\pgfpathlineto{\pgfqpoint{5.358987in}{1.966712in}}%
\pgfpathlineto{\pgfqpoint{5.385394in}{1.918331in}}%
\pgfpathlineto{\pgfqpoint{5.469896in}{1.756798in}}%
\pgfpathlineto{\pgfqpoint{5.491022in}{1.723559in}}%
\pgfpathlineto{\pgfqpoint{5.506866in}{1.702187in}}%
\pgfpathlineto{\pgfqpoint{5.522710in}{1.684201in}}%
\pgfpathlineto{\pgfqpoint{5.538554in}{1.669730in}}%
\pgfpathlineto{\pgfqpoint{5.554398in}{1.658730in}}%
\pgfpathlineto{\pgfqpoint{5.570242in}{1.650985in}}%
\pgfpathlineto{\pgfqpoint{5.586087in}{1.646105in}}%
\pgfpathlineto{\pgfqpoint{5.601931in}{1.643542in}}%
\pgfpathlineto{\pgfqpoint{5.623056in}{1.642520in}}%
\pgfpathlineto{\pgfqpoint{5.660026in}{1.641638in}}%
\pgfpathlineto{\pgfqpoint{5.675870in}{1.639159in}}%
\pgfpathlineto{\pgfqpoint{5.691714in}{1.634014in}}%
\pgfpathlineto{\pgfqpoint{5.702277in}{1.628640in}}%
\pgfpathlineto{\pgfqpoint{5.712840in}{1.621408in}}%
\pgfpathlineto{\pgfqpoint{5.723403in}{1.612079in}}%
\pgfpathlineto{\pgfqpoint{5.733966in}{1.600442in}}%
\pgfpathlineto{\pgfqpoint{5.744528in}{1.586313in}}%
\pgfpathlineto{\pgfqpoint{5.755091in}{1.569542in}}%
\pgfpathlineto{\pgfqpoint{5.770935in}{1.539190in}}%
\pgfpathlineto{\pgfqpoint{5.786779in}{1.502427in}}%
\pgfpathlineto{\pgfqpoint{5.802624in}{1.459249in}}%
\pgfpathlineto{\pgfqpoint{5.818468in}{1.409875in}}%
\pgfpathlineto{\pgfqpoint{5.839593in}{1.335201in}}%
\pgfpathlineto{\pgfqpoint{5.860719in}{1.251974in}}%
\pgfpathlineto{\pgfqpoint{5.892407in}{1.116211in}}%
\pgfpathlineto{\pgfqpoint{5.950503in}{0.864327in}}%
\pgfpathlineto{\pgfqpoint{5.971628in}{0.782877in}}%
\pgfpathlineto{\pgfqpoint{5.987472in}{0.728633in}}%
\pgfpathlineto{\pgfqpoint{6.003316in}{0.681658in}}%
\pgfpathlineto{\pgfqpoint{6.019161in}{0.643053in}}%
\pgfpathlineto{\pgfqpoint{6.029723in}{0.622426in}}%
\pgfpathlineto{\pgfqpoint{6.040286in}{0.606150in}}%
\pgfpathlineto{\pgfqpoint{6.050849in}{0.594398in}}%
\pgfpathlineto{\pgfqpoint{6.061412in}{0.587296in}}%
\pgfpathlineto{\pgfqpoint{6.066693in}{0.585515in}}%
\pgfpathlineto{\pgfqpoint{6.071975in}{0.584920in}}%
\pgfpathlineto{\pgfqpoint{6.071975in}{0.584920in}}%
\pgfusepath{stroke}%
\end{pgfscope}%
\begin{pgfscope}%
\pgfpathrectangle{\pgfqpoint{0.532060in}{0.456901in}}{\pgfqpoint{5.803720in}{2.816433in}}%
\pgfusepath{clip}%
\pgfsetbuttcap%
\pgfsetroundjoin%
\pgfsetlinewidth{1.003750pt}%
\definecolor{currentstroke}{rgb}{0.000000,0.000000,1.000000}%
\pgfsetstrokecolor{currentstroke}%
\pgfsetdash{{1.000000pt}{1.650000pt}}{0.000000pt}%
\pgfpathmoveto{\pgfqpoint{0.795865in}{1.030213in}}%
\pgfpathlineto{\pgfqpoint{0.816991in}{1.031269in}}%
\pgfpathlineto{\pgfqpoint{0.838116in}{1.034437in}}%
\pgfpathlineto{\pgfqpoint{0.859242in}{1.039707in}}%
\pgfpathlineto{\pgfqpoint{0.880367in}{1.047066in}}%
\pgfpathlineto{\pgfqpoint{0.901493in}{1.056496in}}%
\pgfpathlineto{\pgfqpoint{0.922619in}{1.067972in}}%
\pgfpathlineto{\pgfqpoint{0.949026in}{1.085152in}}%
\pgfpathlineto{\pgfqpoint{0.975432in}{1.105415in}}%
\pgfpathlineto{\pgfqpoint{1.001839in}{1.128683in}}%
\pgfpathlineto{\pgfqpoint{1.028246in}{1.154864in}}%
\pgfpathlineto{\pgfqpoint{1.054653in}{1.183853in}}%
\pgfpathlineto{\pgfqpoint{1.086342in}{1.222185in}}%
\pgfpathlineto{\pgfqpoint{1.118030in}{1.264178in}}%
\pgfpathlineto{\pgfqpoint{1.149718in}{1.309592in}}%
\pgfpathlineto{\pgfqpoint{1.186688in}{1.366554in}}%
\pgfpathlineto{\pgfqpoint{1.228939in}{1.436357in}}%
\pgfpathlineto{\pgfqpoint{1.271190in}{1.510498in}}%
\pgfpathlineto{\pgfqpoint{1.324004in}{1.608129in}}%
\pgfpathlineto{\pgfqpoint{1.387381in}{1.730507in}}%
\pgfpathlineto{\pgfqpoint{1.609199in}{2.164223in}}%
\pgfpathlineto{\pgfqpoint{1.656732in}{2.250369in}}%
\pgfpathlineto{\pgfqpoint{1.698983in}{2.322836in}}%
\pgfpathlineto{\pgfqpoint{1.741234in}{2.390672in}}%
\pgfpathlineto{\pgfqpoint{1.778204in}{2.445688in}}%
\pgfpathlineto{\pgfqpoint{1.809892in}{2.489280in}}%
\pgfpathlineto{\pgfqpoint{1.841581in}{2.529319in}}%
\pgfpathlineto{\pgfqpoint{1.873269in}{2.565575in}}%
\pgfpathlineto{\pgfqpoint{1.899676in}{2.592751in}}%
\pgfpathlineto{\pgfqpoint{1.926083in}{2.617049in}}%
\pgfpathlineto{\pgfqpoint{1.952490in}{2.638374in}}%
\pgfpathlineto{\pgfqpoint{1.978897in}{2.656642in}}%
\pgfpathlineto{\pgfqpoint{2.005304in}{2.671779in}}%
\pgfpathlineto{\pgfqpoint{2.026429in}{2.681595in}}%
\pgfpathlineto{\pgfqpoint{2.047555in}{2.689344in}}%
\pgfpathlineto{\pgfqpoint{2.068680in}{2.695006in}}%
\pgfpathlineto{\pgfqpoint{2.089806in}{2.698569in}}%
\pgfpathlineto{\pgfqpoint{2.110931in}{2.700021in}}%
\pgfpathlineto{\pgfqpoint{2.132057in}{2.699361in}}%
\pgfpathlineto{\pgfqpoint{2.153183in}{2.696589in}}%
\pgfpathlineto{\pgfqpoint{2.174308in}{2.691712in}}%
\pgfpathlineto{\pgfqpoint{2.195434in}{2.684744in}}%
\pgfpathlineto{\pgfqpoint{2.216559in}{2.675700in}}%
\pgfpathlineto{\pgfqpoint{2.237685in}{2.664606in}}%
\pgfpathlineto{\pgfqpoint{2.264092in}{2.647895in}}%
\pgfpathlineto{\pgfqpoint{2.290499in}{2.628089in}}%
\pgfpathlineto{\pgfqpoint{2.316906in}{2.605266in}}%
\pgfpathlineto{\pgfqpoint{2.343313in}{2.579516in}}%
\pgfpathlineto{\pgfqpoint{2.369720in}{2.550942in}}%
\pgfpathlineto{\pgfqpoint{2.401408in}{2.513083in}}%
\pgfpathlineto{\pgfqpoint{2.433096in}{2.471535in}}%
\pgfpathlineto{\pgfqpoint{2.464785in}{2.426534in}}%
\pgfpathlineto{\pgfqpoint{2.501754in}{2.370013in}}%
\pgfpathlineto{\pgfqpoint{2.538724in}{2.309580in}}%
\pgfpathlineto{\pgfqpoint{2.580975in}{2.236325in}}%
\pgfpathlineto{\pgfqpoint{2.628508in}{2.149462in}}%
\pgfpathlineto{\pgfqpoint{2.686603in}{2.038469in}}%
\pgfpathlineto{\pgfqpoint{2.776387in}{1.861197in}}%
\pgfpathlineto{\pgfqpoint{2.882015in}{1.653483in}}%
\pgfpathlineto{\pgfqpoint{2.940110in}{1.544109in}}%
\pgfpathlineto{\pgfqpoint{2.987642in}{1.459095in}}%
\pgfpathlineto{\pgfqpoint{3.029893in}{1.387858in}}%
\pgfpathlineto{\pgfqpoint{3.066863in}{1.329464in}}%
\pgfpathlineto{\pgfqpoint{3.103833in}{1.275220in}}%
\pgfpathlineto{\pgfqpoint{3.135521in}{1.232349in}}%
\pgfpathlineto{\pgfqpoint{3.167210in}{1.193080in}}%
\pgfpathlineto{\pgfqpoint{3.198898in}{1.157638in}}%
\pgfpathlineto{\pgfqpoint{3.225305in}{1.131172in}}%
\pgfpathlineto{\pgfqpoint{3.251712in}{1.107608in}}%
\pgfpathlineto{\pgfqpoint{3.278119in}{1.087040in}}%
\pgfpathlineto{\pgfqpoint{3.304526in}{1.069549in}}%
\pgfpathlineto{\pgfqpoint{3.330933in}{1.055204in}}%
\pgfpathlineto{\pgfqpoint{3.352058in}{1.046032in}}%
\pgfpathlineto{\pgfqpoint{3.373184in}{1.038933in}}%
\pgfpathlineto{\pgfqpoint{3.394309in}{1.033926in}}%
\pgfpathlineto{\pgfqpoint{3.415435in}{1.031022in}}%
\pgfpathlineto{\pgfqpoint{3.436561in}{1.030229in}}%
\pgfpathlineto{\pgfqpoint{3.457686in}{1.031550in}}%
\pgfpathlineto{\pgfqpoint{3.478812in}{1.034981in}}%
\pgfpathlineto{\pgfqpoint{3.499937in}{1.040513in}}%
\pgfpathlineto{\pgfqpoint{3.521063in}{1.048132in}}%
\pgfpathlineto{\pgfqpoint{3.542188in}{1.057819in}}%
\pgfpathlineto{\pgfqpoint{3.563314in}{1.069549in}}%
\pgfpathlineto{\pgfqpoint{3.589721in}{1.087040in}}%
\pgfpathlineto{\pgfqpoint{3.616128in}{1.107608in}}%
\pgfpathlineto{\pgfqpoint{3.642535in}{1.131172in}}%
\pgfpathlineto{\pgfqpoint{3.668942in}{1.157638in}}%
\pgfpathlineto{\pgfqpoint{3.695349in}{1.186902in}}%
\pgfpathlineto{\pgfqpoint{3.727037in}{1.225548in}}%
\pgfpathlineto{\pgfqpoint{3.758725in}{1.267835in}}%
\pgfpathlineto{\pgfqpoint{3.790414in}{1.313523in}}%
\pgfpathlineto{\pgfqpoint{3.827383in}{1.370776in}}%
\pgfpathlineto{\pgfqpoint{3.869635in}{1.440871in}}%
\pgfpathlineto{\pgfqpoint{3.911886in}{1.515259in}}%
\pgfpathlineto{\pgfqpoint{3.964700in}{1.613131in}}%
\pgfpathlineto{\pgfqpoint{4.028076in}{1.735692in}}%
\pgfpathlineto{\pgfqpoint{4.244613in}{2.159314in}}%
\pgfpathlineto{\pgfqpoint{4.292146in}{2.245703in}}%
\pgfpathlineto{\pgfqpoint{4.334397in}{2.318435in}}%
\pgfpathlineto{\pgfqpoint{4.376648in}{2.386581in}}%
\pgfpathlineto{\pgfqpoint{4.413618in}{2.441903in}}%
\pgfpathlineto{\pgfqpoint{4.445306in}{2.485780in}}%
\pgfpathlineto{\pgfqpoint{4.476995in}{2.526123in}}%
\pgfpathlineto{\pgfqpoint{4.508683in}{2.562704in}}%
\pgfpathlineto{\pgfqpoint{4.535090in}{2.590161in}}%
\pgfpathlineto{\pgfqpoint{4.561497in}{2.614752in}}%
\pgfpathlineto{\pgfqpoint{4.587904in}{2.636378in}}%
\pgfpathlineto{\pgfqpoint{4.614311in}{2.654955in}}%
\pgfpathlineto{\pgfqpoint{4.640718in}{2.670408in}}%
\pgfpathlineto{\pgfqpoint{4.661843in}{2.680480in}}%
\pgfpathlineto{\pgfqpoint{4.682969in}{2.688489in}}%
\pgfpathlineto{\pgfqpoint{4.704094in}{2.694413in}}%
\pgfpathlineto{\pgfqpoint{4.725220in}{2.698239in}}%
\pgfpathlineto{\pgfqpoint{4.746345in}{2.699955in}}%
\pgfpathlineto{\pgfqpoint{4.767471in}{2.699559in}}%
\pgfpathlineto{\pgfqpoint{4.788597in}{2.697051in}}%
\pgfpathlineto{\pgfqpoint{4.809722in}{2.692437in}}%
\pgfpathlineto{\pgfqpoint{4.830848in}{2.685729in}}%
\pgfpathlineto{\pgfqpoint{4.851973in}{2.676944in}}%
\pgfpathlineto{\pgfqpoint{4.873099in}{2.666104in}}%
\pgfpathlineto{\pgfqpoint{4.899506in}{2.649707in}}%
\pgfpathlineto{\pgfqpoint{4.925913in}{2.630207in}}%
\pgfpathlineto{\pgfqpoint{4.952320in}{2.607682in}}%
\pgfpathlineto{\pgfqpoint{4.978727in}{2.582220in}}%
\pgfpathlineto{\pgfqpoint{5.005134in}{2.553923in}}%
\pgfpathlineto{\pgfqpoint{5.036822in}{2.516382in}}%
\pgfpathlineto{\pgfqpoint{5.068510in}{2.475133in}}%
\pgfpathlineto{\pgfqpoint{5.100199in}{2.430410in}}%
\pgfpathlineto{\pgfqpoint{5.137168in}{2.374185in}}%
\pgfpathlineto{\pgfqpoint{5.174138in}{2.314016in}}%
\pgfpathlineto{\pgfqpoint{5.216389in}{2.241021in}}%
\pgfpathlineto{\pgfqpoint{5.263922in}{2.154394in}}%
\pgfpathlineto{\pgfqpoint{5.322017in}{2.043602in}}%
\pgfpathlineto{\pgfqpoint{5.406519in}{1.876951in}}%
\pgfpathlineto{\pgfqpoint{5.517429in}{1.658567in}}%
\pgfpathlineto{\pgfqpoint{5.575524in}{1.548963in}}%
\pgfpathlineto{\pgfqpoint{5.623056in}{1.463691in}}%
\pgfpathlineto{\pgfqpoint{5.665307in}{1.392176in}}%
\pgfpathlineto{\pgfqpoint{5.702277in}{1.333502in}}%
\pgfpathlineto{\pgfqpoint{5.739247in}{1.278948in}}%
\pgfpathlineto{\pgfqpoint{5.770935in}{1.235787in}}%
\pgfpathlineto{\pgfqpoint{5.802624in}{1.196209in}}%
\pgfpathlineto{\pgfqpoint{5.834312in}{1.160440in}}%
\pgfpathlineto{\pgfqpoint{5.860719in}{1.133690in}}%
\pgfpathlineto{\pgfqpoint{5.887126in}{1.109831in}}%
\pgfpathlineto{\pgfqpoint{5.913533in}{1.088960in}}%
\pgfpathlineto{\pgfqpoint{5.939940in}{1.071158in}}%
\pgfpathlineto{\pgfqpoint{5.966347in}{1.056496in}}%
\pgfpathlineto{\pgfqpoint{5.987472in}{1.047066in}}%
\pgfpathlineto{\pgfqpoint{6.008598in}{1.039707in}}%
\pgfpathlineto{\pgfqpoint{6.029723in}{1.034437in}}%
\pgfpathlineto{\pgfqpoint{6.050849in}{1.031269in}}%
\pgfpathlineto{\pgfqpoint{6.071975in}{1.030213in}}%
\pgfpathlineto{\pgfqpoint{6.071975in}{1.030213in}}%
\pgfusepath{stroke}%
\end{pgfscope}%
\begin{pgfscope}%
\pgfpathrectangle{\pgfqpoint{0.532060in}{0.456901in}}{\pgfqpoint{5.803720in}{2.816433in}}%
\pgfusepath{clip}%
\pgfsetbuttcap%
\pgfsetroundjoin%
\pgfsetlinewidth{1.003750pt}%
\definecolor{currentstroke}{rgb}{0.000000,0.501961,0.000000}%
\pgfsetstrokecolor{currentstroke}%
\pgfsetdash{{1.000000pt}{1.650000pt}}{0.000000pt}%
\pgfpathmoveto{\pgfqpoint{0.795865in}{1.586828in}}%
\pgfpathlineto{\pgfqpoint{0.806428in}{1.587620in}}%
\pgfpathlineto{\pgfqpoint{0.816991in}{1.589993in}}%
\pgfpathlineto{\pgfqpoint{0.832835in}{1.596482in}}%
\pgfpathlineto{\pgfqpoint{0.848679in}{1.606410in}}%
\pgfpathlineto{\pgfqpoint{0.864523in}{1.619652in}}%
\pgfpathlineto{\pgfqpoint{0.880367in}{1.636036in}}%
\pgfpathlineto{\pgfqpoint{0.896212in}{1.655353in}}%
\pgfpathlineto{\pgfqpoint{0.917337in}{1.685239in}}%
\pgfpathlineto{\pgfqpoint{0.938463in}{1.719217in}}%
\pgfpathlineto{\pgfqpoint{0.964870in}{1.766258in}}%
\pgfpathlineto{\pgfqpoint{1.001839in}{1.837612in}}%
\pgfpathlineto{\pgfqpoint{1.070498in}{1.971336in}}%
\pgfpathlineto{\pgfqpoint{1.096904in}{2.017702in}}%
\pgfpathlineto{\pgfqpoint{1.118030in}{2.050969in}}%
\pgfpathlineto{\pgfqpoint{1.139156in}{2.080010in}}%
\pgfpathlineto{\pgfqpoint{1.155000in}{2.098616in}}%
\pgfpathlineto{\pgfqpoint{1.170844in}{2.114232in}}%
\pgfpathlineto{\pgfqpoint{1.186688in}{2.126659in}}%
\pgfpathlineto{\pgfqpoint{1.202532in}{2.135738in}}%
\pgfpathlineto{\pgfqpoint{1.218376in}{2.141353in}}%
\pgfpathlineto{\pgfqpoint{1.228939in}{2.143134in}}%
\pgfpathlineto{\pgfqpoint{1.239502in}{2.143332in}}%
\pgfpathlineto{\pgfqpoint{1.250065in}{2.141946in}}%
\pgfpathlineto{\pgfqpoint{1.260628in}{2.138984in}}%
\pgfpathlineto{\pgfqpoint{1.276472in}{2.131626in}}%
\pgfpathlineto{\pgfqpoint{1.292316in}{2.120855in}}%
\pgfpathlineto{\pgfqpoint{1.308160in}{2.106811in}}%
\pgfpathlineto{\pgfqpoint{1.324004in}{2.089672in}}%
\pgfpathlineto{\pgfqpoint{1.339848in}{2.069659in}}%
\pgfpathlineto{\pgfqpoint{1.360974in}{2.038951in}}%
\pgfpathlineto{\pgfqpoint{1.382100in}{2.004289in}}%
\pgfpathlineto{\pgfqpoint{1.408507in}{1.956612in}}%
\pgfpathlineto{\pgfqpoint{1.445476in}{1.884811in}}%
\pgfpathlineto{\pgfqpoint{1.508853in}{1.761367in}}%
\pgfpathlineto{\pgfqpoint{1.535260in}{1.714772in}}%
\pgfpathlineto{\pgfqpoint{1.556385in}{1.681265in}}%
\pgfpathlineto{\pgfqpoint{1.577511in}{1.651940in}}%
\pgfpathlineto{\pgfqpoint{1.593355in}{1.633095in}}%
\pgfpathlineto{\pgfqpoint{1.609199in}{1.617221in}}%
\pgfpathlineto{\pgfqpoint{1.625044in}{1.604522in}}%
\pgfpathlineto{\pgfqpoint{1.640888in}{1.595159in}}%
\pgfpathlineto{\pgfqpoint{1.656732in}{1.589252in}}%
\pgfpathlineto{\pgfqpoint{1.667295in}{1.587274in}}%
\pgfpathlineto{\pgfqpoint{1.677857in}{1.586878in}}%
\pgfpathlineto{\pgfqpoint{1.688420in}{1.588066in}}%
\pgfpathlineto{\pgfqpoint{1.698983in}{1.590831in}}%
\pgfpathlineto{\pgfqpoint{1.714827in}{1.597901in}}%
\pgfpathlineto{\pgfqpoint{1.730671in}{1.608391in}}%
\pgfpathlineto{\pgfqpoint{1.746516in}{1.622169in}}%
\pgfpathlineto{\pgfqpoint{1.762360in}{1.639058in}}%
\pgfpathlineto{\pgfqpoint{1.778204in}{1.658841in}}%
\pgfpathlineto{\pgfqpoint{1.799329in}{1.689278in}}%
\pgfpathlineto{\pgfqpoint{1.820455in}{1.723714in}}%
\pgfpathlineto{\pgfqpoint{1.846862in}{1.771184in}}%
\pgfpathlineto{\pgfqpoint{1.883832in}{1.842842in}}%
\pgfpathlineto{\pgfqpoint{1.947208in}{1.966464in}}%
\pgfpathlineto{\pgfqpoint{1.973615in}{2.013284in}}%
\pgfpathlineto{\pgfqpoint{1.994741in}{2.047027in}}%
\pgfpathlineto{\pgfqpoint{2.015866in}{2.076635in}}%
\pgfpathlineto{\pgfqpoint{2.031711in}{2.095716in}}%
\pgfpathlineto{\pgfqpoint{2.047555in}{2.111846in}}%
\pgfpathlineto{\pgfqpoint{2.063399in}{2.124817in}}%
\pgfpathlineto{\pgfqpoint{2.079243in}{2.134463in}}%
\pgfpathlineto{\pgfqpoint{2.095087in}{2.140661in}}%
\pgfpathlineto{\pgfqpoint{2.105650in}{2.142837in}}%
\pgfpathlineto{\pgfqpoint{2.116213in}{2.143431in}}%
\pgfpathlineto{\pgfqpoint{2.126776in}{2.142441in}}%
\pgfpathlineto{\pgfqpoint{2.137338in}{2.139872in}}%
\pgfpathlineto{\pgfqpoint{2.153183in}{2.133092in}}%
\pgfpathlineto{\pgfqpoint{2.169027in}{2.122882in}}%
\pgfpathlineto{\pgfqpoint{2.184871in}{2.109372in}}%
\pgfpathlineto{\pgfqpoint{2.200715in}{2.092735in}}%
\pgfpathlineto{\pgfqpoint{2.216559in}{2.073184in}}%
\pgfpathlineto{\pgfqpoint{2.237685in}{2.043021in}}%
\pgfpathlineto{\pgfqpoint{2.258810in}{2.008812in}}%
\pgfpathlineto{\pgfqpoint{2.285217in}{1.961555in}}%
\pgfpathlineto{\pgfqpoint{2.322187in}{1.890046in}}%
\pgfpathlineto{\pgfqpoint{2.385564in}{1.766258in}}%
\pgfpathlineto{\pgfqpoint{2.411971in}{1.719217in}}%
\pgfpathlineto{\pgfqpoint{2.433096in}{1.685239in}}%
\pgfpathlineto{\pgfqpoint{2.454222in}{1.655353in}}%
\pgfpathlineto{\pgfqpoint{2.470066in}{1.636036in}}%
\pgfpathlineto{\pgfqpoint{2.485910in}{1.619652in}}%
\pgfpathlineto{\pgfqpoint{2.501754in}{1.606410in}}%
\pgfpathlineto{\pgfqpoint{2.517599in}{1.596482in}}%
\pgfpathlineto{\pgfqpoint{2.533443in}{1.589993in}}%
\pgfpathlineto{\pgfqpoint{2.544006in}{1.587620in}}%
\pgfpathlineto{\pgfqpoint{2.554568in}{1.586828in}}%
\pgfpathlineto{\pgfqpoint{2.565131in}{1.587620in}}%
\pgfpathlineto{\pgfqpoint{2.575694in}{1.589993in}}%
\pgfpathlineto{\pgfqpoint{2.591538in}{1.596482in}}%
\pgfpathlineto{\pgfqpoint{2.607382in}{1.606410in}}%
\pgfpathlineto{\pgfqpoint{2.623226in}{1.619652in}}%
\pgfpathlineto{\pgfqpoint{2.639071in}{1.636036in}}%
\pgfpathlineto{\pgfqpoint{2.654915in}{1.655353in}}%
\pgfpathlineto{\pgfqpoint{2.676040in}{1.685239in}}%
\pgfpathlineto{\pgfqpoint{2.697166in}{1.719217in}}%
\pgfpathlineto{\pgfqpoint{2.723573in}{1.766258in}}%
\pgfpathlineto{\pgfqpoint{2.760543in}{1.837612in}}%
\pgfpathlineto{\pgfqpoint{2.829201in}{1.971336in}}%
\pgfpathlineto{\pgfqpoint{2.855608in}{2.017702in}}%
\pgfpathlineto{\pgfqpoint{2.876733in}{2.050969in}}%
\pgfpathlineto{\pgfqpoint{2.897859in}{2.080010in}}%
\pgfpathlineto{\pgfqpoint{2.913703in}{2.098616in}}%
\pgfpathlineto{\pgfqpoint{2.929547in}{2.114232in}}%
\pgfpathlineto{\pgfqpoint{2.945391in}{2.126659in}}%
\pgfpathlineto{\pgfqpoint{2.961235in}{2.135738in}}%
\pgfpathlineto{\pgfqpoint{2.977080in}{2.141353in}}%
\pgfpathlineto{\pgfqpoint{2.987642in}{2.143134in}}%
\pgfpathlineto{\pgfqpoint{2.998205in}{2.143332in}}%
\pgfpathlineto{\pgfqpoint{3.008768in}{2.141946in}}%
\pgfpathlineto{\pgfqpoint{3.019331in}{2.138984in}}%
\pgfpathlineto{\pgfqpoint{3.035175in}{2.131626in}}%
\pgfpathlineto{\pgfqpoint{3.051019in}{2.120855in}}%
\pgfpathlineto{\pgfqpoint{3.066863in}{2.106811in}}%
\pgfpathlineto{\pgfqpoint{3.082707in}{2.089672in}}%
\pgfpathlineto{\pgfqpoint{3.098552in}{2.069659in}}%
\pgfpathlineto{\pgfqpoint{3.119677in}{2.038951in}}%
\pgfpathlineto{\pgfqpoint{3.140803in}{2.004289in}}%
\pgfpathlineto{\pgfqpoint{3.167210in}{1.956612in}}%
\pgfpathlineto{\pgfqpoint{3.204179in}{1.884811in}}%
\pgfpathlineto{\pgfqpoint{3.267556in}{1.761367in}}%
\pgfpathlineto{\pgfqpoint{3.293963in}{1.714772in}}%
\pgfpathlineto{\pgfqpoint{3.315089in}{1.681265in}}%
\pgfpathlineto{\pgfqpoint{3.336214in}{1.651940in}}%
\pgfpathlineto{\pgfqpoint{3.352058in}{1.633095in}}%
\pgfpathlineto{\pgfqpoint{3.367902in}{1.617221in}}%
\pgfpathlineto{\pgfqpoint{3.383747in}{1.604522in}}%
\pgfpathlineto{\pgfqpoint{3.399591in}{1.595159in}}%
\pgfpathlineto{\pgfqpoint{3.415435in}{1.589252in}}%
\pgfpathlineto{\pgfqpoint{3.425998in}{1.587274in}}%
\pgfpathlineto{\pgfqpoint{3.436561in}{1.586878in}}%
\pgfpathlineto{\pgfqpoint{3.447123in}{1.588066in}}%
\pgfpathlineto{\pgfqpoint{3.457686in}{1.590831in}}%
\pgfpathlineto{\pgfqpoint{3.473530in}{1.597901in}}%
\pgfpathlineto{\pgfqpoint{3.489374in}{1.608391in}}%
\pgfpathlineto{\pgfqpoint{3.505219in}{1.622169in}}%
\pgfpathlineto{\pgfqpoint{3.521063in}{1.639058in}}%
\pgfpathlineto{\pgfqpoint{3.536907in}{1.658841in}}%
\pgfpathlineto{\pgfqpoint{3.558033in}{1.689278in}}%
\pgfpathlineto{\pgfqpoint{3.579158in}{1.723714in}}%
\pgfpathlineto{\pgfqpoint{3.605565in}{1.771184in}}%
\pgfpathlineto{\pgfqpoint{3.642535in}{1.842842in}}%
\pgfpathlineto{\pgfqpoint{3.705911in}{1.966464in}}%
\pgfpathlineto{\pgfqpoint{3.732318in}{2.013284in}}%
\pgfpathlineto{\pgfqpoint{3.753444in}{2.047027in}}%
\pgfpathlineto{\pgfqpoint{3.774570in}{2.076635in}}%
\pgfpathlineto{\pgfqpoint{3.790414in}{2.095716in}}%
\pgfpathlineto{\pgfqpoint{3.806258in}{2.111846in}}%
\pgfpathlineto{\pgfqpoint{3.822102in}{2.124817in}}%
\pgfpathlineto{\pgfqpoint{3.837946in}{2.134463in}}%
\pgfpathlineto{\pgfqpoint{3.853790in}{2.140661in}}%
\pgfpathlineto{\pgfqpoint{3.864353in}{2.142837in}}%
\pgfpathlineto{\pgfqpoint{3.874916in}{2.143431in}}%
\pgfpathlineto{\pgfqpoint{3.885479in}{2.142441in}}%
\pgfpathlineto{\pgfqpoint{3.896042in}{2.139872in}}%
\pgfpathlineto{\pgfqpoint{3.911886in}{2.133092in}}%
\pgfpathlineto{\pgfqpoint{3.927730in}{2.122882in}}%
\pgfpathlineto{\pgfqpoint{3.943574in}{2.109372in}}%
\pgfpathlineto{\pgfqpoint{3.959418in}{2.092735in}}%
\pgfpathlineto{\pgfqpoint{3.975262in}{2.073184in}}%
\pgfpathlineto{\pgfqpoint{3.996388in}{2.043021in}}%
\pgfpathlineto{\pgfqpoint{4.017514in}{2.008812in}}%
\pgfpathlineto{\pgfqpoint{4.043920in}{1.961555in}}%
\pgfpathlineto{\pgfqpoint{4.080890in}{1.890046in}}%
\pgfpathlineto{\pgfqpoint{4.144267in}{1.766258in}}%
\pgfpathlineto{\pgfqpoint{4.170674in}{1.719217in}}%
\pgfpathlineto{\pgfqpoint{4.191799in}{1.685239in}}%
\pgfpathlineto{\pgfqpoint{4.212925in}{1.655353in}}%
\pgfpathlineto{\pgfqpoint{4.228769in}{1.636036in}}%
\pgfpathlineto{\pgfqpoint{4.244613in}{1.619652in}}%
\pgfpathlineto{\pgfqpoint{4.260458in}{1.606410in}}%
\pgfpathlineto{\pgfqpoint{4.276302in}{1.596482in}}%
\pgfpathlineto{\pgfqpoint{4.292146in}{1.589993in}}%
\pgfpathlineto{\pgfqpoint{4.302709in}{1.587620in}}%
\pgfpathlineto{\pgfqpoint{4.313271in}{1.586828in}}%
\pgfpathlineto{\pgfqpoint{4.323834in}{1.587620in}}%
\pgfpathlineto{\pgfqpoint{4.334397in}{1.589993in}}%
\pgfpathlineto{\pgfqpoint{4.350241in}{1.596482in}}%
\pgfpathlineto{\pgfqpoint{4.366085in}{1.606410in}}%
\pgfpathlineto{\pgfqpoint{4.381929in}{1.619652in}}%
\pgfpathlineto{\pgfqpoint{4.397774in}{1.636036in}}%
\pgfpathlineto{\pgfqpoint{4.413618in}{1.655353in}}%
\pgfpathlineto{\pgfqpoint{4.434743in}{1.685239in}}%
\pgfpathlineto{\pgfqpoint{4.455869in}{1.719217in}}%
\pgfpathlineto{\pgfqpoint{4.482276in}{1.766258in}}%
\pgfpathlineto{\pgfqpoint{4.519246in}{1.837612in}}%
\pgfpathlineto{\pgfqpoint{4.587904in}{1.971336in}}%
\pgfpathlineto{\pgfqpoint{4.614311in}{2.017702in}}%
\pgfpathlineto{\pgfqpoint{4.635436in}{2.050969in}}%
\pgfpathlineto{\pgfqpoint{4.656562in}{2.080010in}}%
\pgfpathlineto{\pgfqpoint{4.672406in}{2.098616in}}%
\pgfpathlineto{\pgfqpoint{4.688250in}{2.114232in}}%
\pgfpathlineto{\pgfqpoint{4.704094in}{2.126659in}}%
\pgfpathlineto{\pgfqpoint{4.719939in}{2.135738in}}%
\pgfpathlineto{\pgfqpoint{4.735783in}{2.141353in}}%
\pgfpathlineto{\pgfqpoint{4.746345in}{2.143134in}}%
\pgfpathlineto{\pgfqpoint{4.756908in}{2.143332in}}%
\pgfpathlineto{\pgfqpoint{4.767471in}{2.141946in}}%
\pgfpathlineto{\pgfqpoint{4.778034in}{2.138984in}}%
\pgfpathlineto{\pgfqpoint{4.793878in}{2.131626in}}%
\pgfpathlineto{\pgfqpoint{4.809722in}{2.120855in}}%
\pgfpathlineto{\pgfqpoint{4.825566in}{2.106811in}}%
\pgfpathlineto{\pgfqpoint{4.841410in}{2.089672in}}%
\pgfpathlineto{\pgfqpoint{4.857255in}{2.069659in}}%
\pgfpathlineto{\pgfqpoint{4.878380in}{2.038951in}}%
\pgfpathlineto{\pgfqpoint{4.899506in}{2.004289in}}%
\pgfpathlineto{\pgfqpoint{4.925913in}{1.956612in}}%
\pgfpathlineto{\pgfqpoint{4.962882in}{1.884811in}}%
\pgfpathlineto{\pgfqpoint{5.026259in}{1.761367in}}%
\pgfpathlineto{\pgfqpoint{5.052666in}{1.714772in}}%
\pgfpathlineto{\pgfqpoint{5.073792in}{1.681265in}}%
\pgfpathlineto{\pgfqpoint{5.094917in}{1.651940in}}%
\pgfpathlineto{\pgfqpoint{5.110761in}{1.633095in}}%
\pgfpathlineto{\pgfqpoint{5.126606in}{1.617221in}}%
\pgfpathlineto{\pgfqpoint{5.142450in}{1.604522in}}%
\pgfpathlineto{\pgfqpoint{5.158294in}{1.595159in}}%
\pgfpathlineto{\pgfqpoint{5.174138in}{1.589252in}}%
\pgfpathlineto{\pgfqpoint{5.184701in}{1.587274in}}%
\pgfpathlineto{\pgfqpoint{5.195264in}{1.586878in}}%
\pgfpathlineto{\pgfqpoint{5.205826in}{1.588066in}}%
\pgfpathlineto{\pgfqpoint{5.216389in}{1.590831in}}%
\pgfpathlineto{\pgfqpoint{5.232233in}{1.597901in}}%
\pgfpathlineto{\pgfqpoint{5.248078in}{1.608391in}}%
\pgfpathlineto{\pgfqpoint{5.263922in}{1.622169in}}%
\pgfpathlineto{\pgfqpoint{5.279766in}{1.639058in}}%
\pgfpathlineto{\pgfqpoint{5.295610in}{1.658841in}}%
\pgfpathlineto{\pgfqpoint{5.316736in}{1.689278in}}%
\pgfpathlineto{\pgfqpoint{5.337861in}{1.723714in}}%
\pgfpathlineto{\pgfqpoint{5.364268in}{1.771184in}}%
\pgfpathlineto{\pgfqpoint{5.401238in}{1.842842in}}%
\pgfpathlineto{\pgfqpoint{5.464615in}{1.966464in}}%
\pgfpathlineto{\pgfqpoint{5.491022in}{2.013284in}}%
\pgfpathlineto{\pgfqpoint{5.512147in}{2.047027in}}%
\pgfpathlineto{\pgfqpoint{5.533273in}{2.076635in}}%
\pgfpathlineto{\pgfqpoint{5.549117in}{2.095716in}}%
\pgfpathlineto{\pgfqpoint{5.564961in}{2.111846in}}%
\pgfpathlineto{\pgfqpoint{5.580805in}{2.124817in}}%
\pgfpathlineto{\pgfqpoint{5.596649in}{2.134463in}}%
\pgfpathlineto{\pgfqpoint{5.612494in}{2.140661in}}%
\pgfpathlineto{\pgfqpoint{5.623056in}{2.142837in}}%
\pgfpathlineto{\pgfqpoint{5.633619in}{2.143431in}}%
\pgfpathlineto{\pgfqpoint{5.644182in}{2.142441in}}%
\pgfpathlineto{\pgfqpoint{5.654745in}{2.139872in}}%
\pgfpathlineto{\pgfqpoint{5.670589in}{2.133092in}}%
\pgfpathlineto{\pgfqpoint{5.686433in}{2.122882in}}%
\pgfpathlineto{\pgfqpoint{5.702277in}{2.109372in}}%
\pgfpathlineto{\pgfqpoint{5.718121in}{2.092735in}}%
\pgfpathlineto{\pgfqpoint{5.733966in}{2.073184in}}%
\pgfpathlineto{\pgfqpoint{5.755091in}{2.043021in}}%
\pgfpathlineto{\pgfqpoint{5.776217in}{2.008812in}}%
\pgfpathlineto{\pgfqpoint{5.802624in}{1.961555in}}%
\pgfpathlineto{\pgfqpoint{5.839593in}{1.890046in}}%
\pgfpathlineto{\pgfqpoint{5.902970in}{1.766258in}}%
\pgfpathlineto{\pgfqpoint{5.929377in}{1.719217in}}%
\pgfpathlineto{\pgfqpoint{5.950503in}{1.685239in}}%
\pgfpathlineto{\pgfqpoint{5.971628in}{1.655353in}}%
\pgfpathlineto{\pgfqpoint{5.987472in}{1.636036in}}%
\pgfpathlineto{\pgfqpoint{6.003316in}{1.619652in}}%
\pgfpathlineto{\pgfqpoint{6.019161in}{1.606410in}}%
\pgfpathlineto{\pgfqpoint{6.035005in}{1.596482in}}%
\pgfpathlineto{\pgfqpoint{6.050849in}{1.589993in}}%
\pgfpathlineto{\pgfqpoint{6.061412in}{1.587620in}}%
\pgfpathlineto{\pgfqpoint{6.071975in}{1.586828in}}%
\pgfpathlineto{\pgfqpoint{6.071975in}{1.586828in}}%
\pgfusepath{stroke}%
\end{pgfscope}%
\begin{pgfscope}%
\pgfpathrectangle{\pgfqpoint{0.532060in}{0.456901in}}{\pgfqpoint{5.803720in}{2.816433in}}%
\pgfusepath{clip}%
\pgfsetbuttcap%
\pgfsetroundjoin%
\pgfsetlinewidth{1.003750pt}%
\definecolor{currentstroke}{rgb}{0.501961,0.000000,0.501961}%
\pgfsetstrokecolor{currentstroke}%
\pgfsetdash{{1.000000pt}{1.650000pt}}{0.000000pt}%
\pgfpathmoveto{\pgfqpoint{0.795865in}{1.698151in}}%
\pgfpathlineto{\pgfqpoint{0.806428in}{1.699470in}}%
\pgfpathlineto{\pgfqpoint{0.816991in}{1.703408in}}%
\pgfpathlineto{\pgfqpoint{0.827554in}{1.709901in}}%
\pgfpathlineto{\pgfqpoint{0.838116in}{1.718846in}}%
\pgfpathlineto{\pgfqpoint{0.848679in}{1.730104in}}%
\pgfpathlineto{\pgfqpoint{0.864523in}{1.750926in}}%
\pgfpathlineto{\pgfqpoint{0.880367in}{1.775803in}}%
\pgfpathlineto{\pgfqpoint{0.901493in}{1.813734in}}%
\pgfpathlineto{\pgfqpoint{0.943744in}{1.896714in}}%
\pgfpathlineto{\pgfqpoint{0.970151in}{1.945885in}}%
\pgfpathlineto{\pgfqpoint{0.985995in}{1.971867in}}%
\pgfpathlineto{\pgfqpoint{1.001839in}{1.994060in}}%
\pgfpathlineto{\pgfqpoint{1.017684in}{2.011677in}}%
\pgfpathlineto{\pgfqpoint{1.028246in}{2.020564in}}%
\pgfpathlineto{\pgfqpoint{1.038809in}{2.026994in}}%
\pgfpathlineto{\pgfqpoint{1.049372in}{2.030866in}}%
\pgfpathlineto{\pgfqpoint{1.059935in}{2.032119in}}%
\pgfpathlineto{\pgfqpoint{1.070498in}{2.030734in}}%
\pgfpathlineto{\pgfqpoint{1.081060in}{2.026732in}}%
\pgfpathlineto{\pgfqpoint{1.091623in}{2.020177in}}%
\pgfpathlineto{\pgfqpoint{1.102186in}{2.011171in}}%
\pgfpathlineto{\pgfqpoint{1.112749in}{1.999858in}}%
\pgfpathlineto{\pgfqpoint{1.128593in}{1.978961in}}%
\pgfpathlineto{\pgfqpoint{1.144437in}{1.954024in}}%
\pgfpathlineto{\pgfqpoint{1.165563in}{1.916037in}}%
\pgfpathlineto{\pgfqpoint{1.244783in}{1.766264in}}%
\pgfpathlineto{\pgfqpoint{1.260628in}{1.742778in}}%
\pgfpathlineto{\pgfqpoint{1.276472in}{1.723636in}}%
\pgfpathlineto{\pgfqpoint{1.287035in}{1.713630in}}%
\pgfpathlineto{\pgfqpoint{1.297597in}{1.706018in}}%
\pgfpathlineto{\pgfqpoint{1.308160in}{1.700921in}}%
\pgfpathlineto{\pgfqpoint{1.318723in}{1.698419in}}%
\pgfpathlineto{\pgfqpoint{1.329286in}{1.698551in}}%
\pgfpathlineto{\pgfqpoint{1.339848in}{1.701315in}}%
\pgfpathlineto{\pgfqpoint{1.350411in}{1.706668in}}%
\pgfpathlineto{\pgfqpoint{1.360974in}{1.714525in}}%
\pgfpathlineto{\pgfqpoint{1.371537in}{1.724763in}}%
\pgfpathlineto{\pgfqpoint{1.387381in}{1.744217in}}%
\pgfpathlineto{\pgfqpoint{1.403225in}{1.767964in}}%
\pgfpathlineto{\pgfqpoint{1.424351in}{1.804828in}}%
\pgfpathlineto{\pgfqpoint{1.456039in}{1.866448in}}%
\pgfpathlineto{\pgfqpoint{1.493009in}{1.937486in}}%
\pgfpathlineto{\pgfqpoint{1.514134in}{1.972672in}}%
\pgfpathlineto{\pgfqpoint{1.529979in}{1.994725in}}%
\pgfpathlineto{\pgfqpoint{1.545823in}{2.012178in}}%
\pgfpathlineto{\pgfqpoint{1.556385in}{2.020944in}}%
\pgfpathlineto{\pgfqpoint{1.566948in}{2.027249in}}%
\pgfpathlineto{\pgfqpoint{1.577511in}{2.030991in}}%
\pgfpathlineto{\pgfqpoint{1.588074in}{2.032113in}}%
\pgfpathlineto{\pgfqpoint{1.598637in}{2.030596in}}%
\pgfpathlineto{\pgfqpoint{1.609199in}{2.026464in}}%
\pgfpathlineto{\pgfqpoint{1.619762in}{2.019783in}}%
\pgfpathlineto{\pgfqpoint{1.630325in}{2.010659in}}%
\pgfpathlineto{\pgfqpoint{1.640888in}{1.999234in}}%
\pgfpathlineto{\pgfqpoint{1.656732in}{1.978191in}}%
\pgfpathlineto{\pgfqpoint{1.672576in}{1.953134in}}%
\pgfpathlineto{\pgfqpoint{1.693702in}{1.915036in}}%
\pgfpathlineto{\pgfqpoint{1.772922in}{1.765419in}}%
\pgfpathlineto{\pgfqpoint{1.788767in}{1.742066in}}%
\pgfpathlineto{\pgfqpoint{1.804611in}{1.723081in}}%
\pgfpathlineto{\pgfqpoint{1.815174in}{1.713192in}}%
\pgfpathlineto{\pgfqpoint{1.825736in}{1.705703in}}%
\pgfpathlineto{\pgfqpoint{1.836299in}{1.700734in}}%
\pgfpathlineto{\pgfqpoint{1.846862in}{1.698362in}}%
\pgfpathlineto{\pgfqpoint{1.857425in}{1.698626in}}%
\pgfpathlineto{\pgfqpoint{1.867988in}{1.701522in}}%
\pgfpathlineto{\pgfqpoint{1.878550in}{1.707002in}}%
\pgfpathlineto{\pgfqpoint{1.889113in}{1.714982in}}%
\pgfpathlineto{\pgfqpoint{1.899676in}{1.725334in}}%
\pgfpathlineto{\pgfqpoint{1.915520in}{1.744944in}}%
\pgfpathlineto{\pgfqpoint{1.931364in}{1.768820in}}%
\pgfpathlineto{\pgfqpoint{1.952490in}{1.805809in}}%
\pgfpathlineto{\pgfqpoint{1.984178in}{1.867499in}}%
\pgfpathlineto{\pgfqpoint{2.015866in}{1.928856in}}%
\pgfpathlineto{\pgfqpoint{2.036992in}{1.965273in}}%
\pgfpathlineto{\pgfqpoint{2.052836in}{1.988560in}}%
\pgfpathlineto{\pgfqpoint{2.068680in}{2.007465in}}%
\pgfpathlineto{\pgfqpoint{2.079243in}{2.017297in}}%
\pgfpathlineto{\pgfqpoint{2.089806in}{2.024724in}}%
\pgfpathlineto{\pgfqpoint{2.100369in}{2.029629in}}%
\pgfpathlineto{\pgfqpoint{2.110931in}{2.031935in}}%
\pgfpathlineto{\pgfqpoint{2.121494in}{2.031604in}}%
\pgfpathlineto{\pgfqpoint{2.132057in}{2.028644in}}%
\pgfpathlineto{\pgfqpoint{2.142620in}{2.023099in}}%
\pgfpathlineto{\pgfqpoint{2.153183in}{2.015059in}}%
\pgfpathlineto{\pgfqpoint{2.163745in}{2.004649in}}%
\pgfpathlineto{\pgfqpoint{2.179590in}{1.984962in}}%
\pgfpathlineto{\pgfqpoint{2.195434in}{1.961022in}}%
\pgfpathlineto{\pgfqpoint{2.216559in}{1.923971in}}%
\pgfpathlineto{\pgfqpoint{2.248248in}{1.862248in}}%
\pgfpathlineto{\pgfqpoint{2.279936in}{1.800930in}}%
\pgfpathlineto{\pgfqpoint{2.301062in}{1.764579in}}%
\pgfpathlineto{\pgfqpoint{2.316906in}{1.741359in}}%
\pgfpathlineto{\pgfqpoint{2.332750in}{1.722532in}}%
\pgfpathlineto{\pgfqpoint{2.343313in}{1.712759in}}%
\pgfpathlineto{\pgfqpoint{2.353875in}{1.705394in}}%
\pgfpathlineto{\pgfqpoint{2.364438in}{1.700553in}}%
\pgfpathlineto{\pgfqpoint{2.375001in}{1.698313in}}%
\pgfpathlineto{\pgfqpoint{2.385564in}{1.698709in}}%
\pgfpathlineto{\pgfqpoint{2.396127in}{1.701735in}}%
\pgfpathlineto{\pgfqpoint{2.406689in}{1.707343in}}%
\pgfpathlineto{\pgfqpoint{2.417252in}{1.715444in}}%
\pgfpathlineto{\pgfqpoint{2.427815in}{1.725911in}}%
\pgfpathlineto{\pgfqpoint{2.443659in}{1.745676in}}%
\pgfpathlineto{\pgfqpoint{2.459503in}{1.769680in}}%
\pgfpathlineto{\pgfqpoint{2.480629in}{1.806792in}}%
\pgfpathlineto{\pgfqpoint{2.512317in}{1.868549in}}%
\pgfpathlineto{\pgfqpoint{2.544006in}{1.929825in}}%
\pgfpathlineto{\pgfqpoint{2.565131in}{1.966111in}}%
\pgfpathlineto{\pgfqpoint{2.580975in}{1.989265in}}%
\pgfpathlineto{\pgfqpoint{2.596819in}{2.008012in}}%
\pgfpathlineto{\pgfqpoint{2.607382in}{2.017726in}}%
\pgfpathlineto{\pgfqpoint{2.617945in}{2.025030in}}%
\pgfpathlineto{\pgfqpoint{2.628508in}{2.029806in}}%
\pgfpathlineto{\pgfqpoint{2.639071in}{2.031981in}}%
\pgfpathlineto{\pgfqpoint{2.649633in}{2.031519in}}%
\pgfpathlineto{\pgfqpoint{2.660196in}{2.028427in}}%
\pgfpathlineto{\pgfqpoint{2.670759in}{2.022756in}}%
\pgfpathlineto{\pgfqpoint{2.681322in}{2.014593in}}%
\pgfpathlineto{\pgfqpoint{2.691884in}{2.004069in}}%
\pgfpathlineto{\pgfqpoint{2.707729in}{1.984228in}}%
\pgfpathlineto{\pgfqpoint{2.723573in}{1.960160in}}%
\pgfpathlineto{\pgfqpoint{2.744698in}{1.922987in}}%
\pgfpathlineto{\pgfqpoint{2.776387in}{1.861198in}}%
\pgfpathlineto{\pgfqpoint{2.808075in}{1.799962in}}%
\pgfpathlineto{\pgfqpoint{2.829201in}{1.763742in}}%
\pgfpathlineto{\pgfqpoint{2.845045in}{1.740656in}}%
\pgfpathlineto{\pgfqpoint{2.860889in}{1.721988in}}%
\pgfpathlineto{\pgfqpoint{2.871452in}{1.712332in}}%
\pgfpathlineto{\pgfqpoint{2.882015in}{1.705091in}}%
\pgfpathlineto{\pgfqpoint{2.892577in}{1.700379in}}%
\pgfpathlineto{\pgfqpoint{2.903140in}{1.698270in}}%
\pgfpathlineto{\pgfqpoint{2.913703in}{1.698798in}}%
\pgfpathlineto{\pgfqpoint{2.924266in}{1.701955in}}%
\pgfpathlineto{\pgfqpoint{2.934828in}{1.707690in}}%
\pgfpathlineto{\pgfqpoint{2.945391in}{1.715913in}}%
\pgfpathlineto{\pgfqpoint{2.955954in}{1.726494in}}%
\pgfpathlineto{\pgfqpoint{2.971798in}{1.746412in}}%
\pgfpathlineto{\pgfqpoint{2.987642in}{1.770544in}}%
\pgfpathlineto{\pgfqpoint{3.008768in}{1.807777in}}%
\pgfpathlineto{\pgfqpoint{3.040456in}{1.869599in}}%
\pgfpathlineto{\pgfqpoint{3.072145in}{1.930792in}}%
\pgfpathlineto{\pgfqpoint{3.093270in}{1.966946in}}%
\pgfpathlineto{\pgfqpoint{3.109114in}{1.989965in}}%
\pgfpathlineto{\pgfqpoint{3.124959in}{2.008553in}}%
\pgfpathlineto{\pgfqpoint{3.135521in}{2.018150in}}%
\pgfpathlineto{\pgfqpoint{3.146084in}{2.025329in}}%
\pgfpathlineto{\pgfqpoint{3.156647in}{2.029977in}}%
\pgfpathlineto{\pgfqpoint{3.167210in}{2.032020in}}%
\pgfpathlineto{\pgfqpoint{3.177772in}{2.031426in}}%
\pgfpathlineto{\pgfqpoint{3.188335in}{2.028205in}}%
\pgfpathlineto{\pgfqpoint{3.198898in}{2.022406in}}%
\pgfpathlineto{\pgfqpoint{3.209461in}{2.014122in}}%
\pgfpathlineto{\pgfqpoint{3.220024in}{2.003484in}}%
\pgfpathlineto{\pgfqpoint{3.235868in}{1.983490in}}%
\pgfpathlineto{\pgfqpoint{3.251712in}{1.959294in}}%
\pgfpathlineto{\pgfqpoint{3.272837in}{1.922001in}}%
\pgfpathlineto{\pgfqpoint{3.304526in}{1.860148in}}%
\pgfpathlineto{\pgfqpoint{3.336214in}{1.798996in}}%
\pgfpathlineto{\pgfqpoint{3.357340in}{1.762910in}}%
\pgfpathlineto{\pgfqpoint{3.373184in}{1.739959in}}%
\pgfpathlineto{\pgfqpoint{3.389028in}{1.721450in}}%
\pgfpathlineto{\pgfqpoint{3.399591in}{1.711912in}}%
\pgfpathlineto{\pgfqpoint{3.410154in}{1.704795in}}%
\pgfpathlineto{\pgfqpoint{3.420716in}{1.700211in}}%
\pgfpathlineto{\pgfqpoint{3.431279in}{1.698234in}}%
\pgfpathlineto{\pgfqpoint{3.441842in}{1.698894in}}%
\pgfpathlineto{\pgfqpoint{3.452405in}{1.702181in}}%
\pgfpathlineto{\pgfqpoint{3.462968in}{1.708043in}}%
\pgfpathlineto{\pgfqpoint{3.473530in}{1.716387in}}%
\pgfpathlineto{\pgfqpoint{3.484093in}{1.727082in}}%
\pgfpathlineto{\pgfqpoint{3.499937in}{1.747153in}}%
\pgfpathlineto{\pgfqpoint{3.515781in}{1.771411in}}%
\pgfpathlineto{\pgfqpoint{3.536907in}{1.808764in}}%
\pgfpathlineto{\pgfqpoint{3.568595in}{1.870648in}}%
\pgfpathlineto{\pgfqpoint{3.600284in}{1.931757in}}%
\pgfpathlineto{\pgfqpoint{3.621409in}{1.967776in}}%
\pgfpathlineto{\pgfqpoint{3.637253in}{1.990660in}}%
\pgfpathlineto{\pgfqpoint{3.653098in}{2.009088in}}%
\pgfpathlineto{\pgfqpoint{3.663660in}{2.018567in}}%
\pgfpathlineto{\pgfqpoint{3.674223in}{2.025623in}}%
\pgfpathlineto{\pgfqpoint{3.684786in}{2.030142in}}%
\pgfpathlineto{\pgfqpoint{3.695349in}{2.032053in}}%
\pgfpathlineto{\pgfqpoint{3.705911in}{2.031327in}}%
\pgfpathlineto{\pgfqpoint{3.716474in}{2.027975in}}%
\pgfpathlineto{\pgfqpoint{3.727037in}{2.022050in}}%
\pgfpathlineto{\pgfqpoint{3.737600in}{2.013645in}}%
\pgfpathlineto{\pgfqpoint{3.748163in}{2.002893in}}%
\pgfpathlineto{\pgfqpoint{3.764007in}{1.982747in}}%
\pgfpathlineto{\pgfqpoint{3.779851in}{1.958425in}}%
\pgfpathlineto{\pgfqpoint{3.800977in}{1.921012in}}%
\pgfpathlineto{\pgfqpoint{3.832665in}{1.859098in}}%
\pgfpathlineto{\pgfqpoint{3.864353in}{1.798033in}}%
\pgfpathlineto{\pgfqpoint{3.885479in}{1.762082in}}%
\pgfpathlineto{\pgfqpoint{3.901323in}{1.739266in}}%
\pgfpathlineto{\pgfqpoint{3.917167in}{1.720918in}}%
\pgfpathlineto{\pgfqpoint{3.927730in}{1.711497in}}%
\pgfpathlineto{\pgfqpoint{3.938293in}{1.704504in}}%
\pgfpathlineto{\pgfqpoint{3.948855in}{1.700050in}}%
\pgfpathlineto{\pgfqpoint{3.959418in}{1.698204in}}%
\pgfpathlineto{\pgfqpoint{3.969981in}{1.698996in}}%
\pgfpathlineto{\pgfqpoint{3.980544in}{1.702413in}}%
\pgfpathlineto{\pgfqpoint{3.991107in}{1.708402in}}%
\pgfpathlineto{\pgfqpoint{4.001669in}{1.716867in}}%
\pgfpathlineto{\pgfqpoint{4.012232in}{1.727676in}}%
\pgfpathlineto{\pgfqpoint{4.028076in}{1.747898in}}%
\pgfpathlineto{\pgfqpoint{4.043920in}{1.772282in}}%
\pgfpathlineto{\pgfqpoint{4.065046in}{1.809754in}}%
\pgfpathlineto{\pgfqpoint{4.102016in}{1.882172in}}%
\pgfpathlineto{\pgfqpoint{4.133704in}{1.942182in}}%
\pgfpathlineto{\pgfqpoint{4.154830in}{1.976636in}}%
\pgfpathlineto{\pgfqpoint{4.170674in}{1.997972in}}%
\pgfpathlineto{\pgfqpoint{4.186518in}{2.014593in}}%
\pgfpathlineto{\pgfqpoint{4.197081in}{2.022756in}}%
\pgfpathlineto{\pgfqpoint{4.207644in}{2.028427in}}%
\pgfpathlineto{\pgfqpoint{4.218206in}{2.031519in}}%
\pgfpathlineto{\pgfqpoint{4.228769in}{2.031981in}}%
\pgfpathlineto{\pgfqpoint{4.239332in}{2.029806in}}%
\pgfpathlineto{\pgfqpoint{4.249895in}{2.025030in}}%
\pgfpathlineto{\pgfqpoint{4.260458in}{2.017726in}}%
\pgfpathlineto{\pgfqpoint{4.271020in}{2.008012in}}%
\pgfpathlineto{\pgfqpoint{4.286864in}{1.989265in}}%
\pgfpathlineto{\pgfqpoint{4.302709in}{1.966111in}}%
\pgfpathlineto{\pgfqpoint{4.323834in}{1.929825in}}%
\pgfpathlineto{\pgfqpoint{4.350241in}{1.879035in}}%
\pgfpathlineto{\pgfqpoint{4.392492in}{1.797073in}}%
\pgfpathlineto{\pgfqpoint{4.413618in}{1.761257in}}%
\pgfpathlineto{\pgfqpoint{4.429462in}{1.738578in}}%
\pgfpathlineto{\pgfqpoint{4.445306in}{1.720392in}}%
\pgfpathlineto{\pgfqpoint{4.455869in}{1.711089in}}%
\pgfpathlineto{\pgfqpoint{4.466432in}{1.704221in}}%
\pgfpathlineto{\pgfqpoint{4.476995in}{1.699895in}}%
\pgfpathlineto{\pgfqpoint{4.487557in}{1.698181in}}%
\pgfpathlineto{\pgfqpoint{4.498120in}{1.699105in}}%
\pgfpathlineto{\pgfqpoint{4.508683in}{1.702652in}}%
\pgfpathlineto{\pgfqpoint{4.519246in}{1.708767in}}%
\pgfpathlineto{\pgfqpoint{4.529808in}{1.717353in}}%
\pgfpathlineto{\pgfqpoint{4.540371in}{1.728275in}}%
\pgfpathlineto{\pgfqpoint{4.556215in}{1.748648in}}%
\pgfpathlineto{\pgfqpoint{4.572060in}{1.773157in}}%
\pgfpathlineto{\pgfqpoint{4.593185in}{1.810746in}}%
\pgfpathlineto{\pgfqpoint{4.630155in}{1.883217in}}%
\pgfpathlineto{\pgfqpoint{4.656562in}{1.933677in}}%
\pgfpathlineto{\pgfqpoint{4.677687in}{1.969425in}}%
\pgfpathlineto{\pgfqpoint{4.693532in}{1.992035in}}%
\pgfpathlineto{\pgfqpoint{4.709376in}{2.010141in}}%
\pgfpathlineto{\pgfqpoint{4.719939in}{2.019384in}}%
\pgfpathlineto{\pgfqpoint{4.730501in}{2.026190in}}%
\pgfpathlineto{\pgfqpoint{4.741064in}{2.030451in}}%
\pgfpathlineto{\pgfqpoint{4.751627in}{2.032100in}}%
\pgfpathlineto{\pgfqpoint{4.762190in}{2.031110in}}%
\pgfpathlineto{\pgfqpoint{4.772752in}{2.027497in}}%
\pgfpathlineto{\pgfqpoint{4.783315in}{2.021319in}}%
\pgfpathlineto{\pgfqpoint{4.793878in}{2.012673in}}%
\pgfpathlineto{\pgfqpoint{4.804441in}{2.001695in}}%
\pgfpathlineto{\pgfqpoint{4.820285in}{1.981246in}}%
\pgfpathlineto{\pgfqpoint{4.836129in}{1.956676in}}%
\pgfpathlineto{\pgfqpoint{4.857255in}{1.919029in}}%
\pgfpathlineto{\pgfqpoint{4.894224in}{1.846533in}}%
\pgfpathlineto{\pgfqpoint{4.920631in}{1.796115in}}%
\pgfpathlineto{\pgfqpoint{4.941757in}{1.760437in}}%
\pgfpathlineto{\pgfqpoint{4.957601in}{1.737896in}}%
\pgfpathlineto{\pgfqpoint{4.973445in}{1.719871in}}%
\pgfpathlineto{\pgfqpoint{4.984008in}{1.710687in}}%
\pgfpathlineto{\pgfqpoint{4.994571in}{1.703943in}}%
\pgfpathlineto{\pgfqpoint{5.005134in}{1.699747in}}%
\pgfpathlineto{\pgfqpoint{5.015696in}{1.698164in}}%
\pgfpathlineto{\pgfqpoint{5.026259in}{1.699220in}}%
\pgfpathlineto{\pgfqpoint{5.036822in}{1.702898in}}%
\pgfpathlineto{\pgfqpoint{5.047385in}{1.709139in}}%
\pgfpathlineto{\pgfqpoint{5.057948in}{1.717845in}}%
\pgfpathlineto{\pgfqpoint{5.068510in}{1.728879in}}%
\pgfpathlineto{\pgfqpoint{5.084354in}{1.749403in}}%
\pgfpathlineto{\pgfqpoint{5.100199in}{1.774035in}}%
\pgfpathlineto{\pgfqpoint{5.121324in}{1.811740in}}%
\pgfpathlineto{\pgfqpoint{5.158294in}{1.884261in}}%
\pgfpathlineto{\pgfqpoint{5.184701in}{1.934634in}}%
\pgfpathlineto{\pgfqpoint{5.205826in}{1.970243in}}%
\pgfpathlineto{\pgfqpoint{5.221671in}{1.992715in}}%
\pgfpathlineto{\pgfqpoint{5.237515in}{2.010659in}}%
\pgfpathlineto{\pgfqpoint{5.248078in}{2.019783in}}%
\pgfpathlineto{\pgfqpoint{5.258640in}{2.026464in}}%
\pgfpathlineto{\pgfqpoint{5.269203in}{2.030596in}}%
\pgfpathlineto{\pgfqpoint{5.279766in}{2.032113in}}%
\pgfpathlineto{\pgfqpoint{5.290329in}{2.030991in}}%
\pgfpathlineto{\pgfqpoint{5.300891in}{2.027249in}}%
\pgfpathlineto{\pgfqpoint{5.311454in}{2.020944in}}%
\pgfpathlineto{\pgfqpoint{5.322017in}{2.012178in}}%
\pgfpathlineto{\pgfqpoint{5.332580in}{2.001088in}}%
\pgfpathlineto{\pgfqpoint{5.348424in}{1.980489in}}%
\pgfpathlineto{\pgfqpoint{5.364268in}{1.955796in}}%
\pgfpathlineto{\pgfqpoint{5.385394in}{1.918033in}}%
\pgfpathlineto{\pgfqpoint{5.422363in}{1.845489in}}%
\pgfpathlineto{\pgfqpoint{5.448770in}{1.795160in}}%
\pgfpathlineto{\pgfqpoint{5.469896in}{1.759621in}}%
\pgfpathlineto{\pgfqpoint{5.485740in}{1.737218in}}%
\pgfpathlineto{\pgfqpoint{5.501584in}{1.719356in}}%
\pgfpathlineto{\pgfqpoint{5.512147in}{1.710291in}}%
\pgfpathlineto{\pgfqpoint{5.522710in}{1.703672in}}%
\pgfpathlineto{\pgfqpoint{5.533273in}{1.699605in}}%
\pgfpathlineto{\pgfqpoint{5.543835in}{1.698154in}}%
\pgfpathlineto{\pgfqpoint{5.554398in}{1.699342in}}%
\pgfpathlineto{\pgfqpoint{5.564961in}{1.703149in}}%
\pgfpathlineto{\pgfqpoint{5.575524in}{1.709517in}}%
\pgfpathlineto{\pgfqpoint{5.586087in}{1.718343in}}%
\pgfpathlineto{\pgfqpoint{5.596649in}{1.729489in}}%
\pgfpathlineto{\pgfqpoint{5.612494in}{1.750162in}}%
\pgfpathlineto{\pgfqpoint{5.628338in}{1.774917in}}%
\pgfpathlineto{\pgfqpoint{5.649463in}{1.812736in}}%
\pgfpathlineto{\pgfqpoint{5.686433in}{1.885304in}}%
\pgfpathlineto{\pgfqpoint{5.712840in}{1.935587in}}%
\pgfpathlineto{\pgfqpoint{5.733966in}{1.971057in}}%
\pgfpathlineto{\pgfqpoint{5.749810in}{1.993390in}}%
\pgfpathlineto{\pgfqpoint{5.765654in}{2.011171in}}%
\pgfpathlineto{\pgfqpoint{5.776217in}{2.020177in}}%
\pgfpathlineto{\pgfqpoint{5.786779in}{2.026732in}}%
\pgfpathlineto{\pgfqpoint{5.797342in}{2.030734in}}%
\pgfpathlineto{\pgfqpoint{5.807905in}{2.032119in}}%
\pgfpathlineto{\pgfqpoint{5.818468in}{2.030866in}}%
\pgfpathlineto{\pgfqpoint{5.829031in}{2.026994in}}%
\pgfpathlineto{\pgfqpoint{5.839593in}{2.020564in}}%
\pgfpathlineto{\pgfqpoint{5.850156in}{2.011677in}}%
\pgfpathlineto{\pgfqpoint{5.860719in}{2.000476in}}%
\pgfpathlineto{\pgfqpoint{5.876563in}{1.979727in}}%
\pgfpathlineto{\pgfqpoint{5.892407in}{1.954912in}}%
\pgfpathlineto{\pgfqpoint{5.913533in}{1.917036in}}%
\pgfpathlineto{\pgfqpoint{5.955784in}{1.834073in}}%
\pgfpathlineto{\pgfqpoint{5.982191in}{1.784847in}}%
\pgfpathlineto{\pgfqpoint{5.998035in}{1.758809in}}%
\pgfpathlineto{\pgfqpoint{6.013879in}{1.736545in}}%
\pgfpathlineto{\pgfqpoint{6.029723in}{1.718846in}}%
\pgfpathlineto{\pgfqpoint{6.040286in}{1.709901in}}%
\pgfpathlineto{\pgfqpoint{6.050849in}{1.703408in}}%
\pgfpathlineto{\pgfqpoint{6.061412in}{1.699470in}}%
\pgfpathlineto{\pgfqpoint{6.071975in}{1.698151in}}%
\pgfpathlineto{\pgfqpoint{6.071975in}{1.698151in}}%
\pgfusepath{stroke}%
\end{pgfscope}%
\begin{pgfscope}%
\pgfsetrectcap%
\pgfsetmiterjoin%
\pgfsetlinewidth{0.602250pt}%
\definecolor{currentstroke}{rgb}{0.000000,0.000000,0.000000}%
\pgfsetstrokecolor{currentstroke}%
\pgfsetdash{}{0pt}%
\pgfpathmoveto{\pgfqpoint{0.532060in}{0.456901in}}%
\pgfpathlineto{\pgfqpoint{0.532060in}{3.273333in}}%
\pgfusepath{stroke}%
\end{pgfscope}%
\begin{pgfscope}%
\pgfsetrectcap%
\pgfsetmiterjoin%
\pgfsetlinewidth{0.602250pt}%
\definecolor{currentstroke}{rgb}{0.000000,0.000000,0.000000}%
\pgfsetstrokecolor{currentstroke}%
\pgfsetdash{}{0pt}%
\pgfpathmoveto{\pgfqpoint{0.532060in}{0.456901in}}%
\pgfpathlineto{\pgfqpoint{6.335780in}{0.456901in}}%
\pgfusepath{stroke}%
\end{pgfscope}%
\begin{pgfscope}%
\pgfsetbuttcap%
\pgfsetroundjoin%
\pgfsetlinewidth{2.007500pt}%
\definecolor{currentstroke}{rgb}{0.000000,0.000000,0.000000}%
\pgfsetstrokecolor{currentstroke}%
\pgfsetdash{{7.400000pt}{3.200000pt}}{0.000000pt}%
\pgfpathmoveto{\pgfqpoint{2.116838in}{3.565000in}}%
\pgfpathlineto{\pgfqpoint{2.516838in}{3.565000in}}%
\pgfusepath{stroke}%
\end{pgfscope}%
\begin{pgfscope}%
\definecolor{textcolor}{rgb}{0.000000,0.000000,0.000000}%
\pgfsetstrokecolor{textcolor}%
\pgfsetfillcolor{textcolor}%
\pgftext[x=2.700171in,y=3.506667in,left,base]{\color{textcolor}{\rmfamily\fontsize{12.000000}{14.400000}\selectfont\catcode`\^=\active\def^{\ifmmode\sp\else\^{}\fi}\catcode`\%=\active\def%{\%}Исходный сигнал}}%
\end{pgfscope}%
\begin{pgfscope}%
\pgfsetrectcap%
\pgfsetroundjoin%
\pgfsetlinewidth{2.007500pt}%
\definecolor{currentstroke}{rgb}{1.000000,0.000000,0.000000}%
\pgfsetstrokecolor{currentstroke}%
\pgfsetdash{}{0pt}%
\pgfpathmoveto{\pgfqpoint{2.116838in}{3.329356in}}%
\pgfpathlineto{\pgfqpoint{2.516838in}{3.329356in}}%
\pgfusepath{stroke}%
\end{pgfscope}%
\begin{pgfscope}%
\definecolor{textcolor}{rgb}{0.000000,0.000000,0.000000}%
\pgfsetstrokecolor{textcolor}%
\pgfsetfillcolor{textcolor}%
\pgftext[x=2.700171in,y=3.271022in,left,base]{\color{textcolor}{\rmfamily\fontsize{12.000000}{14.400000}\selectfont\catcode`\^=\active\def^{\ifmmode\sp\else\^{}\fi}\catcode`\%=\active\def%{\%}Аппроксимация рядом Фурье}}%
\end{pgfscope}%
\begin{pgfscope}%
\pgfsetbuttcap%
\pgfsetroundjoin%
\pgfsetlinewidth{1.003750pt}%
\definecolor{currentstroke}{rgb}{0.000000,0.000000,1.000000}%
\pgfsetstrokecolor{currentstroke}%
\pgfsetdash{{1.000000pt}{1.650000pt}}{0.000000pt}%
\pgfpathmoveto{\pgfqpoint{2.116838in}{3.093386in}}%
\pgfpathlineto{\pgfqpoint{2.516838in}{3.093386in}}%
\pgfusepath{stroke}%
\end{pgfscope}%
\begin{pgfscope}%
\definecolor{textcolor}{rgb}{0.000000,0.000000,0.000000}%
\pgfsetstrokecolor{textcolor}%
\pgfsetfillcolor{textcolor}%
\pgftext[x=2.700171in,y=3.035052in,left,base]{\color{textcolor}{\rmfamily\fontsize{12.000000}{14.400000}\selectfont\catcode`\^=\active\def^{\ifmmode\sp\else\^{}\fi}\catcode`\%=\active\def%{\%}Гармоника k=1}}%
\end{pgfscope}%
\begin{pgfscope}%
\pgfsetbuttcap%
\pgfsetroundjoin%
\pgfsetlinewidth{1.003750pt}%
\definecolor{currentstroke}{rgb}{0.000000,0.501961,0.000000}%
\pgfsetstrokecolor{currentstroke}%
\pgfsetdash{{1.000000pt}{1.650000pt}}{0.000000pt}%
\pgfpathmoveto{\pgfqpoint{2.116838in}{2.857741in}}%
\pgfpathlineto{\pgfqpoint{2.516838in}{2.857741in}}%
\pgfusepath{stroke}%
\end{pgfscope}%
\begin{pgfscope}%
\definecolor{textcolor}{rgb}{0.000000,0.000000,0.000000}%
\pgfsetstrokecolor{textcolor}%
\pgfsetfillcolor{textcolor}%
\pgftext[x=2.700171in,y=2.799408in,left,base]{\color{textcolor}{\rmfamily\fontsize{12.000000}{14.400000}\selectfont\catcode`\^=\active\def^{\ifmmode\sp\else\^{}\fi}\catcode`\%=\active\def%{\%}Гармоника k=3}}%
\end{pgfscope}%
\begin{pgfscope}%
\pgfsetbuttcap%
\pgfsetroundjoin%
\pgfsetlinewidth{1.003750pt}%
\definecolor{currentstroke}{rgb}{0.501961,0.000000,0.501961}%
\pgfsetstrokecolor{currentstroke}%
\pgfsetdash{{1.000000pt}{1.650000pt}}{0.000000pt}%
\pgfpathmoveto{\pgfqpoint{2.116838in}{2.622097in}}%
\pgfpathlineto{\pgfqpoint{2.516838in}{2.622097in}}%
\pgfusepath{stroke}%
\end{pgfscope}%
\begin{pgfscope}%
\definecolor{textcolor}{rgb}{0.000000,0.000000,0.000000}%
\pgfsetstrokecolor{textcolor}%
\pgfsetfillcolor{textcolor}%
\pgftext[x=2.700171in,y=2.563763in,left,base]{\color{textcolor}{\rmfamily\fontsize{12.000000}{14.400000}\selectfont\catcode`\^=\active\def^{\ifmmode\sp\else\^{}\fi}\catcode`\%=\active\def%{\%}Гармоника k=5}}%
\end{pgfscope}%
\end{pgfpicture}%
\makeatother%
\endgroup%

    \caption{Аппроксимация входного сигнала отрезком ряда Фурье}
    \label{fig:plot_fourier_approx}
\end{figure}